\documentclass[UTF8, 12pt, fontset=fandol, oneside]{ctexart}
\usepackage{researchnotes}

\title{第三章\quad 线性空间}
\author{}
\date{}

\begin{document}

\maketitle

\section*{§3.1\quad 数域}

读者已经在中学里学到过各种数集，如整数集、有理数集、实数集及复数集. 这些数集不仅是一些数的符号的集合，更重要的是在其上定义了运算. 常见的运算有：加、减、乘、除. 我们常记整数集为 $\mathbb{Z}$，有理数集为 $\mathbb{Q}$，实数集为 $\mathbb{R}$，复数集为 $\mathbb{C}$，则 $\mathbb{Z},\mathbb{Q},\mathbb{R}$ 都是 $\mathbb{C}$ 的子集，即 $\mathbb{C}$ 的一部分. 注意到在整数集 $\mathbb{Z}$ 中，任意两个元素相加、相减或相乘以后仍属于 $\mathbb{Z}$，但是两个整数相除 (除数不为零) 则并不一定属于 $\mathbb{Z}$，它可能是一个分数. 这就是说，整数集 $\mathbb{Z}$ 在加法、减法与乘法下封闭，但在除法下不封闭. 在有理数集 $\mathbb{Q}$ 中，加、减、乘、除都是封闭的. 在实数集及复数集中也如此. 我们把数集的这些特性抽象出来，作如下的定义.

\noindent\textbf{定义 3.1.1}\quad 设 $\mathbb{K}$ 是复数集 $\mathbb{C}$ 的子集且至少有两个不同的元素，如果 $\mathbb{K}$ 中任意两个数的加法、减法、乘法及除法 (除数不为零) 仍属于 $\mathbb{K}$，则称 $\mathbb{K}$ 是一个数域.

根据这个定义，有理数集、实数集及复数集都是数域，而整数集不是数域. 通常我们把在加法、减法、乘法下封闭 (不一定除法封闭) 的数集称为数环，因此整数集是数环但不是数域.

数域是一个比较广泛的概念. 除了已知的有理数域、实数域及复数域外，还有没有其他数域 (即在四则运算下封闭的数集)？我们来看下面的例子.

所有形如
\[
a+b\sqrt{2}
\]
的数，其中 $a,b$ 都是有理数，构成一个数域. 这个数域通常用 $\mathbb{Q}(\sqrt{2})$ 来表示. 现在我们来验证它确是一个数域. 首先注意到：
\[
(a+b\sqrt{2})\pm(c+d\sqrt{2})=(a\pm c)+(b\pm d)\sqrt{2}.
\]
由于两个有理数的和与差仍是有理数，$(a\pm c)+(b\pm d)\sqrt{2}\in\mathbb{Q}(\sqrt{2})$. 也就是说 $\mathbb{Q}(\sqrt{2})$ 在加法、减法下封闭. 又
\[
(a+b\sqrt{2})(c+d\sqrt{2})=(ac+2bd)+(ad+bc)\sqrt{2}.
\]
可见 $\mathbb{Q}(\sqrt{2})$ 在乘法下也封闭. 最后若 $c+d\sqrt{2}\neq 0$，即 $c,d$ 不同时为零，则
\[
\frac{a+b\sqrt{2}}{c+d\sqrt{2}}
=\frac{ac-2bd}{c^2-2d^2}+\frac{bc-ad}{c^2-2d^2}\sqrt{2}.
\]
注意到 $c^2-2d^2\neq 0$，且 $\dfrac{ac-2bd}{c^2-2d^2}$，$\dfrac{bc-ad}{c^2-2d^2}$ 是有理数，因此
\[
\frac{a+b\sqrt{2}}{c+d\sqrt{2}}\in\mathbb{Q}(\sqrt{2}).
\]
这就是说 $\mathbb{Q}(\sqrt{2})$ 在加法、减法、乘法及除法下均封闭，因此它是一个数域.

又如，设 $\pi$ 是圆周率，则所有形如
\begin{equation}
\frac{a_0+a_1\pi+\cdots+a_n\pi^n}{b_0+b_1\pi+\cdots+b_m\pi^m}
\tag{3.1.1}
\end{equation}
的数的全体构成一个数域，其中 $a_i,b_j(i=0,1,\cdots,n;\ j=0,1,\cdots,m)$ 都是有理数，但 $b_j$ 不全为零，$m,n$ 可以为任意非负整数. 这里需要注意的是
\[
b_0+b_1\pi+\cdots+b_m\pi^m=0
\]
当且仅当 $b_0=b_1=\cdots=b_m=0$，即 $\pi$ 是一个超越数. 有了这一点，读者不难自己验证形如 (3.1.1) 式的数全体构成一个数域.

若限制 $a,b$ 是任意整数，则由形如 $a+b\sqrt{2}$ 的所有实数构成的集合只是一个数环而不是一个数域. 这是因为 1 和 2 属于这个数集而 $\dfrac{1}{2}$ 不属于该数集.

任一数域必包含 0 及 1 这两个数. 事实上任意两个相同数之差为 0，两个相同的非零数 (由定义数域必包含非零元) 的商为 1. 因此 0,1 是每个数域都必须拥有的数. 不仅如此，我们有下列命题.

\noindent\textbf{定理 3.1.1}\quad 任一数域必包含有理数域 $\mathbb{Q}$.

\noindent\textbf{证明}\quad 由上面知道，1 必属于任一数域. 将 1 连加 $n$ 次，则 $n$ 也应属于该数域，因此任一正整数属于该数域. 又 $0-n=-n$，因此 $-n$ 也应在此数域中. 因而整数全体都必须在这个数域之中. 最后，若 $m\neq 0$，$m,n$ 为整数，则由除法封闭性可知 $\dfrac{n}{m}$ 也应属于该数域，即任一有理数都应在此数域中. $\square$

定理 3.1.1 告诉我们，有理数域是一个“最小”的数域.

\subsection*{习题 3.1}

1. 判断下列数集是不是数域并说明理由：

(1) 所有偶数全体；

(2) 所有形如 $a+b\sqrt{3}$ 的数的全体，其中 $a,b$ 为有理数；

(3) 所有形如 $a+b\sqrt{-1}$ 的数的全体，其中 $a,b$ 为有理数；

(4) 所有形如 $a\sqrt[3]{2}$ 的数的全体，其中 $a$ 为有理数.

2. 验证：所有形如 (3.1.1) 式所示的数的全体构成一个数域.

3. 验证：所有形如 $a+b\sqrt[3]{2}+c\sqrt[3]{4}$ 的数构成一个数域，其中 $a,b,c$ 是有理数.

4. 验证：所有形如 $a+b\sqrt{2}+c\sqrt{3}+d\sqrt{6}$ 的数构成一个数域，其中 $a,b,c,d$ 是有理数.

\section*{§3.2\quad 行向量和列向量}

读者已经学过平面直角坐标系的概念. 在一个平面上，如果建立了一个直角坐标系，那么平面上任一点 $C$ 均可以用两个有序实数 $(a,b)$ 来表示，其中 $a$ 称为点 $C$ 的横坐标，$b$ 称为点 $C$ 的纵坐标. 反过来，给定两个有序实数 $(a,b)$，必有平面上一点与之对应. 这样平面上的点与有序实数对之间可建立起一个一一对应. 读者还学过向量 (矢量) 的概念. 所谓平面上的向量是指平面上的一条有向线段，其中一端称为起点，另一端称为终点. 如图 3.1 所示，若将点 $O$ 与点 $C$ 连起来并用一个箭头表示这条线段的方向，则 $\overrightarrow{OC}$ 就是一个向量，它的起点为 $O$，终点为 $C$.

现在我们把平面上所有以原点 $O$ 为起点的向量组成的集合记为 $V$. 显然 $V$ 与平面 $xOy$ 上的点之间有一个一一对应：平面 $xOy$ 上的任一点 $C$ 对应于向量 $\overrightarrow{OC}$；反之，任一以 $O$ 为起点的向量对应于平面上一点，即该向量的终点. 起点为 $O$ 而终点也为 $O$ 的向量称为零向量，它对应于原点 $O$. 我们刚才已经提到，平面上的点与有序实数对之间有一个一一对应关系，因此不难看出，平面上任一以原点 $O$ 为始点的向量均可对应一个实数对，即向量 $\overrightarrow{OC}$ 可以用点 $C$ 的坐标 $(a,b)$ 来唯一确定. 这样，我们可以用实数对来代替平面上以原点为始点的全体向量. 或者更直接地，把实数对 $(a,b)$ 就定义为平面上的以原点为始点的向量. 这种把向量“代数化”的方法有着明显的好处：一是可以用代数的工具来研究几何对象；二是它可以推广到更一般的情形，即所谓的 $n$ 维向量. 这种推广不仅是形式上的，而且对于数学的发展及应用起着极其重要的作用. 下面我们就抛开向量的几何形象，来定义 $n$ 维向量的概念.

\noindent\textbf{定义 3.2.1}\quad 设 $\mathbb{K}$ 是一个数域，$a_1,a_2,\cdots,a_n$ 是 $\mathbb{K}$ 中的元素，由 $a_1,a_2,\cdots,a_n$ 组成的有序数组 $(a_1,a_2,\cdots,a_n)$ 称为数域 $\mathbb{K}$ 上的一个 $n$ 维行向量.

对这个定义，我们需要说明以下几点：

第一，我们称 $\alpha=(a_1,a_2,\cdots,a_n)$ 是 $\mathbb{K}$ 上的 $n$ 维行向量. 当然也有 $\mathbb{K}$ 上的 $n$ 维列向量的概念. 如果把 $\mathbb{K}$ 上的 $n$ 个数 $a_1,a_2,\cdots,a_n$ 依次排成一列，就称为 $\mathbb{K}$ 上的一个 $n$ 维列向量：
\[
\left(\begin{array}{c}
a_1\\
a_2\\
\vdots\\
a_n
\end{array}\right).
\]
注意不要把行向量与列向量混为一谈. 即使一个行向量中的元素与一个列向量中的元素对应相等，也不认为它们是一回事. 当然我们将会看到，行向量与列向量有相同的性质，这种相似性的本质我们将在后面加以阐明. 在不引起混淆的情况下，行向量、列向量统称为向量.

第二，两个行向量 $\alpha=(a_1,a_2,\cdots,a_n)$，$\beta=(b_1,b_2,\cdots,b_n)$，仅当 $a_i=b_i(i=1,2,\cdots,n)$ 时相等. 如二维向量 $(1,2)$ 与 $(2,1)$ 是两个不同的向量，虽然它们都由 $1,2$ 两个数组成. 因此两个向量相等不仅要求它们的元素相同，而且要求元素出现的次序也相同.

第三，读者可能已经看出，$n$ 维行向量也可以看成一个 $1\times n$ 矩阵. $n$ 维列向量可以看成是一个 $n\times 1$ 矩阵. 事实上我们确实可以这样看. 对一个 $m\times n$ 矩阵 $A=(a_{ij})$，我们在上一章中已经定义过它的第 $i$ 行为第 $i$ 行向量，第 $j$ 列为第 $j$ 列向量. 用矩阵的观点来看上面的三点说明，那就是不言而喻的了.

对数域 $\mathbb{K}$ 上的 $n$ 维行向量 (或 $n$ 维列向量)，我们可以定义加法、减法及数乘. 这些定义与矩阵的相应运算的定义相同.

若 $\alpha=(a_1,a_2,\cdots,a_n)$，$\beta=(b_1,b_2,\cdots,b_n)$，定义
\[
\begin{aligned}
\alpha+\beta&=(a_1+b_1,a_2+b_2,\cdots,a_n+b_n),\\
\alpha-\beta&=(a_1-b_1,a_2-b_2,\cdots,a_n-b_n).
\end{aligned}
\]
若 $k\in\mathbb{K}$，定义
\[
k\alpha=(ka_1,ka_2,\cdots,ka_n).
\]
由于 $\mathbb{K}$ 是数域，不难看出 $\mathbb{K}$ 上 $n$ 维行向量的和、差及数乘 (该数取自 $\mathbb{K}$) 仍然是 $\mathbb{K}$ 上的 $n$ 维行向量. 对列向量的加法、减法与数乘也可类似定义. 从这里可以看出，如果把向量看成矩阵，其运算就是相应的矩阵运算.

若一个 $n$ 维向量的所有元素都等于零，就称之为零向量，记为 $0$. 但需注意一个 $n$ 维行 (列) 零向量指的是一个 $1\times n(n\times 1)$ 零矩阵. 又若 $\alpha=(a_1,a_2,\cdots,a_n)$，记 $-\alpha=(-a_1,-a_2,\cdots,-a_n)$，称 $-\alpha$ 为 $\alpha$ 的负向量.

\noindent\textbf{向量运算规则}

(1) 加法交换律：$\alpha+\beta=\beta+\alpha$；

(2) 加法结合律：$(\alpha+\beta)+\gamma=\alpha+(\beta+\gamma)$；

(3) $\alpha+0=\alpha$；

(4) $\alpha+(-\alpha)=0$；

(5) $1\cdot\alpha=\alpha$；

(6) $k(\alpha+\beta)=k\alpha+k\beta,\ k\in\mathbb{K}$；

(7) $(k+l)\alpha=k\alpha+l\alpha,\ k,l\in\mathbb{K}$；

(8) $k(l\alpha)=(kl)\alpha$.

上述规则也可以看成是矩阵的相应运算规则.

作为例子，我们来看一下二维实向量的加法与数乘的几何意义. 设在平面直角坐标系内有两个向量 $\alpha=(a_1,a_2)$，$\beta=(b_1,b_2)$. 如图 3.2 所示，则 $\alpha+\beta=(a_1+b_1,a_2+b_2)$. 这和用平行四边形法则求两个向量之和的结果完全一致.

再看图 3.3，设 $\alpha=(a_1,a_2)$，$k$ 是一个实数，则 $k\alpha=(ka_1,ka_2)$ 表示将 $\alpha$ 伸长 $k$ 倍. 当 $k>0$ 时表示在原方向上伸长 $k$ 倍，当 $k<0$ 时表示在反方向上伸长 $-k$ 倍.

对三维实向量，也有类似的几何意义.

由数域 $\mathbb{K}$ 上的 $n$ 维行向量全体组成的集合称为域 $\mathbb{K}$ 上的 $n$ 维行向量空间. 由 $\mathbb{K}$ 上的 $n$ 维列向量全体组成的集合称为 $\mathbb{K}$ 上的 $n$ 维列向量空间. 实数域 $\mathbb{R}$ 上的二维空间与三维空间就是我们所熟悉的平面及三维空间.

\subsection*{习题 3.2}

1. 已知向量 $\alpha=(1,1,0,-1)$，$\beta=(-2,1,0,0)$，$\gamma=(-1,-2,0,1)$，试求下列向量：
\[
\alpha+\beta+\gamma;\quad 3\alpha-\beta+5\gamma.
\]

2. 已知 $\beta=(1,0,1)$，$\gamma=(1,1,-1)$，求解下列向量方程：
\[
3x+\beta=\gamma.
\]

\section*{§3.3\quad 线性空间}

在上一节中，我们定义了行向量空间及列向量空间的概念，它们可以看成是现实的实二维空间与实三维空间的推广. 现在我们要做进一步的抽象，引进一般的向量空间的概念.

\noindent\textbf{定义 3.3.1}\quad 设 $\mathbb{K}$ 是一个数域，$V$ 是一个集合，在 $V$ 上定义了一个加法 ``$+$''，即对 $V$ 中任意两个元素 $\alpha,\beta$，总存在 $V$ 中唯一的元素 $\gamma$ 与之对应，记为 $\gamma=\alpha+\beta$. 在数域 $\mathbb{K}$ 与 $V$ 之间定义了一种运算，称为数乘，即对 $\mathbb{K}$ 中任一数 $k$ 及 $V$ 中任一元素 $\alpha$，在 $V$ 中总有唯一的元素 $\delta$ 与之对应，记为 $\delta=k\alpha$. 若上述加法及数乘满足下列运算规则：

(1) 加法交换律：$\alpha+\beta=\beta+\alpha$；

(2) 加法结合律：$(\alpha+\beta)+\gamma=\alpha+(\beta+\gamma)$；

(3) 在 $V$ 中存在一个元素 $0$，对于 $V$ 中任一元素 $\alpha$，都有 $\alpha+0=\alpha$；

(4) 对于 $V$ 中每个元素 $\alpha$，存在元素 $\beta$，使 $\alpha+\beta=0$；

(5) $1\cdot\alpha=\alpha$；

(6) $k(\alpha+\beta)=k\alpha+k\beta$；

(7) $(k+l)\alpha=k\alpha+l\alpha$；

(8) $k(l\alpha)=(kl)\alpha$，

其中 $\alpha,\beta,\gamma$ 是 $V$ 中任意的元素，$k,l$ 是 $\mathbb{K}$ 中任意的数，则集合 $V$ 称为数域 $\mathbb{K}$ 上的线性空间或向量空间. $V$ 中的元素称为向量，$V$ 中适合 (3) 的元素 $0$ 称为零向量. 对 $V$ 中的元素 $\alpha$，适合 $\alpha+\beta=0$ 的元素 $\beta$ 称为 $\alpha$ 的负向量，记为 $-\alpha$.

读者可能要问：为什么要引进抽象的线性空间？我们先来看几个线性空间的例子.

\noindent\textbf{例 3.3.1}\quad 上一节中数域 $\mathbb{K}$ 上 $n$ 维行向量集合 (列向量集合) 是 $\mathbb{K}$ 上的线性空间，这个空间我们记之为 $\mathbb{K}^n(\mathbb{K}_n)$. 我们以后将看到这种线性空间具有普遍的代表性.

\noindent\textbf{例 3.3.2}\quad 系数取自数域 $\mathbb{K}$ 上的一元多项式全体，记为 $\mathbb{K}[x]$，按照通常的方式定义两个多项式的加法 (同次项系数相加) 及一个数与一个多项式的数乘 (将此数乘以多项式的每一个系数)，则不难验证 $\mathbb{K}[x]$ 是 $\mathbb{K}$ 上的线性空间. 在 $\mathbb{K}[x]$ 中，取次数小于等于 $n$ 的多项式全体，记这个集合为 $\mathbb{K}_n[x]$，则 $\mathbb{K}_n[x]$ 也是 $\mathbb{K}$ 上的线性空间.

\noindent\textbf{例 3.3.3}\quad 闭区间 $[0,1]$ 上的连续函数全体记为 $C[0,1]$，将函数的加法及数乘定义为
\[
(f+g)(x)=f(x)+g(x);\quad (kf)(x)=kf(x),
\]
则 $C[0,1]$ 是实数域 $\mathbb{R}$ 上的线性空间.

\noindent\textbf{例 3.3.4}\quad 数域 $\mathbb{K}$ 上 $m\times n$ 矩阵全体在矩阵的加法与数乘下也构成 $\mathbb{K}$ 上的线性空间.

\noindent\textbf{例 3.3.5}\quad 复数域 $\mathbb{C}$ 可看成是实数域 $\mathbb{R}$ 上的线性空间，这时 $\mathbb{C}$ 上向量的加法就是复数的加法. $\mathbb{R}$ 中元素对 $\mathbb{C}$ 中向量 (即复数) 的乘法就是通常的数的乘法. 一般来说，若两个数域 $\mathbb{K}_1\subseteq\mathbb{K}_2$，则 $\mathbb{K}_2$ 可以看成是 $\mathbb{K}_1$ 上的线性空间，向量就是 $\mathbb{K}_2$ 中的数，向量的加法就是数的加法，数乘就是 $\mathbb{K}_1$ 中的数乘以 $\mathbb{K}_2$ 中的数. 特别地，数域 $\mathbb{K}$ 也可以看成是 $\mathbb{K}$ 自身上的线性空间.

我们称实数域 $\mathbb{R}$ 上的线性空间为实线性空间，称复数域 $\mathbb{C}$ 上的线性空间为复线性空间.

从上面的例子可以看出，抽象线性空间概念的引入使我们扩大了视野，它把众多不同研究对象的共同特点用线性空间这一概念加以概括，从而极大地扩大了代数学理论的应用范围. 在这一章里，我们将用线性空间的理论来进一步讨论线性方程组的解，线性空间的理论是线性代数的核心.

现在我们来研究线性空间的一些最基本的性质.

\noindent\textbf{命题 3.3.1}\quad 零向量是唯一的.

\noindent\textbf{证明}\quad 假设 $0_1,0_2$ 是线性空间 $V$ 中的两个零向量，则
\[
0_1=0_1+0_2=0_2.
\]
这就证明了唯一性. $\square$

\noindent\textbf{命题 3.3.2}\quad 负向量也是唯一的.

\noindent\textbf{证明}\quad 设 $\alpha$ 是 $V$ 中的向量，$\beta_1,\beta_2$ 也是 $V$ 中的向量，且
\[
\alpha+\beta_1=0,\quad \alpha+\beta_2=0,
\]
则
\[
\begin{aligned}
\beta_1
&=\beta_1+0=\beta_1+(\alpha+\beta_2)=(\beta_1+\alpha)+\beta_2\\
&=(\alpha+\beta_1)+\beta_2=0+\beta_2=\beta_2.
\end{aligned}
\]
这说明负向量是唯一的. $\square$

\noindent\textbf{命题 3.3.3}\quad 对任意的 $\alpha,\beta,\gamma\in V$，有

(1) 从 $\alpha+\beta=\alpha+\gamma$ 可推出 $\beta=\gamma$，即加法消去律成立；

(2) $0\cdot\alpha=0$，这里左边的 $0$ 表示数零，右边的 $0$ 表示零向量；

(3) $k\cdot 0=0$；

(4) $(-1)\alpha=-\alpha$；

(5) 若 $k\alpha=0$，则 $\alpha=0$ 或 $k=0$.

\noindent\textbf{证明}\quad (1) $(-\alpha)+(\alpha+\beta)=(-\alpha)+(\alpha+\gamma)$，再由结合律得
\[
((-\alpha)+\alpha)+\beta= ((-\alpha)+\alpha)+\gamma,
\]
即
\[
0+\beta=0+\gamma,
\]
于是 $\beta=\gamma$.

(2) $0\cdot\alpha=(0+0)\alpha=0\cdot\alpha+0\cdot\alpha$，再由 (1) 两边消去 $0\cdot\alpha$ 即得 $0=0\cdot\alpha$.

(3) $k\cdot 0=k(0+0)=k\cdot 0+k\cdot 0$，两边消去 $k\cdot 0$ 即得 $0=k\cdot 0$.

(4) $\alpha+(-1)\alpha=1\cdot\alpha+(-1)\alpha=(1+(-1))\alpha=0\cdot\alpha=0$，因此
\[
(-1)\alpha=-\alpha.
\]

(5) 假设 $k\neq 0$ 且 $k\alpha=0$，则 $k^{-1}$ 存在，故
\[
\alpha=1\cdot\alpha=(k^{-1}\cdot k)\alpha=k^{-1}(k\alpha)=k^{-1}\cdot 0=0.
\]
$\square$

\noindent\textbf{注}\quad (1) 在 $V$ 中我们定义减法为
\[
\alpha-\beta=\alpha+(-1)\beta.
\]
由于消去律成立，$V$ 中元素的等式可以进行 ``移项''，因此，如若
\[
\alpha+\beta=\gamma,
\]
则
\[
\alpha=\gamma-\beta,
\]
或
\[
\alpha+\beta-\gamma=0,
\]
等等. 在形式上与数的加减法运算完全一样.

(2) 由于 $V$ 中元素满足加法结合律，$(\alpha+\beta)+\gamma=\alpha+(\beta+\gamma)$，我们可以不用括号，直接把上述元素写为 $\alpha+\beta+\gamma$. 一般地，几个向量相加，也可以不用括号，因为从结合律非常容易推出当几个向量相加时，相加的先后次序不影响最后的结果. 比如 4 个向量做加法时，有
\[
((\alpha_1+\alpha_2)+\alpha_3)+\alpha_4=(\alpha_1+\alpha_2)+(\alpha_3+\alpha_4)=\alpha_1+(\alpha_2+(\alpha_3+\alpha_4)).
\]
上述向量可写为 $\alpha_1+\alpha_2+\alpha_3+\alpha_4$.

\subsection*{习题 3.3}

1. 判断下列集合是否是实数域 $\mathbb{R}$ 上的线性空间：

(1) $V$ 是次数等于 $n(n\geq 1)$ 的实系数多项式全体，加法和数乘就是多项式的加法和数乘；

(2) $V$ 是 $n$ 阶实上三角阵全体，加法和数乘就是矩阵的加法和数乘；

(3) $V$ 是 $[0,1]$ 区间上可导函数全体，加法和数乘就是函数的加法和数乘；

(4) $V$ 是 $n$ 阶实矩阵全体，数乘就是矩阵的数乘，加法 $\oplus$ 定义为 $A\oplus B=AB-BA$，其中等式右边是矩阵的乘法和减法；

(5) $V$ 是 $n$ 阶实矩阵全体，数乘就是矩阵的数乘，加法 $\oplus$ 定义为 $A\oplus B=AB+BA$，其中等式右边是矩阵的乘法和加法；

(6) $V$ 是以 $0$ 为极限的实数数列全体，即 $V=\{\{a_n\}\mid \lim_{n\to\infty}a_n=0\}$，定义两个数列的加法 $\oplus$ 及数乘 $\circ$ 为：$\{a_n\}\oplus\{b_n\}=\{a_n+b_n\}$，$k\circ\{a_n\}=\{ka_n\}$，其中等式右边分别是数的加法和乘法；

(7) $V$ 是正实数全体 $\mathbb{R}^+$，加法 $\oplus$ 定义为 $a\oplus b=ab$，数乘 $\circ$ 定义为 $k\circ a=a^k$，其中等式右边分别是数的乘法和乘方；

(8) $V$ 是实数对全体 $\{(a,b)\mid a,b\in\mathbb{R}\}$，加法 $\oplus$ 定义为 $(a_1,b_1)\oplus(a_2,b_2)=(a_1+a_2,b_1+b_2+a_1a_2)$，数乘 $\circ$ 定义为 $k\circ(a,b)=(ka,kb+\dfrac{k(k-1)}{2}a^2)$，其中等式右边分别是数的加法和乘法；

(9) $V$ 是实数对全体 $\{(a,b)\mid a,b\in\mathbb{R}\}$，加法 $\oplus$ 定义为 $(a_1,b_1)\oplus(a_2,b_2)=(a_1+a_2,b_1-b_2)$，数乘 $\circ$ 定义为 $k\circ(a,b)=(ka,kb)$，其中等式右边分别是数的加减法和乘法；

(10) $V$ 是实数对全体 $\{(a,b)\mid a,b\in\mathbb{R}\}$，加法 $\oplus$ 定义为 $(a_1,b_1)\oplus(a_2,b_2)=(a_1+a_2,0)$，数乘 $\circ$ 定义为 $k\circ(a,b)=(ka,0)$，其中等式右边分别是数的加法和乘法.

2. 求证在线性空间中下列等式成立：

(1) $-(-\alpha)=\alpha$；

(2) $-(k\alpha)=(-k)\alpha=k(-\alpha)$；

(3) $k(\alpha-\beta)=k\alpha-k\beta$.

\section*{§3.4\quad 向量的线性关系}

设有 $n$ 个未知数 $m$ 个方程式的线性方程组：
\begin{equation}
\left\{
\begin{array}{l}
a_{11}x_1+a_{12}x_2+\cdots+a_{1n}x_n=b_1,\\
a_{21}x_1+a_{22}x_2+\cdots+a_{2n}x_n=b_2,\\
\cdots\cdots\cdots\cdots\\
a_{m1}x_1+a_{m2}x_2+\cdots+a_{mn}x_n=b_m.
\end{array}
\right.
\tag{3.4.1}
\end{equation}
这个方程组我们曾经用矩阵来表示过. 现在我们要用向量来表示该方程组，设方程组的增广矩阵为
\[
\widetilde{A}=
\left(\begin{array}{ccccc}
a_{11}&a_{12}&\cdots&a_{1n}&b_1\\
a_{21}&a_{22}&\cdots&a_{2n}&b_2\\
\vdots&\vdots&&\vdots&\vdots\\
a_{m1}&a_{m2}&\cdots&a_{mn}&b_m
\end{array}\right),
\]
分别用 $\alpha_1,\alpha_2,\cdots,\alpha_n,\beta$ 表示上述矩阵的列向量，即
\[
\alpha_1=\left(\begin{array}{c}
a_{11}\\
a_{21}\\
\vdots\\
a_{m1}
\end{array}\right),\quad
\alpha_2=\left(\begin{array}{c}
a_{12}\\
a_{22}\\
\vdots\\
a_{m2}
\end{array}\right),\quad \cdots,\quad
\alpha_n=\left(\begin{array}{c}
a_{1n}\\
a_{2n}\\
\vdots\\
a_{mn}
\end{array}\right);\quad
\beta=\left(\begin{array}{c}
b_1\\
b_2\\
\vdots\\
b_m
\end{array}\right),
\]
则方程组 (3.4.1) 等价于下列向量形式的方程式：
\begin{equation}
x_1\alpha_1+x_2\alpha_2+\cdots+x_n\alpha_n=\beta.
\tag{3.4.2}
\end{equation}

\noindent\textbf{定义 3.4.1}\quad 设 $V$ 是数域 $\mathbb{K}$ 上的线性空间，$\alpha_1,\alpha_2,\cdots,\alpha_n$ 和 $\beta$ 均是 $V$ 中的向量，若存在 $\mathbb{K}$ 中的 $n$ 个数 $k_1,k_2,\cdots,k_n$，使
\[
\beta=k_1\alpha_1+k_2\alpha_2+\cdots+k_n\alpha_n,
\]
则称 $\beta$ 是 $\alpha_1,\alpha_2,\cdots,\alpha_n$ 的线性组合或 $\beta$ 可由 $\alpha_1,\alpha_2,\cdots,\alpha_n$ 线性表示.

显而易见，方程组 (3.4.1) 有解当且仅当向量 $\beta$ 可以表示为向量 $\alpha_1,\alpha_2,\cdots,\alpha_n$ 的线性组合.

\noindent\textbf{例 3.4.1}\quad 设 $V$ 是三维行向量空间，下列向量
\[
\alpha_1=(1,0,1),\quad \alpha_2=(-1,2,2),\quad \beta=(1,2,4)
\]
适合关系式
\[
\beta=2\alpha_1+\alpha_2,
\]
因此 $\beta$ 是 $\alpha_1,\alpha_2$ 的线性组合.

\noindent\textbf{例 3.4.2}\quad 设 $e_1=(1,0,\cdots,0),e_2=(0,1,\cdots,0),\cdots,e_n=(0,0,\cdots,1)$ 是 $n$ 维标准单位行向量，则对任一 $n$ 维行向量 $\alpha=(a_1,a_2,\cdots,a_n)$，
\[
\alpha=a_1e_1+a_2e_2+\cdots+a_ne_n,
\]
即任一 $n$ 维行向量 $\alpha$ 均可由 $e_1,e_2,\cdots,e_n$ 线性表示.

再看齐次线性方程组：
\begin{equation}
\left\{
\begin{array}{l}
a_{11}x_1+a_{12}x_2+\cdots+a_{1n}x_n=0,\\
a_{21}x_1+a_{22}x_2+\cdots+a_{2n}x_n=0,\\
\cdots\cdots\cdots\cdots\\
a_{m1}x_1+a_{m2}x_2+\cdots+a_{mn}x_n=0.
\end{array}
\right.
\tag{3.4.3}
\end{equation}
这个方程组等价于下列向量形式的方程式：
\begin{equation}
x_1\alpha_1+x_2\alpha_2+\cdots+x_n\alpha_n=0.
\tag{3.4.4}
\end{equation}

\noindent\textbf{定义 3.4.2}\quad 设 $V$ 是数域 $\mathbb{K}$ 上的线性空间，$\alpha_1,\alpha_2,\cdots,\alpha_n$ 是 $V$ 中的 $n$ 个向量，若存在 $\mathbb{K}$ 中不全为零的 $n$ 个数 $k_1,k_2,\cdots,k_n$，使
\[
k_1\alpha_1+k_2\alpha_2+\cdots+k_n\alpha_n=0,
\]
则称 $\alpha_1,\alpha_2,\cdots,\alpha_n$ 线性相关. 反之，若不存在 $\mathbb{K}$ 中不全为零的数 $k_1,k_2,\cdots,k_n$，使上式成立，则称 $\alpha_1,\alpha_2,\cdots,\alpha_n$ 线性无关或线性独立.

于是，方程组 (3.4.3) 有非零解 (即在解中至少有一个数不等于零) 的充分必要条件是向量 $\alpha_1,\alpha_2,\cdots,\alpha_n$ 线性相关.

\noindent\textbf{注}\quad (1) 在线性相关、线性无关的定义中，数 $k_1,k_2,\cdots,k_n$ 必须取自数域 $\mathbb{K}$. 举例来说，若把复数域看成是实数域上的线性空间，那么 $1$ 与 $\mathrm{i}=\sqrt{-1}$ 是两个线性无关的向量，因为不存在不全为零的实数 $a,b$ 使 $a+b\mathrm{i}=0$. 但是如果允许 $a,b$ 取复数，比如取 $a=1,b=\mathrm{i}$，就有 $a+b\mathrm{i}=0$.

(2) 线性无关还可这样等价地定义：若存在 $k_1,k_2,\cdots,k_n\in\mathbb{K}$，使
\[
k_1\alpha_1+k_2\alpha_2+\cdots+k_n\alpha_n=0,
\]
则必有 $k_1=k_2=\cdots=k_n=0$.

\noindent\textbf{例 3.4.3}\quad 设 $V$ 是三维行向量空间，下列向量
\[
\alpha_1=(1,2,3),\quad \alpha_2=(2,-1,-4),\quad \alpha_3=(1,1,1)
\]
满足关系式
\[
3\alpha_1+\alpha_2-5\alpha_3=0,
\]
因此 $\alpha_1,\alpha_2,\alpha_3$ 线性相关.

\noindent\textbf{例 3.4.4}\quad 同例 3.4.2 的假设，则 $e_1,e_2,\cdots,e_n$ 线性无关.

\noindent\textbf{证明}\quad 假设
\[
k_1e_1+k_2e_2+\cdots+k_ne_n=0,
\]
则 $(k_1,k_2,\cdots,k_n)=0$，即 $k_1=k_2=\cdots=k_n=0$，因此这 $n$ 个向量线性无关. $\square$

\noindent\textbf{注}\quad 对 $n$ 维标准单位列向量也有同例 3.4.2 和例 3.4.4 一样的结论.

\noindent\textbf{例 3.4.5}\quad 设 $V$ 是数域 $\mathbb{K}$ 上的线性空间，若 $S$ 是只含一个向量 $\alpha$ 的向量组，则 $S$ 线性相关的充分必要条件是 $\alpha=0$. 又若 $S$ 是含有零向量的向量组，则 $S$ 必线性相关.

\noindent\textbf{证明}\quad 从 $k\alpha=0$ 及 $k\neq 0$ 可推出 $\alpha=0$，反之亦然. 因此单个向量 $\alpha$ 线性相关的充分必要条件是 $\alpha=0$.

设 $0,\alpha_2,\cdots,\alpha_m$ 是一个向量组，则
\[
1\cdot 0+0\cdot\alpha_2+\cdots+0\cdot\alpha_m=0,
\]
即这组向量线性相关. $\square$

在线性空间理论中，线性组合、线性相关及线性无关是向量之间最基本的关系，统称为向量的线性关系. 接下去，我们要探讨向量线性关系最基本的性质.

\noindent\textbf{定理 3.4.1}\quad 若 $\alpha_1,\alpha_2,\cdots,\alpha_m$ 是一组线性相关的向量，则任一包含这组向量的向量组必线性相关. 又若 $\alpha_1,\alpha_2,\cdots,\alpha_m$ 是一组线性无关的向量，则从这一组向量中任意取出一组向量必线性无关.

\noindent\textbf{证明}\quad 设 $\alpha_1,\alpha_2,\cdots,\alpha_m,\alpha_{m+1},\cdots,\alpha_n$ 是包含 $\alpha_1,\alpha_2,\cdots,\alpha_m$ 的一组向量. 若 $\alpha_1,\alpha_2,\cdots,\alpha_m$ 线性相关，则存在不全为零的一组数 $k_1,k_2,\cdots,k_m$，使
\[
k_1\alpha_1+k_2\alpha_2+\cdots+k_m\alpha_m=0.
\]
令 $k_{m+1}=\cdots=k_n=0$，则仍有
\[
k_1\alpha_1+k_2\alpha_2+\cdots+k_m\alpha_m+k_{m+1}\alpha_{m+1}+\cdots+k_n\alpha_n=0.
\]
因此 $\alpha_1,\alpha_2,\cdots,\alpha_n$ 线性相关. 另一个论断显然和已证明的结论是等价的. $\square$

\noindent\textbf{定理 3.4.2}\quad 设 $\alpha_1,\alpha_2,\cdots,\alpha_m$ 是线性空间 $V$ 中的向量，则 $\alpha_1,\alpha_2,\cdots,\alpha_m$ 线性相关的充分必要条件是其中至少有一个向量可以表示为其余向量的线性组合.

\noindent\textbf{证明}\quad 设 $\alpha_1,\alpha_2,\cdots,\alpha_m$ 线性相关，则存在不全为零的一组数 $k_1,k_2,\cdots,k_m$，使 $k_1\alpha_1+k_2\alpha_2+\cdots+k_m\alpha_m=0$，其中有某个 $k_i\neq 0$. 于是
\[
\alpha_i=-\frac{k_1}{k_i}\alpha_1-\cdots-\frac{k_{i-1}}{k_i}\alpha_{i-1}-\frac{k_{i+1}}{k_i}\alpha_{i+1}-\cdots-\frac{k_m}{k_i}\alpha_m,
\]
即 $\alpha_i$ 是其余 $m-1$ 个向量的线性组合.

反过来，若
\[
\alpha_i=b_1\alpha_1+\cdots+b_{i-1}\alpha_{i-1}+b_{i+1}\alpha_{i+1}+\cdots+b_m\alpha_m,
\]
则
\[
b_1\alpha_1+\cdots+b_{i-1}\alpha_{i-1}+(-1)\alpha_i+b_{i+1}\alpha_{i+1}+\cdots+b_m\alpha_m=0,
\]
即 $\alpha_1,\alpha_2,\cdots,\alpha_m$ 线性相关. $\square$

\noindent\textbf{定理 3.4.3}\quad 设 $\alpha_1,\alpha_2,\cdots,\alpha_m,\beta$ 是线性空间 $V$ 中的向量. 已知 $\beta$ 可表示为 $\alpha_1,\alpha_2,\cdots,\alpha_m$ 的线性组合，即
\[
\beta=k_1\alpha_1+k_2\alpha_2+\cdots+k_m\alpha_m,
\]
则表示唯一的充分必要条件是向量 $\alpha_1,\alpha_2,\cdots,\alpha_m$ 线性无关.

\noindent\textbf{证明}\quad 假设向量 $\alpha_1,\alpha_2,\cdots,\alpha_m$ 线性无关且另外有一个表示：
\[
\beta=b_1\alpha_1+b_2\alpha_2+\cdots+b_m\alpha_m,
\]
则将已知的两个表示式相减得到
\[
(b_1-k_1)\alpha_1+(b_2-k_2)\alpha_2+\cdots+(b_m-k_m)\alpha_m=0.
\]
因为 $\alpha_1,\alpha_2,\cdots,\alpha_m$ 线性无关，故 $b_1-k_1=b_2-k_2=\cdots=b_m-k_m=0$，即 $b_i=k_i(i=1,2,\cdots,m)$. 也就是说 $\beta$ 只能用唯一一种方式表示为 $\alpha_1,\alpha_2,\cdots,\alpha_m$ 的线性组合.

反之，若向量 $\alpha_1,\alpha_2,\cdots,\alpha_m$ 线性相关，即存在不全为零的数 $c_1,c_2,\cdots,c_m$，使得
\[
c_1\alpha_1+c_2\alpha_2+\cdots+c_m\alpha_m=0,
\]
则除了已知的表示外，$\beta$ 还有另外一个不同的表示：
\[
\beta=(k_1+c_1)\alpha_1+(k_2+c_2)\alpha_2+\cdots+(k_m+c_m)\alpha_m.
\]
$\square$

\noindent\textbf{注}\quad 由上述定理知道，如果方程组 (3.4.1) 或方程式 (3.4.2) 有解，则有唯一解的充分必要条件是向量组 $\alpha_1,\alpha_2,\cdots,\alpha_n$ 线性无关.

\noindent\textbf{定理 3.4.4}\quad 设向量组 $A=\{\alpha_1,\alpha_2,\cdots,\alpha_m\}$，$B=\{\beta_1,\beta_2,\cdots,\beta_n\}$ 和 $C=\{\gamma_1,\gamma_2,\cdots,\gamma_p\}$ 满足：$A$ 中任一向量都是 $B$ 中向量的线性组合，$B$ 中任一向量都是 $C$ 中向量的线性组合，则 $A$ 中任一向量都是 $C$ 中向量的线性组合.

\noindent\textbf{证明}\quad 设
\[
\alpha_i=\sum_{j=1}^{n}a_{ij}\beta_j\ (1\leq i\leq m);\quad \beta_j=\sum_{k=1}^{p}b_{jk}\gamma_k\ (1\leq j\leq n),
\]
则可得
\[
\alpha_i=\sum_{j=1}^{n}a_{ij}\left(\sum_{k=1}^{p}b_{jk}\gamma_k\right)=\sum_{k=1}^{p}\left(\sum_{j=1}^{n}a_{ij}b_{jk}\right)\gamma_k,
\]
即任一 $\alpha_i$ 都是 $\gamma_1,\gamma_2,\cdots,\gamma_p$ 的线性组合. $\square$

\noindent\textbf{例 3.4.6}\quad 若 $\alpha=(a_1,a_2,\cdots,a_n)$，$\beta=(b_1,b_2,\cdots,b_n)$ 是两个 $n$ 维行向量，则 $\alpha,\beta$ 线性相关的充分必要条件是 $a_i,b_i$ 成比例.

\noindent\textbf{证明}\quad 假设 $\alpha,\beta$ 线性相关，由定理 3.4.2，不妨设 $\beta$ 是 $\alpha$ 的线性组合. 令 $\beta=k\alpha$，则
\[
(ka_1,ka_2,\cdots,ka_n)=(b_1,b_2,\cdots,b_n).
\]
因此 $ka_i=b_i(1\leq i\leq n)$，即 $a_i,b_i$ 成比例.

反之，若 $ka_i=b_i(1\leq i\leq n)$，则 $k\alpha-\beta=0$，即 $\alpha,\beta$ 线性相关. $\square$

请读者思考下面两个问题：

(1) 在三维几何空间中，3 个以原点为始点的向量线性相关、线性无关的几何意义是什么？在 Descartes (笛卡尔) 平面上，两个以原点为始点的向量线性相关的几何意义又是什么？在 Descartes 平面上，3 个以原点为始点的向量是否一定线性相关？

(2) 我们先看下面的线性方程组：
\[
\left\{
\begin{array}{rcl}
x_1+x_2+x_3&=&3,\\
2x_1-x_2-3x_3&=&-2,\\
4x_1+x_2-x_3&=&4.
\end{array}
\right.
\]
这个方程组的第三个方程式是多余的，因为它可以由第一个方程式乘以 2 加上第二个方程式得到. 用向量的语言来说，就是方程组增广矩阵
\begin{equation}
\widetilde{A}=
\left(\begin{array}{rrrr}
1&1&1&3\\
2&-1&-3&-2\\
4&1&-1&4
\end{array}\right)
\tag{3.4.5}
\end{equation}
的第三个行向量可以表示为第一和第二个行向量的线性组合. 现假设有线性方程组 (3.4.1)，其增广矩阵的行向量线性相关，你能得到什么结论？

\subsection*{习题 3.4}

1. 确定下列三维实向量是线性相关还是线性无关：

(1) $(-1,3,1),(2,1,0),(1,4,1)$；

(2) $(2,3,0),(-1,4,0),(0,0,2)$.

2. 问 $a$ 取何值时，下列实向量线性相关？
\[
(a,1,1),\ (1,a,1),\ (1,1,a).
\]

3. $\alpha_1,\alpha_2,\cdots,\alpha_m$ 是一组线性无关的向量，$c$ 是非零常数，问 $c\alpha_1,c\alpha_2,\cdots,c\alpha_m$ 是否线性无关？

4. 若 $\alpha_1,\alpha_2$ 线性相关，$\beta_1,\beta_2$ 线性相关，问 $\alpha_1+\beta_1$ 和 $\alpha_2+\beta_2$ 是否必线性相关？

5. 若 $\alpha_1$ 和 $\alpha_2$ 线性无关，$\beta$ 是另一个向量，问 $\alpha_1+\beta$ 和 $\alpha_2+\beta$ 是否必线性无关？

6. 若 $\alpha,\beta$ 线性无关，$\alpha,\gamma$ 线性无关，$\beta,\gamma$ 线性无关，问 $\alpha,\beta,\gamma$ 是否必线性无关？

7. 设 $\alpha_1,\alpha_2,\cdots,\alpha_m$ 是线性空间 $V$ 中一组线性无关的向量，$\beta$ 是 $V$ 中的向量. 求证：或者 $\alpha_1,\alpha_2,\cdots,\alpha_m,\beta$ 线性无关，或者 $\beta$ 是 $\alpha_1,\alpha_2,\cdots,\alpha_m$ 的线性组合.

8. 设向量 $\beta$ 可由向量 $\alpha_1,\alpha_2,\cdots,\alpha_m$ 线性表示，但不能由其中任何个数少于 $m$ 的部分向量线性表示，求证：这 $m$ 个向量线性无关.

9. 设线性空间 $V$ 中向量 $\alpha_1,\alpha_2,\cdots,\alpha_r$ 线性无关，已知有序向量组 $\{\beta,\alpha_1,\alpha_2,\cdots,\alpha_r\}$ 线性相关，求证：最多只有一个 $\alpha_i$ 可以表示为前面向量的线性组合.

10. 设 $n$ 维列向量 $\alpha_1,\alpha_2,\cdots,\alpha_m$ 线性无关，$A$ 为 $n$ 阶可逆阵，求证：$A\alpha_1,A\alpha_2,\cdots,A\alpha_m$ 线性无关.

11. 设 $A$ 是 $n\times m$ 矩阵，$B$ 是 $m\times n$ 矩阵，满足 $AB=I_n$，求证：$B$ 的 $n$ 个列向量线性无关.

12. 设 $\{\alpha_i=(a_{i1},a_{i2},\cdots,a_{in}),1\leq i\leq m\}$ 是一组 $n$ 维行向量，$1\leq j_1<j_2<\cdots<j_t\leq n$ 是给定的 $t(t<n)$ 个指标. 定义 $\widetilde{\alpha_i}=(a_{ij_1},a_{ij_2},\cdots,a_{ij_t})$，称 $\widetilde{\alpha_i}$ 为 $\alpha_i$ 的 $t$ 维缩短向量. 求证：

(1) 若 $\alpha_1,\alpha_2,\cdots,\alpha_m$ 线性相关，则 $\widetilde{\alpha_1},\widetilde{\alpha_2},\cdots,\widetilde{\alpha_m}$ 也线性相关；

(2) 设 $n$ 维行向量 $\alpha=(a_1,a_2,\cdots,a_n)$ 是 $\alpha_1,\alpha_2,\cdots,\alpha_m$ 的线性组合，则 $\widetilde{\alpha}$ 也是 $\widetilde{\alpha_1},\widetilde{\alpha_2},\cdots,\widetilde{\alpha_m}$ 的线性组合.

\section*{§3.5\quad 向量组的秩}

在上一节最后，我们在思考题 (2) 中发现，矩阵 (3.4.5) 的 3 个行向量线性相关，第三个行向量可以用其余两个行向量线性表示. 这表明原线性方程组的第三个方程式是多余的. 我们也不难证明第一和第二两个行向量是线性无关的，因此如果将原方程组的第三个方程式去掉以后，剩下的两个方程式再也不能去掉了，否则得到的方程组将和原方程组不同解. 一般来说，给定一组向量，如果线性相关，这时必有某个向量可以用其余向量线性表示，我们将它去掉. 不断地重复这个过程直到剩下的向量线性无关为止，剩下的向量就称为原向量组的极大无关组. 我们给它下一个严格的定义.

\noindent\textbf{定义 3.5.1}\quad 在线性空间 $V$ 中，向量的集合称为向量族，向量的有限集合称为向量组. 设 $S$ 是向量族，若在 $S$ 中存在一组向量 $\{\alpha_1,\alpha_2,\cdots,\alpha_r\}$ 满足如下条件：

(1) $\alpha_1,\alpha_2,\cdots,\alpha_r$ 线性无关；

(2) $S$ 中任意一个向量都可以用 $\alpha_1,\alpha_2,\cdots,\alpha_r$ 线性表示，

则称 $\{\alpha_1,\alpha_2,\cdots,\alpha_r\}$ 是向量族 $S$ 的极大线性无关组，简称极大无关组.

\noindent\textbf{注}\quad 上述定义 (2) 表明若将 $S$ 中任一向量 $\alpha$ 加入 $\{\alpha_1,\alpha_2,\cdots,\alpha_r\}$，则向量组 $\{\alpha_1,\alpha_2,\cdots,\alpha_r,\alpha\}$ 一定线性相关. 正是在这个意义上，我们称 $\{\alpha_1,\alpha_2,\cdots,\alpha_r\}$ 为极大线性无关组.

\noindent\textbf{命题 3.5.1}\quad 设 $S$ 是一个向量组且至少包含一个非零向量，则 $S$ 的极大无关组一定存在.

\noindent\textbf{证明}\quad 设 $S$ 所含向量的个数为 $k$，对 $k$ 用归纳法进行证明. 若 $k=1$，则 $S$ 由一个非零向量 $\alpha$ 组成，于是 $\{\alpha\}$ 就是 $S$ 的极大无关组. 一般地，若 $S$ 中的 $k$ 个向量线性无关，那么这 $k$ 个向量就构成了 $S$ 的极大无关组. 若这 $k$ 个向量线性相关，由定理 3.4.2 知至少有一个向量是其余向量的线性组合，不妨设为 $\alpha$. 考虑向量组 $S\setminus\{\alpha\}$，它所含向量的个数为 $k-1$ 且至少包含一个非零向量 (否则容易推出 $S$ 只包含零向量)，由归纳假设知 $S\setminus\{\alpha\}$ 存在极大无关组 $\{\alpha_1,\alpha_2,\cdots,\alpha_r\}$. 由定理 3.4.4 知 $\alpha$ 也可由 $\alpha_1,\alpha_2,\cdots,\alpha_r$ 线性表示，因此 $\{\alpha_1,\alpha_2,\cdots,\alpha_r\}$ 也是 $S$ 的极大无关组. $\square$

一个向量族的极大无关组唯一吗？我们来看一个简单的例子.

\noindent\textbf{例 3.5.1}\quad 设有向量组 $S=\{(1,0),(0,1),(1,1)\}$. 这 3 个向量线性相关，但不难验证 $S$ 的 3 个子集 $\{(1,0),(0,1)\}$，$\{(1,0),(1,1)\}$ 以及 $\{(0,1),(1,1)\}$ 都是 $S$ 的极大无关组. 因此一般来说，向量族的极大无关组并不唯一.

虽然极大无关组不唯一，但是我们在例 3.5.1 中发现，向量组 $S$ 的每个极大无关组所含向量的个数是相同的. 我们要问：这个结论对一般的向量组还对吗？

我们不妨来分析一下：假设已知向量族 $S$ 有两个极大无关组 $A,B$. 由极大无关组的定义，$A$ 和 $B$ 都是线性无关的向量组，且 $A$ 中每个向量可以用 $B$ 中向量线性表示，$B$ 中每个向量也可以用 $A$ 中向量线性表示. 我们希望证明：两个线性无关的向量组如果能够互相线性表示，则它们含有相同个数的向量. 这只需证明如下命题.

\noindent\textbf{引理 3.5.1}\quad 设 $A,B$ 是 $V$ 中两组向量，$A$ 含有 $r$ 个向量，$B$ 含有 $s$ 个向量，且 $A$ 中每个向量均可用 $B$ 中向量线性表示. 如果 $A$ 中向量线性无关，则 $r\leq s$.

\noindent\textbf{证明}\quad 我们用反证法. 设
\[
A=\{\alpha_1,\alpha_2,\cdots,\alpha_r\},\quad
B=\{\beta_1,\beta_2,\cdots,\beta_s\}.
\]
假设 $r>s$，我们来推出矛盾.

由已知，$A$ 中向量 $\alpha_1$ 可由 $B$ 中向量的线性组合来表示，即存在数 $\lambda_1,\lambda_2,\cdots,\lambda_s$，使
\begin{equation}
\alpha_1=\lambda_1\beta_1+\lambda_2\beta_2+\cdots+\lambda_s\beta_s.
\tag{3.5.1}
\end{equation}
因为 $A$ 中向量线性无关，故 $\alpha_1\neq 0$，从而 $\lambda_i$ 中至少有一个不为零，不妨假设 $\lambda_1\neq 0$. 由 (3.5.1) 式解出 $\beta_1$：
\begin{equation}
\beta_1=\frac{1}{\lambda_1}\alpha_1-\frac{\lambda_2}{\lambda_1}\beta_2-\cdots-\frac{\lambda_s}{\lambda_1}\beta_s.
\tag{3.5.2}
\end{equation}
但对任意的 $\alpha_i(i=2,3,\cdots,r)$，已知 $\alpha_i$ 可由 $\beta_1,\beta_2,\cdots,\beta_s$ 的线性组合表示，将 (3.5.2) 式代入 $\alpha_i$ 的表示式，则 $\alpha_i$ 可由 $\alpha_1,\beta_2,\cdots,\beta_s$ 的线性组合表示. 这样，我们可将 $B$ 中的向量 $\beta_1$ 换成 $\alpha_1$，这时 $A$ 中任一向量仍可用 $B$ 中向量的线性组合来表示.

现在我们用归纳法，设 $B$ 中向量已经换成 $\{\alpha_1,\cdots,\alpha_k,\beta_{k+1},\cdots,\beta_s\}$，且 $A$ 中任一向量都可以用 $\{\alpha_1,\cdots,\alpha_k,\beta_{k+1},\cdots,\beta_s\}$ 的线性组合表示. 假设 $k<r$，则 $\alpha_{k+1}$ 可表示为
\[
\alpha_{k+1}=\mu_1\alpha_1+\cdots+\mu_k\alpha_k+\mu_{k+1}\beta_{k+1}+\cdots+\mu_s\beta_s,
\]
其中至少有一个 $\mu_i(i=k+1,\cdots,s)$ 不为零. 这是因为若 $\mu_{k+1}=\cdots=\mu_s=0$，则 $\alpha_{k+1}$ 可用 $\alpha_1,\cdots,\alpha_k$ 线性表示，这与 $A$ 中向量线性无关矛盾. 不失一般性，可设 $\mu_{k+1}\neq 0$. 用与上述相同的论证，又可将 $\beta_{k+1}$ 换成 $\alpha_{k+1}$，得到向量组 $\{\alpha_1,\cdots,\alpha_{k+1},\beta_{k+2},\cdots,\beta_s\}$，且 $A$ 中任一向量都可以用这组向量的线性组合表示. 这一事实表明，我们可将 $A$ 中向量依次换入 $B$. 但 $r>s$，因此可将 $A$ 中 $s$ 个向量换入 $B$ 中. 不妨设 $B$ 经调换以后的向量组为 $\{\alpha_1,\alpha_2,\cdots,\alpha_s\}$，则 $A$ 中向量 $\alpha_r$ 也可用 $\alpha_1,\alpha_2,\cdots,\alpha_s$ 的线性组合来表示，从而向量组 $\{\alpha_1,\alpha_2,\cdots,\alpha_s,\alpha_r\}$ 线性相关，引出矛盾. $\square$

引理 3.5.1 的逆否命题也非常有用，为了记住它，我们用一句话来概括：“多”若可以用“少”来线性表示，则“多”线性相关.

\noindent\textbf{引理 3.5.2}\quad 设 $A,B$ 都是线性无关的向量组，又 $A$ 中任一向量可用 $B$ 中向量的线性组合来表示，$B$ 中任一向量也可用 $A$ 中向量的线性组合来表示，则这两组向量所含的向量个数相等.

\noindent\textbf{证明}\quad 设 $A$ 有 $r$ 个向量，$B$ 有 $s$ 个向量. 由引理 3.5.1 得 $r\leq s$. 同理又有 $s\leq r$，故 $r=s$. $\square$

\noindent\textbf{定理 3.5.1}\quad 设 $A$ 与 $B$ 都是向量族 $S$ 的极大线性无关组，则 $A$ 与 $B$ 所含的向量个数相等.

\noindent\textbf{证明}\quad 由定义 3.5.1 及引理 3.5.2 即得. $\square$

\noindent\textbf{定义 3.5.2}\quad 向量族 $S$ 的极大无关组所含的向量个数称为 $S$ 的秩，记作 $\mathrm{r}(S)$ 或 $\operatorname{rank}(S)$.

向量族的秩可以看成是向量族线性无关程度的度量.

\noindent\textbf{定义 3.5.3}\quad 若向量组 $A$ 和 $B$ 可以互相线性表示，则称这两个向量组等价.

\noindent\textbf{定理 3.5.2}\quad 等价的向量组有相同的秩.

\noindent\textbf{证明}\quad 设 $A_1$ 和 $B_1$ 分别是 $A$ 和 $B$ 的极大无关组，$A_1$ 有 $r$ 个向量，$B_1$ 有 $s$ 个向量，则 $\mathrm{r}(A)=r,\mathrm{r}(B)=s$. 因为 $B_1$ 中向量均可用 $A$ 中向量线性表示，而 $A$ 中向量均可用 $A_1$ 中向量线性表示，故由定理 3.4.4，$B_1$ 中向量均可用 $A_1$ 中向量线性表示，再由引理 3.5.1 可得 $s\leq r$. 同理可证 $r\leq s$，于是 $r=s$. $\square$

如果我们考虑的向量族是整个的线性空间，其极大无关组就是所谓的基.

\noindent\textbf{定义 3.5.4}\quad 设 $V$ 是数域 $\mathbb{K}$ 上的线性空间，若在 $V$ 中存在线性无关的向量 $e_1,e_2,\cdots,e_n$，使得 $V$ 中任一向量均可表示为这组向量的线性组合，则称 $\{e_1,e_2,\cdots,e_n\}$ 是 $V$ 的一组基，线性空间 $V$ 称为 $n$ 维线性空间 (具有维数 $n$). 如果不存在有限个向量组成 $V$ 的一组基，则称 $V$ 是无限维线性空间.

\noindent\textbf{注}\quad 对任一无限维线性空间，也有基的概念. 无限维线性空间基的存在性证明超出了本课程的范围.

显然，$V$ 中的一组极大线性无关组就是 $V$ 的一组基. $n$ 维线性空间任一组基都含有 $n$ 个向量. 如果 $V$ 是数域 $\mathbb{K}$ 上的 $n$ 维线性空间，则记为 $\dim_{\mathbb{K}}V=n$.

\noindent\textbf{例 3.5.2}\quad $n$ 维标准单位行向量 $\{e_1=(1,0,\cdots,0),e_2=(0,1,\cdots,0),\cdots,e_n=(0,0,\cdots,1)\}$ 是 $n$ 维行向量空间 $\mathbb{K}^n$ 的一组基 (证明见例 3.4.2 和例 3.4.4)，因此 $\dim_{\mathbb{K}}\mathbb{K}^n=n$. 同理，$n$ 维标准单位列向量 $\{e'_1,e'_2,\cdots,e'_n\}$ 是 $n$ 维列向量空间 $\mathbb{K}_n$ 的一组基，因此 $\dim_{\mathbb{K}}\mathbb{K}_n=n$.

\noindent\textbf{例 3.5.3}\quad 将复数域 $\mathbb{C}$ 看成是实数域 $\mathbb{R}$ 上的线性空间，容易验证 $\{1,\mathrm{i}=\sqrt{-1}\}$ 是一组基，因此 $\dim_{\mathbb{R}}\mathbb{C}=2$.

\noindent\textbf{推论 3.5.1}\quad $n$ 维线性空间 $V$ 中任一超过 $n$ 个向量的向量组必线性相关.

\noindent\textbf{证明}\quad 设 $V$ 的一组基为 $\{e_1,e_2,\cdots,e_n\}$，且 $\alpha_1,\alpha_2,\cdots,\alpha_m$ 为 $V$ 中 $m(m>n)$ 个向量，则 $\alpha_i$ 均可由 $e_1,e_2,\cdots,e_n$ 的线性组合来表示，由引理 3.5.1 的逆否命题知 $\alpha_1,\alpha_2,\cdots,\alpha_m$ 线性相关. $\square$

\noindent\textbf{定理 3.5.3}\quad 设 $V$ 是 $n$ 维线性空间，$e_1,e_2,\cdots,e_n$ 是 $V$ 中 $n$ 个向量. 若它们适合下列条件之一，则 $\{e_1,e_2,\cdots,e_n\}$ 是 $V$ 的一组基.

(1) $e_1,e_2,\cdots,e_n$ 线性无关；

(2) $V$ 中任一向量均可由 $e_1,e_2,\cdots,e_n$ 线性表示.

\noindent\textbf{证明}\quad (1) 因为 $V$ 中任意 $n+1$ 个向量一定线性相关，故对 $V$ 中任一向量 $v$，向量组 $e_1,e_2,\cdots,e_n,v$ 线性相关. 于是存在不全为零的数 $a_1,a_2,\cdots,a_n,c$，使
\[
a_1e_1+a_2e_2+\cdots+a_ne_n+cv=0,
\]
其中 $c\neq 0$. 事实上若 $c=0$，因为 $e_1,e_2,\cdots,e_n$ 线性无关，将导致 $a_1=a_2=\cdots=a_n=c=0$，与假设矛盾. 由 $c\neq 0$ 可得
\[
v=-\frac{a_1}{c}e_1-\frac{a_2}{c}e_2-\cdots-\frac{a_n}{c}e_n.
\]
因此 $v$ 可用向量组 $\{e_1,e_2,\cdots,e_n\}$ 线性表示，即 $\{e_1,e_2,\cdots,e_n\}$ 是 $V$ 的一组基.

(2) 由命题 3.5.1，不妨设 $\{e_1,e_2,\cdots,e_r\}$ 是 $\{e_1,e_2,\cdots,e_n\}$ 的极大无关组，其中 $r\leq n$. 由定理 3.4.4 知 $V$ 中任一向量均可由 $e_1,e_2,\cdots,e_r$ 线性表示，因此 $\{e_1,e_2,\cdots,e_r\}$ 是 $V$ 的一组基. 特别地，$r=\dim V=n$，即 $\{e_1,e_2,\cdots,e_n\}$ 是 $V$ 的一组基. $\square$

\noindent\textbf{定理 3.5.4}\quad 设 $V$ 是 $n$ 维线性空间，$v_1,v_2,\cdots,v_m$ 是 $V$ 中 $m(m<n)$ 个线性无关的向量，又假设 $\{e_1,e_2,\cdots,e_n\}$ 是 $V$ 的一组基，则必可在 $\{e_1,e_2,\cdots,e_n\}$ 中选出 $n-m$ 个向量，使之和 $v_1,v_2,\cdots,v_m$ 一起组成 $V$ 的一组基.

\noindent\textbf{证明}\quad 将 $e_i(i=1,\cdots,n)$ 依次放入 $\{v_1,v_2,\cdots,v_m\}$，则必有一个 $e_i$，使 $v_1,v_2,\cdots,v_m,e_i$ 线性无关. 这是因为若任一 $e_i$ 加入 $v_1,v_2,\cdots,v_m$ 后线性相关，则每个 $e_i$ 可用 $v_1,v_2,\cdots,v_m$ 线性表示，将和引理 3.5.1 的结论矛盾. 现不妨设 $i=m+1$. 若 $m+1<n$，又可从 $e_1,e_2,\cdots,e_n$ 中找到一个向量，加入 $\{v_1,v_2,\cdots,v_m,e_{m+1}\}$ 后仍线性无关. 不断这样做下去，便可将 $v_1,v_2,\cdots,v_m$ 扩张成为 $V$ 的一组基. $\square$

\noindent\textbf{注}\quad 定理 3.5.4 通常称为基扩张定理，常用的形式是：$n$ 维线性空间 $V$ 中任意 $m(m<n)$ 个线性无关的向量均可扩张为 $V$ 的一组基，或 $V$ 的任意一个子空间 (子空间的概念见 §3.7) 的基均可扩张为 $V$ 的一组基.

\subsection*{习题 3.5}

1. 若 $\alpha_1,\alpha_2,\cdots,\alpha_n$ 是线性空间 $V$ 中的向量，且 $V$ 中任一向量均可用唯一的方式表示为 $\alpha_1,\alpha_2,\cdots,\alpha_n$ 的线性组合，求证：$\{\alpha_1,\alpha_2,\cdots,\alpha_n\}$ 是 $V$ 的一组基.

2. 设 $\{\alpha_1,\alpha_2,\cdots,\alpha_n\}$ 是线性空间 $V$ 的一组基，$\beta_1,\beta_2,\cdots,\beta_n$ 是 $V$ 中 $n$ 个向量，若 $\alpha_1,\alpha_2,\cdots,\alpha_n$ 中任一向量均可由 $\beta_1,\beta_2,\cdots,\beta_n$ 线性表示，求证：$\{\beta_1,\beta_2,\cdots,\beta_n\}$ 也是 $V$ 的一组基.

3. 设 $V$ 是数域 $\mathbb{K}$ 上次数不超过 $n$ 的多项式全体构成的实线性空间，求证：$\{1,x,x^2,\cdots,x^n\}$ 是 $V$ 的一组基，并且 $\{1,x+1,(x+1)^2,\cdots,(x+1)^n\}$ 也是 $V$ 的一组基.

4. 设 $V$ 是数域 $\mathbb{K}$ 上次数小于 $n$ 的多项式全体构成的线性空间，$a_1,a_2,\cdots,a_n$ 是 $\mathbb{K}$ 中互不相同的 $n$ 个数，$f(x)=(x-a_1)(x-a_2)\cdots(x-a_n)$，$f_i(x)=f(x)/(x-a_i)$，求证：$\{f_1(x),f_2(x),\cdots,f_n(x)\}$ 组成 $V$ 的一组基.

5. 设 $V$ 是数域 $\mathbb{K}$ 上 $m\times n$ 矩阵组成的线性空间，令 $E_{ij}(1\leq i\leq m,1\leq j\leq n)$ 是第 $(i,j)$ 元素为 1、其余元素为 0 的 $m\times n$ 矩阵，求证：全体 $E_{ij}$ 组成了 $V$ 的一组基，从而 $V$ 是 $mn$ 维线性空间.

6. 设 $V$ 是数域 $\mathbb{K}$ 上 $n$ 阶上三角阵全体组成的线性空间，求 $V$ 的一组基和维数.

7. 设 $V$ 是数域 $\mathbb{K}$ 上 $n$ 阶对称阵全体组成的线性空间，求 $V$ 的一组基和维数.

8. 设 $V$ 是数域 $\mathbb{K}$ 上 $n$ 阶反对称阵全体组成的线性空间，求 $V$ 的一组基和维数.

9. 设 $V_1=\{A\in M_n(\mathbb{C})\mid \overline{A}'=A\}$ 为 $n$ 阶 Hermite 矩阵全体，$V_2=\{A\in M_n(\mathbb{C})\mid \overline{A}'=-A\}$ 为 $n$ 阶斜 Hermite 矩阵全体，求证：在矩阵加法和实数关于矩阵的数乘下，$V_1,V_2$ 成为实数域 $\mathbb{R}$ 上的线性空间，并且具有相同的维数.

10. 设 $V=\{a+b\sqrt[3]{2}+c\sqrt[3]{4}\}$，其中 $a,b,c$ 均是有理数，证明：$V$ 是有理数域上的线性空间并求其维数.

11. 设 $V=\{a+b\sqrt{2}+c\sqrt{3}+d\sqrt{6}\}$，其中 $a,b,c,d$ 均是有理数，证明：$V$ 是有理数域上的线性空间并求其维数.

\section*{§3.6\quad 矩阵的秩}

\noindent\textbf{定义 3.6.1}\quad 设 $A$ 是 $m\times n$ 矩阵，则 $A$ 的 $m$ 个行向量的秩称为 $A$ 的行秩；$A$ 的 $n$ 个列向量的秩称为 $A$ 的列秩.

\noindent\textbf{注}\quad 我们很快将证明矩阵的行秩等于它的列秩.

在 §3.5 开始时，我们从线性方程组引出了向量组秩的概念. 注意到交换线性方程组中的方程式，用非零常数乘以某一个方程式以及某个方程式乘以一个常数加到另外一个方程式上去，这 3 种变换并不改变线性方程组的同解性. 因此我们有理由猜想：矩阵的行秩和列秩在初等变换下是不变的.

\noindent\textbf{定理 3.6.1}\quad 矩阵的行秩与列秩在初等变换下不变.

\noindent\textbf{证明}\quad 我们分两步走. 第一步证明矩阵的行秩在初等行变换下不变，列秩在初等列变换下不变. 第二步证明列秩在初等行变换下不变，行秩在初等列变换下不变.

第一步，设 $A=(a_{ij})_{m\times n}$，为简单起见将它写成行分块的形状：
\[
A=\left(\begin{array}{c}
\alpha_1\\
\alpha_2\\
\vdots\\
\alpha_m
\end{array}\right),
\]
其中 $\alpha_i=(a_{i1},a_{i2},\cdots,a_{in})(i=1,2,\cdots,m)$ 是 $A$ 的第 $i$ 个行向量. 对换 $A$ 的任意两行并不改变 $A$ 的行向量组，因此也不改变 $A$ 的行秩. 这表明 $A$ 的行秩在第一类初等行变换下不变. 又若以一个非零常数 $k$ 乘以 $A$ 的第 $i$ 行，则 $A$ 变成
\[
A_1=\left(\begin{array}{c}
\alpha_1\\
\vdots\\
k\alpha_i\\
\vdots\\
\alpha_m
\end{array}\right).
\]
显然，$A_1$ 的 $m$ 个行向量可用 $A$ 的 $m$ 个行向量的线性组合来表示. 反之 $A$ 的 $m$ 个行向量也可用 $A_1$ 的行向量的线性组合来表示，因此 $A$ 的行秩与 $A_1$ 的行秩相等. 接下来看第三类初等行变换. 将矩阵 $A$ 的第 $i$ 行乘以 $k$ 后加到第 $j$ 行上去，矩阵 $A$ 变成了下列矩阵：
\[
A_2=\left(\begin{array}{c}
\alpha_1\\
\vdots\\
\alpha_i\\
\vdots\\
k\alpha_i+\alpha_j\\
\vdots\\
\alpha_m
\end{array}\right).
\]
显然，$A_2$ 的行向量是 $A$ 的行向量的线性组合. 反之，
\[
\alpha_j=(-k)\alpha_i+(k\alpha_i+\alpha_j).
\]
因此 $A$ 的行向量也是 $A_2$ 的行向量的线性组合，从而 $A$ 与 $A_2$ 的行秩相等. 这就证明了 $A$ 的行秩在初等行变换下不变. 同理，$A$ 的列秩在初等列变换下也不变.

第二步，我们证明 $A$ 的列秩在初等行变换下不变. 由于 $A$ 的初等行变换等价于用一个初等矩阵左乘以 $A$，我们只需证明对任一初等矩阵 $Q$，$QA$ 与 $A$ 的列秩相等就可以了. 现把 $A$ 写成列分块的形状：
\[
A=(\beta_1,\beta_2,\cdots,\beta_n),
\]
其中 $\beta_j$ 是 $A$ 的第 $j$ 个列向量. 由分块矩阵的乘法得
\[
QA=(Q\beta_1,Q\beta_2,\cdots,Q\beta_n).
\]
设 $A$ 的列向量的极大无关组为 $\beta_{j_1},\cdots,\beta_{j_r}$，现在我们证明 $\{Q\beta_{j_1},\cdots,Q\beta_{j_r}\}$ 是 $QA$ 的列向量的极大无关组.

先证明 $Q\beta_{j_1},\cdots,Q\beta_{j_r}$ 线性无关. 设有 $\lambda_1,\lambda_2,\cdots,\lambda_r\in\mathbb{K}$，使
\[
\lambda_1Q\beta_{j_1}+\lambda_2Q\beta_{j_2}+\cdots+\lambda_rQ\beta_{j_r}=0,
\]
则
\[
Q(\lambda_1\beta_{j_1}+\lambda_2\beta_{j_2}+\cdots+\lambda_r\beta_{j_r})=0.
\]
但 $Q$ 是非异阵，在上式两边左乘 $Q^{-1}$ 即得
\[
\lambda_1\beta_{j_1}+\lambda_2\beta_{j_2}+\cdots+\lambda_r\beta_{j_r}=0.
\]
再由 $\beta_{j_1},\cdots,\beta_{j_r}$ 线性无关即得 $\lambda_1=\cdots=\lambda_r=0$. 这就证明了 $Q\beta_{j_1},\cdots,Q\beta_{j_r}$ 是一组线性无关的向量.

再证明任一 $Q\beta_j$ 均可表示为 $Q\beta_{j_1},\cdots,Q\beta_{j_r}$ 的线性组合. 由于 $\beta_{j_1},\cdots,\beta_{j_r}$ 是 $A$ 的列向量的极大无关组，故
\[
\beta_j=\mu_1\beta_{j_1}+\mu_2\beta_{j_2}+\cdots+\mu_r\beta_{j_r}.
\]
上式两边左乘 $Q$ 即得
\[
Q\beta_j=\mu_1Q\beta_{j_1}+\mu_2Q\beta_{j_2}+\cdots+\mu_rQ\beta_{j_r}.
\]
由上面的论证知道 $A$ 与 $QA$ 的列向量的极大无关组都有相同个数的向量，因此 $A$ 与 $QA$ 的列秩相等. 同理可证明 $A$ 的行秩在初等列变换下不变. $\square$

\noindent\textbf{推论 3.6.1}\quad 任一矩阵的行秩等于列秩.

\noindent\textbf{证明}\quad 任一矩阵 $A$ 经初等变换后均可变成下列分块对角阵：
\[
B=\left(\begin{array}{cc}
I_r&O\\
O&O
\end{array}\right),
\]
其中 $B$ 是分块矩阵，$I_r$ 为 $r$ 阶单位阵. 显然，$B$ 的行秩与列秩都等于 $r$，因此 $A$ 的行秩与列秩都等于 $r$. $\square$

有了这个推论，我们今后将矩阵的行秩与列秩统称为矩阵的秩. 矩阵 $A$ 的秩用 $\mathrm{r}(A)$ 或 $\operatorname{rank}(A)$ 来表示.

\noindent\textbf{推论 3.6.2}\quad 设 $A$ 是 $m\times n$ 矩阵且 $A$ 的第 $j_1,\cdots,$ 第 $j_r$ 列向量是 $A$ 的列向量的极大无关组，则对任意的 $m$ 阶非异阵 $Q$，矩阵 $QA$ 的第 $j_1,\cdots,$ 第 $j_r$ 列向量也是 $QA$ 的列向量的极大无关组.

\noindent\textbf{证明}\quad 非异阵 $Q$ 是若干个初等矩阵的乘积，由定理 3.6.1 的证明即得结论. $\square$

\noindent\textbf{命题 3.6.1}\quad 设 $A$ 是阶梯形矩阵，则 $A$ 的秩等于其非零行的个数，且阶梯点所在的列向量是 $A$ 的列向量的极大无关组.

\noindent\textbf{证明}\quad 设阶梯形矩阵 $A$ 有 $r$ 个非零行，其阶梯点依次是 $a_{1k_1},a_{2k_2},\cdots,a_{rk_r}$：
\[
A=\left(\begin{array}{ccccccccc}
0&\cdots&a_{1k_1}&\cdots&\cdots&\cdots&\cdots&\cdots&\cdots\\
0&\cdots&0&\cdots&a_{2k_2}&\cdots&\cdots&\cdots&\cdots\\
\vdots&&\vdots&&\vdots&&\vdots&\vdots&\\
0&\cdots&0&\cdots&0&\cdots&a_{rk_r}&\cdots&\cdots
\end{array}\right).
\]

先用第三类初等列变换以及阶梯点上的元素依次消去同行的其他非零元素；再用第二类初等列变换将阶梯点上的元素全部变成 1；最后用列对换依次将 $r$ 个阶梯点换到第 $(1,1),(2,2),\cdots,(r,r)$ 位置，从而得到相抵标准型：
\[
\left(\begin{array}{cc}
I_r&O\\
O&O
\end{array}\right).
\]
由定理 3.6.1 可得 $\mathrm{r}(A)=r$.

对于第二个结论，将 $r$ 个阶梯点所在的列向量取出，排成一个新的矩阵：
\[
\left(\begin{array}{cccc}
a_{1k_1}&\cdots&\cdots&\cdots\\
0&a_{2k_2}&\cdots&\cdots\\
\vdots&\vdots&&\vdots\\
0&0&\cdots&a_{rk_r}\\
&&O&
\end{array}\right).
\]
利用同样的方法可将此矩阵化为相抵标准型：
\[
\left(\begin{array}{c}
I_r\\
O
\end{array}\right).
\]
因此 $r$ 个阶梯点所在的列向量组的秩等于 $r$，即为 $A$ 的列秩，从而阶梯点所在的列向量是 $A$ 的列向量的极大无关组. $\square$

由定理 3.6.1 以及命题 3.6.1，我们得到求一个矩阵秩的方法：用初等行变换将一个矩阵 $A$ 化为阶梯形矩阵 $B$，则矩阵 $B$ 的非零行的个数就是矩阵 $A$ 的秩.

\noindent\textbf{例 3.6.1}\quad 求下列矩阵的秩以及列向量的极大无关组：
\[
A=\left(\begin{array}{rrrr}
1&2&3&4\\
-1&-1&-1&-1\\
1&3&5&7
\end{array}\right).
\]

\noindent\textbf{解}\quad 通过初等行变换可将 $A$ 化为如下阶梯形矩阵：
\[
\left(\begin{array}{rrrr}
1&2&3&4\\
-1&-1&-1&-1\\
1&3&5&7
\end{array}\right)
\to
\left(\begin{array}{rrrr}
1&2&3&4\\
0&1&2&3\\
0&1&2&3
\end{array}\right)
\to
\left(\begin{array}{rrrr}
1&2&3&4\\
0&1&2&3\\
0&0&0&0
\end{array}\right).
\]
因此 $A$ 的秩为 2. 由命题 3.6.1 和推论 3.6.2 可得 $A$ 的第一列和第二列是其列向量的极大无关组.

\noindent\textbf{推论 3.6.3}\quad 对任意一个秩为 $r$ 的 $m\times n$ 矩阵 $A$，总存在 $m$ 阶非异阵 $P$ 和 $n$ 阶非异阵 $Q$，使得
\begin{equation}
PAQ=\left(\begin{array}{cc}
I_r&O\\
O&O
\end{array}\right).
\tag{3.6.1}
\end{equation}

\noindent\textbf{证明}\quad 由推论 3.6.1 的证明即得结论. $\square$

\noindent\textbf{推论 3.6.4}\quad 任一矩阵 $A$ 的转置 $A'$ 与 $A$ 有相同的秩.

\noindent\textbf{证明}\quad 由推论 3.6.1 即得结论. $\square$

\noindent\textbf{推论 3.6.5}\quad 任一矩阵与非异阵相乘，其秩不变.

\noindent\textbf{证明}\quad 任一非异阵均可化为有限个初等矩阵的积，由此即得结论. $\square$

若 $n$ 阶方阵 $A$ 的秩等于 $n$，则称 $A$ 为满秩阵. 根据矩阵秩的定义，满秩条件等价于 $A$ 的 $n$ 个行向量线性无关，也等价于 $A$ 的 $n$ 个列向量线性无关.

\noindent\textbf{推论 3.6.6}\quad $n$ 阶方阵 $A$ 为非异阵的充分必要条件是 $A$ 为满秩阵.

\noindent\textbf{证明}\quad 若 $A$ 为非异阵，则由推论 3.6.5 可得 $\mathrm{r}(A)=\mathrm{r}(AI_n)=\mathrm{r}(I_n)=n$，即 $A$ 为满秩阵. 若 $A$ 为满秩阵，则由推论 3.6.3 知 $A$ 经过初等变换可化为单位阵 $I_n$，从而 $A$ 为非异阵. $\square$

由这个推论，非异阵又称为满秩阵.

\noindent\textbf{推论 3.6.7}\quad 两个 $m\times n$ 矩阵等价的充分必要条件是它们具有相同的秩.

\noindent\textbf{证明}\quad 设矩阵 $A,B$ 秩都等于 $r$，则它们都等价于 (3.6.1) 式的矩阵，因此 $A$ 和 $B$ 等价. 反之，由于秩在初等变换下不变，从 $A,B$ 等价可知它们的秩相同. $\square$

我们已经知道，一个 $n$ 阶方阵秩等于 $n$ 的充分必要条件是该矩阵的行列式不等于零. 这个结论很容易被推广到秩为 $r$ 的 $m\times n$ 矩阵上. 从推论 3.6.3 中我们发现，在 $PAQ$（它的秩等于 $r$）中，有一个 $r$ 阶子行列式 $|I_r|$ 其值不等于 0，而没有值不等于 0 的超过 $r$ 阶的子行列式. 这个结论对一般的矩阵还对吗？

我们先解释一下子式的概念. 设 $A=(a_{ij})$ 是一个 $m\times n$ 矩阵. 任取 $A$ 的 $k$ 行与 $k$ 列，位于这些行与这些列的交叉处的元素按原来的顺序构成一个 $k$ 阶行列式，称为 $A$ 的一个 $k$ 阶子式.

例如，在矩阵
\[
A=\left(\begin{array}{rrrrr}
1&2&3&1&0\\
-1&0&1&-2&4\\
0&1&-3&5&7
\end{array}\right)
\]
中取第一行、第二行及第一列、第三列得到的二阶子式为
\[
\left|\begin{array}{rr}
1&3\\
-1&1
\end{array}\right|,
\]
取第二行、第三行及第一列、第五列得到的二阶子式为
\[
\left|\begin{array}{rr}
-1&4\\
0&7
\end{array}\right|.
\]

\noindent\textbf{定理 3.6.2}\quad 设 $m\times n$ 矩阵 $A=(a_{ij})$ 有一个 $r$ 阶子式不等于零，且 $A$ 中任意 $r+1$ 阶子式（如存在）都等于零，则 $\mathrm{r}(A)=r$. 反之，若 $\mathrm{r}(A)=r$，则 $A$ 中必有一个 $r$ 阶子式不等于零，而所有 $r+1$ 阶子式都等于零.

\noindent\textbf{证明}\quad 设 $\mathrm{r}(A)=r$，则 $A$ 中任意 $r+1$ 行都线性相关，由 §3.4 习题 12 的结论知 $A$ 的任意 $r+1$ 阶子式的行向量也线性相关，再由推论 3.6.6 可知这些 $r+1$ 阶子式的值均为零. 再证明 $A$ 至少有一个 $r$ 阶子式不等于零. 因为 $A$ 的秩为 $r$，$A$ 中有 $r$ 行线性无关. 不失一般性，设为前 $r$ 行. 把这 $r$ 行取出得到一个矩阵：
\[
B=\left(\begin{array}{cccc}
a_{11}&a_{12}&\cdots&a_{1n}\\
\vdots&\vdots&&\vdots\\
a_{r1}&a_{r2}&\cdots&a_{rn}
\end{array}\right).
\]
显然 $\mathrm{r}(B)=r$，因此 $B$ 有 $r$ 列线性无关，同样不妨设为前 $r$ 列，则由 $B$ 的前 $r$ 列组成的行列式不等于零，即 $A$ 有一个 $r$ 阶子式不等于零.

反之，设 $A$ 有一个 $r$ 阶子式不为零而 $A$ 的所有 $r+1$ 阶子式全等于零，这时由 Laplace 定理可知，$A$ 的所有高于 $r$ 阶的子式均等于零. 设 $\mathrm{r}(A)=t$，则由前面的论述可知 $t\geq r$，否则 $A$ 的 $r$ 阶子式无一不为零. 但 $t$ 也不能大于 $r$，否则 $A$ 就要有一个大于 $r$ 阶的子式不等于零而与假设矛盾，因此 $t=r$. $\square$

\noindent\textbf{例 3.6.2}\quad 设 $C=\left(\begin{array}{cc}A&O\\O&B\end{array}\right)$，求证：$\mathrm{r}(C)=\mathrm{r}(A)+\mathrm{r}(B)$.

\noindent\textbf{证明}\quad 设 $A,B$ 的秩分别为 $r_1,r_2$，则存在可逆阵 $P_1,Q_1$ 和可逆阵 $P_2,Q_2$，使
\[
P_1AQ_1=\left(\begin{array}{cc}
I_{r_1}&O\\
O&O
\end{array}\right),\qquad
P_2BQ_2=\left(\begin{array}{cc}
I_{r_2}&O\\
O&O
\end{array}\right).
\]
于是
\[
\left(\begin{array}{cc}
P_1&O\\
O&P_2
\end{array}\right)
\left(\begin{array}{cc}
A&O\\
O&B
\end{array}\right)
\left(\begin{array}{cc}
Q_1&O\\
O&Q_2
\end{array}\right)
=
\left(\begin{array}{cc}
P_1AQ_1&O\\
O&P_2BQ_2
\end{array}\right)
=
\left(\begin{array}{cccc}
I_{r_1}&O&O&O\\
O&O&O&O\\
O&O&I_{r_2}&O\\
O&O&O&O
\end{array}\right).
\]
因此 $\mathrm{r}(C)=r_1+r_2$. $\square$

\noindent\textbf{例 3.6.3}\quad 求证：$\mathrm{r}(AB)\leq \min\{\mathrm{r}(A),\mathrm{r}(B)\}$.

\noindent\textbf{证明}\quad 设 $A$ 是 $m\times n$ 矩阵，$B$ 是 $n\times s$ 矩阵. 将矩阵 $B$ 按列分块，$B=(\beta_1,\beta_2,\cdots,\beta_s)$，则 $AB=(A\beta_1,A\beta_2,\cdots,A\beta_s)$. 若 $B$ 的列向量的极大无关组为 $\{\beta_{j_1},\beta_{j_2},\cdots,\beta_{j_r}\}$，则 $B$ 的任一列向量 $\beta_j$ 可用 $\{\beta_{j_1},\beta_{j_2},\cdots,\beta_{j_r}\}$ 线性表示. 于是任一 $A\beta_j$ 也可用 $\{A\beta_{j_1},A\beta_{j_2},\cdots,A\beta_{j_r}\}$ 来线性表示. 因此向量组 $\{A\beta_1,A\beta_2,\cdots,A\beta_s\}$ 的秩不超过 $r$，即 $\mathrm{r}(AB)\leq \mathrm{r}(B)$. 同理，对矩阵 $A$ 用行分块的方法可以证明 $\mathrm{r}(AB)\leq \mathrm{r}(A)$. $\square$

\noindent\textbf{例 3.6.4}\quad 求证：$n$ 阶矩阵 $A$ 是幂等阵（即 $A^2=A$）的充分必要条件是
\[
\mathrm{r}(A)+\mathrm{r}(I_n-A)=n.
\]

\noindent\textbf{证明}\quad 在下列矩阵的分块初等变换中矩阵的秩保持不变：
\[
\left(\begin{array}{cc}
A&O\\
O&I-A
\end{array}\right)
\to
\left(\begin{array}{cc}
A&A\\
O&I-A
\end{array}\right)
\to
\left(\begin{array}{cc}
A&A\\
A&I
\end{array}\right)
\to
\left(\begin{array}{cc}
A-A^2&A\\
O&I
\end{array}\right)
\to
\left(\begin{array}{cc}
A-A^2&O\\
O&I
\end{array}\right).
\]
因此
\[
\mathrm{r}\left(\begin{array}{cc}
A&O\\
O&I-A
\end{array}\right)
=
\mathrm{r}\left(\begin{array}{cc}
A-A^2&O\\
O&I
\end{array}\right),
\]
即 $\mathrm{r}(A)+\mathrm{r}(I-A)=\mathrm{r}(A-A^2)+n$. 由此即得结论. $\square$

对矩阵求秩的方法也可以用来求向量组的秩，方法是将向量组拼成一个矩阵，用初等变换求出矩阵的秩. 由于矩阵的秩就是其行向量组或列向量组的秩，因此就得到了向量组的秩. 我们也可以利用命题 3.6.1 和推论 3.6.2 来求向量组的极大无关组，注意此时应将向量组按列分块的方式拼成矩阵，并用初等行变换将矩阵变为阶梯形，这样就可以得到原向量组的极大无关组了. 通常求向量组的秩和极大无关组可以同时进行.

\noindent\textbf{例 3.6.5}\quad 求向量组的秩和极大无关组：
\[
\{(1,0,2),(2,-1,3),(3,-2,4),(4,-3,5)\}.
\]

\noindent\textbf{解}\quad 将上述向量按列分块的方式拼成矩阵：
\[
\left(\begin{array}{rrrr}
1&2&3&4\\
0&-1&-2&-3\\
2&3&4&5
\end{array}\right).
\]
对上述矩阵进行初等行变换化为阶梯形矩阵：
\[
\left(\begin{array}{rrrr}
1&2&3&4\\
0&-1&-2&-3\\
2&3&4&5
\end{array}\right)
\to
\left(\begin{array}{rrrr}
1&2&3&4\\
0&-1&-2&-3\\
0&-1&-2&-3
\end{array}\right)
\to
\left(\begin{array}{rrrr}
1&2&3&4\\
0&-1&-2&-3\\
0&0&0&0
\end{array}\right).
\]
得到矩阵的秩显然为 2，所以向量组的秩也等于 2. 根据阶梯点所在的位置可知向量组的极大无关组为 $\{(1,0,2),(2,-1,3)\}$.

要判断向量组是否线性相关，用定义来做将是一件很麻烦的事. 现在我们可以用求矩阵的秩的方法来判断. 具体来说，我们将要判断的向量组拼成一个矩阵，然后求出矩阵的秩. 若矩阵的秩（实际上也是向量组的秩）等于向量的个数，则向量组线性无关；若秩小于向量的个数，则向量组线性相关.

\noindent\textbf{例 3.6.6}\quad 判定下列向量组是否线性无关：
\[
\{(-1,3,1),(2,1,0),(1,4,1)\}.
\]

\noindent\textbf{解}\quad 将向量组拼成矩阵并用初等变换求秩：
\[
\left(\begin{array}{rrr}
-1&3&1\\
2&1&0\\
1&4&1
\end{array}\right)
\to
\left(\begin{array}{rrr}
-1&3&1\\
0&7&2\\
0&7&2
\end{array}\right)
\to
\left(\begin{array}{rrr}
-1&3&1\\
0&7&2\\
0&0&0
\end{array}\right).
\]
矩阵的秩等于 2，小于向量组中向量的个数，因此向量组线性相关.

\noindent\textbf{注}\quad 如果用向量组拼成的矩阵是一个方阵，则也可用行列式法来判断这个矩阵的秩是否等于向量组中向量的个数. 比如计算出上面矩阵的行列式后我们发现它等于零，因此矩阵的秩小于 3，即向量组线性相关.

\noindent\textbf{例 3.6.7}\quad 判断下列向量是否线性相关：
\[
(1,2,-1),(0,2,2),(2,1,3).
\]

\noindent\textbf{解}\quad 计算由这 3 个向量组成的矩阵的行列式
\[
\left|\begin{array}{rrr}
1&2&-1\\
0&2&2\\
2&1&3
\end{array}\right|
=
\left|\begin{array}{rrr}
1&2&-1\\
0&2&2\\
0&-3&5
\end{array}\right|
=16\neq 0.
\]
所以这 3 个向量线性无关.

\subsection*{习题 3.6}

1. 用初等变换法求下列矩阵的秩：
\[
\left(\begin{array}{rrr}
1&2&0\\
0&1&1\\
-1&2&3
\end{array}\right);\quad
\left(\begin{array}{rrr}
1&2&3\\
0&-1&-1\\
3&4&7
\end{array}\right);\quad
\left(\begin{array}{rrrr}
1&2&3&4\\
2&3&4&1\\
3&4&1&2\\
4&1&2&3
\end{array}\right);\quad
\left(\begin{array}{rrrr}
-1&2&1&0\\
1&-2&-1&0\\
-1&0&1&1\\
-2&0&2&2
\end{array}\right).
\]

2. 用矩阵的初等变换法求下列向量组的秩：

(1) $(2,1,3,0,4),(-1,2,3,1,0),(3,-1,0,-1,4)$；

(2) $(1,2,3,4),(0,-1,2,3),(2,3,8,11),(2,3,6,8)$.

3. 用矩阵的初等变换法判断下列向量组是否线性相关：

(1) $(5,1,2),(-3,1,3),(2,2,3)$；

(2) $(1,2,3),(3,6,9),(2,1,1)$；

(3) $(1,-2,2,3),(-2,4,-1,3),(0,6,2,3)$.

4. 已知向量组 $\{\alpha_1=(1,2,3,4),\alpha_2=(2,3,4,5),\alpha_3=(3,4,5,6),\alpha_4=(4,5,6,7)\}$，用矩阵的初等变换法求该向量组的一个极大无关组.

5. 证明下列矩阵秩的公式：

(1) 若 $k\neq 0$，则 $\mathrm{r}(kA)=\mathrm{r}(A)$；

(2) $\mathrm{r}\left(\begin{array}{cc}A&C\\O&B\end{array}\right)\geq \mathrm{r}(A)+\mathrm{r}(B)$，$\mathrm{r}\left(\begin{array}{cc}A&O\\D&B\end{array}\right)\geq \mathrm{r}(A)+\mathrm{r}(B)$；

(3) $\mathrm{r}(A\ B)\leq \mathrm{r}(A)+\mathrm{r}(B)$，$\mathrm{r}\left(\begin{array}{c}A\\B\end{array}\right)\leq \mathrm{r}(A)+\mathrm{r}(B)$；

(4) $\mathrm{r}(A+B)\leq \mathrm{r}(A)+\mathrm{r}(B)$，$\mathrm{r}(A-B)\leq \mathrm{r}(A)+\mathrm{r}(B)$；

(5) $\mathrm{r}(A-B)\geq |\mathrm{r}(A)-\mathrm{r}(B)|$.

6. 证明：一个矩阵添加一行或一列，其秩不变或增加 1.

7. 设 $A$ 是 $m\times n$ 矩阵，$B$ 是 $n\times t$ 矩阵，求证：
\[
\mathrm{r}(AB)\geq \mathrm{r}(A)+\mathrm{r}(B)-n.
\]

8. 求证：$n$ 阶矩阵 $A$ 是对合阵（即 $A^2=I_n$）的充分必要条件是
\[
\mathrm{r}(I_n+A)+\mathrm{r}(I_n-A)=n.
\]

9. 设 $A$ 是 $n$ 阶矩阵，求证：$\mathrm{r}(A)+\mathrm{r}(I_n+A)\geq n$.

10. 设 $A$ 是 $m\times n$ 矩阵，求证：

(1) 若 $\mathrm{r}(A)=n$，则必存在秩为 $n$ 的 $n\times m$ 矩阵 $B$，使得 $BA=I_n$；

(2) 若 $\mathrm{r}(A)=m$，则必存在秩为 $m$ 的 $n\times m$ 矩阵 $C$，使得 $AC=I_m$.

11. 设 $m\times n$ 矩阵 $A$ 的秩为 $r$，证明：

(1) $A=BC$，其中 $B$ 是 $m\times r$ 矩阵且 $\mathrm{r}(B)=r$，$C$ 是 $r\times n$ 矩阵且 $\mathrm{r}(C)=r$，这种分解称为 $A$ 的满秩分解；

(2) 若 $A$ 有两个满秩分解 $A=B_1C_1=B_2C_2$，则存在 $r$ 阶非异阵 $P$，使得 $B_2=B_1P$，$C_2=P^{-1}C_1$.

\section*{§3.7\quad 坐标向量}

读者已经学过解析几何. 解析几何就是用代数工具来研究几何问题，做到这一点最根本的是要建立坐标系，将平面（或空间）上的点和有序实数组对应起来. 我们在线性空间中引进基的目的就是为了要在其中引进“坐标”.

\noindent\textbf{引理 3.7.1}\quad 设 $\{e_1,e_2,\cdots,e_n\}$ 是 $n$ 维线性空间 $V$ 的一组基，且
\[
\alpha=a_1e_1+a_2e_2+\cdots+a_ne_n=b_1e_1+b_2e_2+\cdots+b_ne_n,
\]
则 $a_1=b_1,a_2=b_2,\cdots,a_n=b_n$.

\noindent\textbf{证明}\quad 由假设得
\[
(a_1-b_1)e_1+(a_2-b_2)e_2+\cdots+(a_n-b_n)e_n=0.
\]
但 $e_1,e_2,\cdots,e_n$ 线性无关，因此 $a_i-b_i=0$，即 $a_i=b_i(i=1,2,\cdots,n)$. $\square$

这个引理表明，如果取定 $V$ 中的一组基，则 $V$ 中任一向量可以而且只可以用一种方式表示为 $e_1,e_2,\cdots,e_n$ 的线性组合. 如果我们固定基向量的次序为 $\{e_1,e_2,\cdots,e_n\}$，则 $\alpha$ 唯一地对应 $\mathbb{K}$ 中的一组有序数 $(a_1,a_2,\cdots,a_n)$. 我们称这组有序数为 $\alpha$ 在基 $\{e_1,e_2,\cdots,e_n\}$ 下的坐标向量，其中 $a_i$ 称为第 $i$ 个坐标. 反过来，$\mathbb{K}$ 中的任一组有序的 $n$ 个数 $(a_1,a_2,\cdots,a_n)$ 也唯一地对应 $V$ 中的一个向量 $a_1e_1+a_2e_2+\cdots+a_ne_n$. 我们把坐标向量看成是 $n$ 维行向量，于是在 $V$ 与 $\mathbb{K}^n$ 之间存在一个一一对应的映射 $\varphi$：
\[
a_1e_1+a_2e_2+\cdots+a_ne_n\mapsto (a_1,a_2,\cdots,a_n).
\]

我们知道，$V$ 不仅是一个集合，而且在其中存在着一个代数结构，即向量之间有加减法与数乘. 同样的道理，$\mathbb{K}^n$ 也是一个向量空间，其向量之间也有运算关系. 我们希望知道上述一一对应和这两个线性空间中向量的运算有着怎样的联系.

先来看加法. 设
\[
\alpha=a_1e_1+a_2e_2+\cdots+a_ne_n,\qquad
\beta=b_1e_1+b_2e_2+\cdots+b_ne_n,
\]
则
\[
\alpha+\beta=(a_1+b_1)e_1+(a_2+b_2)e_2+\cdots+(a_n+b_n)e_n.
\]
于是 $\alpha+\beta$ 对应于
\[
(a_1+b_1,a_2+b_2,\cdots,a_n+b_n).
\]
当我们把 $(a_1,a_2,\cdots,a_n)$ 看成行向量空间 $\mathbb{K}^n$ 中的向量时，
\[
(a_1+b_1,a_2+b_2,\cdots,a_n+b_n)=(a_1,a_2,\cdots,a_n)+(b_1,b_2,\cdots,b_n),
\]
因此
\[
\varphi(\alpha+\beta)=\varphi(\alpha)+\varphi(\beta).
\]
这里需要注意的是 $\alpha+\beta$ 中的“$+$”是在 $V$ 中的加法，而 $\varphi(\alpha)+\varphi(\beta)$ 中的“$+$”是 $\mathbb{K}^n$ 中的加法.

另一方面，若 $k\in\mathbb{K}$，则
\[
k\alpha=ka_1e_1+ka_2e_2+\cdots+ka_ne_n,
\]
因此
\[
\varphi(k\alpha)=(ka_1,ka_2,\cdots,ka_n)=k(a_1,a_2,\cdots,a_n)=k\varphi(\alpha).
\]
上面的分析表明，当我们在 $V$ 中引进一组基后，可以建立起 $V$ 中的向量与 $\mathbb{K}$ 上的 $n$ 维行向量之间的一一对应. 这个对应保持了线性运算，即
\[
\varphi(\alpha+\beta)=\varphi(\alpha)+\varphi(\beta);\quad \varphi(k\alpha)=k\varphi(\alpha).
\]

\noindent\textbf{定义 3.7.1}\quad 设 $V,U$ 是数域 $\mathbb{K}$ 上的两个线性空间，若存在 $V$ 到 $U$ 上的一个一一对应的映射 $\varphi$，使得对任意 $V$ 中向量 $\alpha,\beta$ 以及 $\mathbb{K}$ 中的数 $k$，均有
\[
\varphi(\alpha+\beta)=\varphi(\alpha)+\varphi(\beta);\quad \varphi(k\alpha)=k\varphi(\alpha),
\]
则称 $V$ 与 $U$ 这两个线性空间同构，记为 $V\cong U$.

上面的论证可归结为下列定理.

\noindent\textbf{定理 3.7.1}\quad 数域 $\mathbb{K}$ 上的任一 $n$ 维线性空间 $V$ 均与 $\mathbb{K}$ 上的 $n$ 维行向量空间 $\mathbb{K}^n$ 同构.

同构的线性空间顾名思义其代数结构是相同的，那么它们中向量的线性关系应该是一致的. 事实上同构的线性空间有如下定理所表述的性质.

\noindent\textbf{定理 3.7.2}\quad (1) 设 $\varphi:V\to U$ 为线性空间的同构，则
\[
\varphi(0)=0.
\]

(2) $\varphi$ 将线性相关的向量组映成线性相关的向量组，将线性无关的向量组映成线性无关的向量组.

(3) 同构关系是一个等价关系，即

(i) $V\cong V$；

(ii) 若 $V\cong U$，则 $U\cong V$；

(iii) 若 $V\cong U,U\cong W$，则 $V\cong W$.

(4) 数域 $\mathbb{K}$ 上的两个有限维线性空间同构的充分必要条件是它们具有相同的维数.

\noindent\textbf{证明}\quad (1) 显然有
\[
\varphi(0)=\varphi(0+0)=\varphi(0)+\varphi(0).
\]
消去一个 $\varphi(0)$ 就有 $\varphi(0)=0$.

(2) 若 $\alpha_1,\alpha_2,\cdots,\alpha_m$ 是 $V$ 中线性相关的向量，则存在不全为零的数 $k_1,k_2,\cdots,k_m$，使
\[
k_1\alpha_1+k_2\alpha_2+\cdots+k_m\alpha_m=0,
\]
于是由 (1)，有
\[
\varphi(k_1\alpha_1+k_2\alpha_2+\cdots+k_m\alpha_m)=0.
\]
上式左端为
\[
\varphi(k_1\alpha_1)+\varphi(k_2\alpha_2)+\cdots+\varphi(k_m\alpha_m)=k_1\varphi(\alpha_1)+k_2\varphi(\alpha_2)+\cdots+k_m\varphi(\alpha_m).
\]
因此 $\varphi(\alpha_1),\varphi(\alpha_2),\cdots,\varphi(\alpha_m)$ 是一组线性相关的向量.

另一方面，若 $\alpha_1,\alpha_2,\cdots,\alpha_m$ 线性无关且如果 $\varphi(\alpha_1),\varphi(\alpha_2),\cdots,\varphi(\alpha_m)$ 在 $U$ 中线性相关，则存在一组不全为零的数 $c_1,c_2,\cdots,c_m$，使
\[
c_1\varphi(\alpha_1)+c_2\varphi(\alpha_2)+\cdots+c_m\varphi(\alpha_m)=0.
\]
但上式左端等于
\[
\varphi(c_1\alpha_1+c_2\alpha_2+\cdots+c_m\alpha_m).
\]
由于 $\varphi$ 是一一对应的映射且已证明 $V$ 中的零向量与 $U$ 中的零向量对应，因此
\[
c_1\alpha_1+c_2\alpha_2+\cdots+c_m\alpha_m=0.
\]
但 $\alpha_1,\alpha_2,\cdots,\alpha_m$ 线性无关，这就引出了矛盾，故 $\varphi(\alpha_1),\varphi(\alpha_2),\cdots,\varphi(\alpha_m)$ 必线性无关.

(3) (i) 显然 $V$ 与自身同构，这时取 $\varphi$ 为恒同映射，即
\[
\varphi(\alpha)=\alpha,\quad \alpha\in V.
\]

(ii) 设 $\varphi$ 是 $V\to U$ 上的一一对应，$\varphi^{-1}$ 是其逆对应：$U\to V$，则 $\varphi^{-1}$ 也是一一对应. 设 $x,y$ 是 $U$ 中的向量，由于 $\varphi$ 是一一对应，故存在 $\alpha,\beta\in V$，使
\[
\varphi(\alpha)=x,\quad \varphi(\beta)=y,
\]
也就是
\[
\alpha=\varphi^{-1}(x),\quad \beta=\varphi^{-1}(y).
\]
由 $\varphi$ 是同构可知
\[
\varphi(\alpha+\beta)=x+y,\quad \varphi(k\alpha)=kx,
\]
因此
\[
\varphi^{-1}(x+y)=\alpha+\beta=\varphi^{-1}(x)+\varphi^{-1}(y),
\]
\[
\varphi^{-1}(kx)=k\alpha=k\varphi^{-1}(x).
\]
这表明 $\varphi^{-1}$ 是 $U\to V$ 上的同构，故 $U\cong V$.

(iii) 若 $\varphi$ 是 $V\to U$ 上的同构，$\psi$ 是 $U\to W$ 上的同构. 令 $\xi=\psi\varphi$，即对任意的 $\alpha\in V$，有
\[
\xi(\alpha)=\psi(\varphi(\alpha)),
\]
则 $\xi$ 是 $V\to W$ 上的一一对应，且
\[
\begin{aligned}
\xi(\alpha+\beta)&=\psi(\varphi(\alpha+\beta))=\psi(\varphi(\alpha)+\varphi(\beta))\\
&=\psi(\varphi(\alpha))+\psi(\varphi(\beta))\\
&=\xi(\alpha)+\xi(\beta),
\end{aligned}
\]
\[
\begin{aligned}
\xi(k\alpha)&=\psi(\varphi(k\alpha))=\psi(k\varphi(\alpha))\\
&=k\psi(\varphi(\alpha))=k\xi(\alpha).
\end{aligned}
\]
这就是说 $\xi$ 是 $V\to W$ 上的同构.

(4) 设 $V$ 与 $U$ 是 $\mathbb{K}$ 上的两个线性空间，$V\cong U$. 设 $\dim V=n$，$\{e_1,e_2,\cdots,e_n\}$ 是 $V$ 的一组基，则 $\{\varphi(e_1),\varphi(e_2),\cdots,\varphi(e_n)\}$ 是 $U$ 的一组线性无关的向量. 又若 $x\in U$，则由于 $\varphi$ 是一一对应，因此存在 $\alpha\in V$，使 $x=\varphi(\alpha)$. 设
\[
\alpha=k_1e_1+k_2e_2+\cdots+k_ne_n,
\]
则
\[
x=\varphi(\alpha)=k_1\varphi(e_1)+k_2\varphi(e_2)+\cdots+k_n\varphi(e_n).
\]
即 $U$ 中任一向量可表示为 $\varphi(e_1),\varphi(e_2),\cdots,\varphi(e_n)$ 的线性组合，故 $\{\varphi(e_1),\varphi(e_2),\cdots,\varphi(e_n)\}$ 是 $U$ 的一组基. 因此 $\dim U=n=\dim V$.

反之，若 $\dim U=\dim V=n$，则 $U$ 与 $V$ 皆同构于 $n$ 维行向量空间 $\mathbb{K}^n$，由 (3) 可知 $V$ 和 $U$ 同构. $\square$

由同构的定义我们可以看出，两个同构的线性空间不仅元素之间有一个一一对应关系，而且这个对应保持了线性关系. 因此，在一个线性空间中由线性关系获得的性质在与之同构的线性空间中也成立. 又由上述定理知道，数域 $\mathbb{K}$ 上的任一 $n$ 维线性空间同构于 $\mathbb{K}^n$. $\mathbb{K}^n$ 是 $\mathbb{K}$ 上的行向量空间，它比较具体，容易捉摸，它是一般 $n$ 维线性空间的“模型”. 我们常常通过对 $\mathbb{K}^n$ 的研究来探讨一般 $n$ 维线性空间的性质. 显然，上述讨论也适用于 $n$ 维列向量空间，即 $\mathbb{K}_n$ 也是一个合适的模型. 我们可视讨论的方便采用合适的模型. 在本书中，我们将比较多地采用列向量空间 $\mathbb{K}_n$. 如果 $\alpha=a_1e_1+a_2e_2+\cdots+a_ne_n$，我们称列向量
\[
\left(\begin{array}{c}
a_1\\
a_2\\
\vdots\\
a_n
\end{array}\right)
\]
为向量 $\alpha$ 在基 $\{e_1,e_2,\cdots,e_n\}$ 下的坐标向量.

既然同构保持了向量的线性关系，那么向量组的秩在同构关系下也应该保持，即我们有下面的推论.

\noindent\textbf{推论 3.7.1}\quad 设 $\{e_1,e_2,\cdots,e_n\}$ 是线性空间 $V$ 的基，$\alpha_1,\alpha_2,\cdots,\alpha_m$ 是 $V$ 中的向量. 它们在这组基下的坐标向量依次为 $\widetilde{\alpha}_1,\widetilde{\alpha}_2,\cdots,\widetilde{\alpha}_m$，则向量组 $\widetilde{\alpha}_1,\widetilde{\alpha}_2,\cdots,\widetilde{\alpha}_m$ 和向量组 $\alpha_1,\alpha_2,\cdots,\alpha_m$ 有相同的秩.

\noindent\textbf{证明}\quad 同构映射将 $\{\alpha_1,\alpha_2,\cdots,\alpha_m\}$ 的极大无关组映为 $\{\widetilde{\alpha}_1,\widetilde{\alpha}_2,\cdots,\widetilde{\alpha}_m\}$ 的极大无关组，所以这两组向量的秩相同. $\square$

这个推论可以将一般的线性空间中向量组的求秩问题归结为行向量组或列向量组的求秩问题，我们已经知道后者可以用矩阵来处理. 特别，判断向量组是否线性相关也可以这样做.

\noindent\textbf{例 3.7.1}\quad 设 $\{e_1,e_2,e_3\}$ 是线性空间 $V$ 的基，又
\[
\left\{\begin{array}{l}
\alpha_1=e_1+2e_2+3e_3,\\
\alpha_2=2e_1-e_2-3e_3,\\
\alpha_3=e_1-3e_2-6e_3,
\end{array}\right.
\]
求向量组 $\{\alpha_1,\alpha_2,\alpha_3\}$ 的秩并判断它们是否线性相关.

\noindent\textbf{解}\quad 因为矩阵
\[
\left(\begin{array}{rrr}
1&2&3\\
2&-1&-3\\
1&-3&-6
\end{array}\right)
\]
的秩等于 2，所以向量组 $\{\alpha_1,\alpha_2,\alpha_3\}$ 的秩等于 2. 这 3 个向量线性相关.

\subsection*{习题 3.7}

1. 求向量 $\alpha=(a_1,a_2,\cdots,a_n)$ 在基
\[
\{\beta_1=(1,1,\cdots,1,1),\beta_2=(1,1,\cdots,1,0),\cdots,\beta_n=(1,0,\cdots,0,0)\}
\]
下的坐标.

2. 设 $\{e_1,e_2,\cdots,e_n\}$ 是线性空间 $V$ 的一组基，问：$\{e_1,e_1+e_2,\cdots,e_1+e_2+\cdots+e_n\}$ 是否也是 $V$ 的基？

3. 已知向量组 $\{\alpha_1,\alpha_2,\cdots,\alpha_s\}(s>1)$ 是线性空间 $V$ 的一组基，设 $\beta_1=\alpha_1+\alpha_2,\beta_2=\alpha_2+\alpha_3,\cdots,\beta_s=\alpha_s+\alpha_1$. 讨论向量 $\beta_1,\beta_2,\cdots,\beta_s$ 的线性相关性.

4. 设 $\{e_1,e_2,e_3,e_4\}$ 是线性空间 $V$ 的一组基，已知
\[
\left\{\begin{array}{l}
\alpha_1=e_1+e_2+e_3+3e_4,\\
\alpha_2=-e_1-3e_2+5e_3+e_4,\\
\alpha_3=3e_1+2e_2-e_3+4e_4,\\
\alpha_4=-2e_1-6e_2+10e_3+2e_4,
\end{array}\right.
\]
求 $\alpha_1,\alpha_2,\alpha_3,\alpha_4$ 的一个极大无关组.

5. 设 $a_i(i=1,2,\cdots,n)$ 是 $n$ 个不同的数，$\{e_1,e_2,\cdots,e_n\}$ 是线性空间 $V$ 的一组基，已知
\[
\left\{\begin{array}{l}
\alpha_1=e_1+a_1e_2+\cdots+a_1^{n-1}e_n,\\
\alpha_2=e_1+a_2e_2+\cdots+a_2^{n-1}e_n,\\
\cdots\cdots\cdots\cdots\\
\alpha_n=e_1+a_ne_2+\cdots+a_n^{n-1}e_n,
\end{array}\right.
\]
求证：$\{\alpha_1,\alpha_2,\cdots,\alpha_n\}$ 也是 $V$ 的一组基.

\section*{§3.8\quad 基变换与过渡矩阵}

我们在上一节中引进了基与坐标的概念. 我们将在这一节中考虑这样一个问题：如果线性空间的基发生了变动，同一个向量的坐标将发生怎样的变化？

\noindent\textbf{定义 3.8.1}\quad 设 $\{e_1,e_2,\cdots,e_n\}$ 是数域 $\mathbb{K}$ 上线性空间 $V$ 的一组基，$\{f_1,f_2,\cdots,f_n\}$ 是另一组基，则 $f_1,f_2,\cdots,f_n$ 可用 $e_1,e_2,\cdots,e_n$ 的下列线性组合表示：
\begin{equation}
\left\{\begin{array}{l}
f_1=a_{11}e_1+a_{12}e_2+\cdots+a_{1n}e_n,\\
f_2=a_{21}e_1+a_{22}e_2+\cdots+a_{2n}e_n,\\
\cdots\cdots\cdots\\
f_n=a_{n1}e_1+a_{n2}e_2+\cdots+a_{nn}e_n.
\end{array}\right.
\tag{3.8.1}
\end{equation}
上述表示式中 $e_i$ 的系数组成了一个元素在 $\mathbb{K}$ 上的 $n$ 阶矩阵，这个矩阵的转置
\[
A=\left(\begin{array}{cccc}
a_{11}&a_{21}&\cdots&a_{n1}\\
a_{12}&a_{22}&\cdots&a_{n2}\\
\vdots&\vdots&&\vdots\\
a_{1n}&a_{2n}&\cdots&a_{nn}
\end{array}\right)
\]
称为从基 $\{e_1,e_2,\cdots,e_n\}$ 到基 $\{f_1,f_2,\cdots,f_n\}$ 的过渡矩阵.

现设
\[
\alpha=\lambda_1e_1+\lambda_2e_2+\cdots+\lambda_ne_n=\mu_1f_1+\mu_2f_2+\cdots+\mu_nf_n,
\]
将 (3.8.1) 式代入得
\[
\begin{aligned}
\alpha&=\mu_1\left(\sum_{j=1}^n a_{1j}e_j\right)+\mu_2\left(\sum_{j=1}^n a_{2j}e_j\right)+\cdots+\mu_n\left(\sum_{j=1}^n a_{nj}e_j\right)\\
&=\left(\sum_{i=1}^n \mu_i a_{i1}\right)e_1+\left(\sum_{i=1}^n \mu_i a_{i2}\right)e_2+\cdots+\left(\sum_{i=1}^n \mu_i a_{in}\right)e_n.
\end{aligned}
\]
由于 $\alpha=\lambda_1e_1+\lambda_2e_2+\cdots+\lambda_ne_n$，故
\begin{equation}
\lambda_j=\mu_1a_{1j}+\mu_2a_{2j}+\cdots+\mu_na_{nj}\quad (j=1,2,\cdots,n).
\tag{3.8.2}
\end{equation}
上式可用矩阵来表示：
\begin{equation}
\left(\begin{array}{c}
\lambda_1\\
\lambda_2\\
\vdots\\
\lambda_n
\end{array}\right)
=
\left(\begin{array}{cccc}
a_{11}&a_{21}&\cdots&a_{n1}\\
a_{12}&a_{22}&\cdots&a_{n2}\\
\vdots&\vdots&&\vdots\\
a_{1n}&a_{2n}&\cdots&a_{nn}
\end{array}\right)
\left(\begin{array}{c}
\mu_1\\
\mu_2\\
\vdots\\
\mu_n
\end{array}\right).
\tag{3.8.3}
\end{equation}
上式表明了同一个向量在不同基下的坐标向量之间的关系.

反之，若 $V$ 中向量 $\alpha$ 在基 $\{e_1,e_2,\cdots,e_n\}$ 下的坐标向量 $(\lambda_1,\lambda_2,\cdots,\lambda_n)'$ 和在基 $\{f_1,f_2,\cdots,f_n\}$ 下的坐标向量 $(\mu_1,\mu_2,\cdots,\mu_n)'$ 适合如 (3.8.3) 式的关系，注意到 $f_1$ 在 $\{f_1,f_2,\cdots,f_n\}$ 下的坐标向量为 $(1,0,\cdots,0)'$，由 (3.8.3) 式可知它在 $\{e_1,e_2,\cdots,e_n\}$ 下的坐标向量为 $(a_{11},a_{12},\cdots,a_{1n})'$，这即是说
\[
f_1=a_{11}e_1+a_{12}e_2+\cdots+a_{1n}e_n.
\]
同理，对 $f_i$ 有
\[
f_i=a_{i1}e_1+a_{i2}e_2+\cdots+a_{in}e_n.
\]
因此，矩阵 $A$ 就是从基 $\{e_1,e_2,\cdots,e_n\}$ 到基 $\{f_1,f_2,\cdots,f_n\}$ 的过渡矩阵.

我们已经知道如果 (3.8.1) 式中的系数矩阵（或等价于它的转置即过渡矩阵 $A$）可逆，那么向量组 $\{f_1,f_2,\cdots,f_n\}$ 必线性无关，从而是线性空间 $V$ 的一组基. 那么反过来，过渡矩阵它是否一定可逆？另外一个问题是：如果已知从基 $\{e_1,e_2,\cdots,e_n\}$ 到基 $\{f_1,f_2,\cdots,f_n\}$ 的过渡矩阵为 $A$，则从 $\{f_1,f_2,\cdots,f_n\}$ 到 $\{e_1,e_2,\cdots,e_n\}$ 的过渡矩阵是什么？我们有理由猜想：过渡矩阵必是可逆阵，且从 $\{f_1,f_2,\cdots,f_n\}$ 到 $\{e_1,e_2,\cdots,e_n\}$ 的过渡矩阵就是 $A^{-1}$. 现设
\begin{equation}
\left\{\begin{array}{l}
e_1=b_{11}f_1+b_{12}f_2+\cdots+b_{1n}f_n,\\
e_2=b_{21}f_1+b_{22}f_2+\cdots+b_{2n}f_n,\\
\cdots\cdots\cdots\\
e_n=b_{n1}f_1+b_{n2}f_2+\cdots+b_{nn}f_n.
\end{array}\right.
\tag{3.8.4}
\end{equation}
令
\[
B=\left(\begin{array}{cccc}
b_{11}&b_{21}&\cdots&b_{n1}\\
b_{12}&b_{22}&\cdots&b_{n2}\\
\vdots&\vdots&&\vdots\\
b_{1n}&b_{2n}&\cdots&b_{nn}
\end{array}\right),
\]
即矩阵 $B$ 是从基 $\{f_1,f_2,\cdots,f_n\}$ 到基 $\{e_1,e_2,\cdots,e_n\}$ 的过渡矩阵.

\noindent\textbf{定理 3.8.1}\quad 上述 $n$ 阶方阵 $A,B$ 互为逆阵，即
\[
B=A^{-1}.
\]

\noindent\textbf{证明}\quad 设
\[
\alpha=\lambda_1e_1+\lambda_2e_2+\cdots+\lambda_ne_n=\mu_1f_1+\mu_2f_2+\cdots+\mu_nf_n,
\]
则
\[
\left(\begin{array}{c}
\lambda_1\\
\lambda_2\\
\vdots\\
\lambda_n
\end{array}\right)
=
\left(\begin{array}{cccc}
a_{11}&a_{21}&\cdots&a_{n1}\\
a_{12}&a_{22}&\cdots&a_{n2}\\
\vdots&\vdots&&\vdots\\
a_{1n}&a_{2n}&\cdots&a_{nn}
\end{array}\right)
\left(\begin{array}{c}
\mu_1\\
\mu_2\\
\vdots\\
\mu_n
\end{array}\right),
\]
\[
\left(\begin{array}{c}
\mu_1\\
\mu_2\\
\vdots\\
\mu_n
\end{array}\right)
=
\left(\begin{array}{cccc}
b_{11}&b_{21}&\cdots&b_{n1}\\
b_{12}&b_{22}&\cdots&b_{n2}\\
\vdots&\vdots&&\vdots\\
b_{1n}&b_{2n}&\cdots&b_{nn}
\end{array}\right)
\left(\begin{array}{c}
\lambda_1\\
\lambda_2\\
\vdots\\
\lambda_n
\end{array}\right).
\]
因此
\[
\left(\begin{array}{c}
\lambda_1\\
\lambda_2\\
\vdots\\
\lambda_n
\end{array}\right)
=
\left(\begin{array}{cccc}
a_{11}&a_{21}&\cdots&a_{n1}\\
a_{12}&a_{22}&\cdots&a_{n2}\\
\vdots&\vdots&&\vdots\\
a_{1n}&a_{2n}&\cdots&a_{nn}
\end{array}\right)
\left(\begin{array}{cccc}
b_{11}&b_{21}&\cdots&b_{n1}\\
b_{12}&b_{22}&\cdots&b_{n2}\\
\vdots&\vdots&&\vdots\\
b_{1n}&b_{2n}&\cdots&b_{nn}
\end{array}\right)
\left(\begin{array}{c}
\lambda_1\\
\lambda_2\\
\vdots\\
\lambda_n
\end{array}\right).
\]
但 $\lambda_i$ 可取 $\mathbb{K}$ 中任意数，因此
\[
AB=I_n. \quad \square
\]

接下去一个问题是：我们假设从基 $\{e_1,e_2,\cdots,e_n\}$ 到基 $\{f_1,f_2,\cdots,f_n\}$ 的过渡矩阵为 $A$，从基 $\{f_1,f_2,\cdots,f_n\}$ 到基 $\{g_1,g_2,\cdots,g_n\}$ 的过渡矩阵为 $B$，那么从 $\{e_1,e_2,\cdots,e_n\}$ 到 $\{g_1,g_2,\cdots,g_n\}$ 的过渡矩阵是什么？

假设
\[
\alpha=\lambda_1e_1+\lambda_2e_2+\cdots+\lambda_ne_n
=\mu_1f_1+\mu_2f_2+\cdots+\mu_nf_n
=\xi_1g_1+\xi_2g_2+\cdots+\xi_ng_n,
\]
则
\[
\left(\begin{array}{c}
\lambda_1\\
\lambda_2\\
\vdots\\
\lambda_n
\end{array}\right)
=A
\left(\begin{array}{c}
\mu_1\\
\mu_2\\
\vdots\\
\mu_n
\end{array}\right),\qquad
\left(\begin{array}{c}
\mu_1\\
\mu_2\\
\vdots\\
\mu_n
\end{array}\right)
=B
\left(\begin{array}{c}
\xi_1\\
\xi_2\\
\vdots\\
\xi_n
\end{array}\right),
\]
于是
\[
\left(\begin{array}{c}
\lambda_1\\
\lambda_2\\
\vdots\\
\lambda_n
\end{array}\right)
=AB
\left(\begin{array}{c}
\xi_1\\
\xi_2\\
\vdots\\
\xi_n
\end{array}\right).
\]
这就是说 $AB$ 是从基 $\{e_1,e_2,\cdots,e_n\}$ 到基 $\{g_1,g_2,\cdots,g_n\}$ 的过渡矩阵.

\noindent\textbf{例 3.8.1}\quad 设在 $\mathbb{K}^3$ 中有两组基 $f_1=(1,0,-1),f_2=(2,1,1),f_3=(1,1,1)$ 和 $g_1=(0,1,1),g_2=(-1,1,0),g_3=(1,2,1)$. 求从 $\{f_1,f_2,f_3\}$ 到 $\{g_1,g_2,g_3\}$ 的过渡矩阵.

\noindent\textbf{解}\quad 这道题如果直接做，将面临解一个九元一次方程组，非常麻烦. 我们利用上面的结论可以很快得到所要求的结果.

设 $e_1=(1,0,0),e_2=(0,1,0),e_3=(0,0,1)$，则从 $\{e_1,e_2,e_3\}$ 到 $\{f_1,f_2,f_3\}$ 的过渡矩阵为
\[
A=\left(\begin{array}{rrr}
1&2&1\\
0&1&1\\
-1&1&1
\end{array}\right),
\]
从 $\{e_1,e_2,e_3\}$ 到 $\{g_1,g_2,g_3\}$ 的过渡矩阵为
\[
B=\left(\begin{array}{rrr}
0&-1&1\\
1&1&2\\
1&0&1
\end{array}\right),
\]
则从 $\{f_1,f_2,f_3\}$ 到 $\{g_1,g_2,g_3\}$ 的过渡矩阵为 $A^{-1}B$. 用第二章例 2.5.4 的方法求矩阵 $A^{-1}B$：
\[
(A\mid B)=
\left(\begin{array}{rrr|rrr}
1&2&1&0&-1&1\\
0&1&1&1&1&2\\
-1&1&1&1&0&1
\end{array}\right)
\to
\left(\begin{array}{rrr|rrr}
1&2&1&0&-1&1\\
0&1&1&1&1&2\\
0&3&2&1&-1&2
\end{array}\right)
\to
\left(\begin{array}{rrr|rrr}
1&0&-1&-2&-3&-3\\
0&1&1&1&1&2\\
0&0&-1&-2&-4&-4
\end{array}\right)
\]
\[
\to
\left(\begin{array}{rrr|rrr}
1&0&0&0&1&1\\
0&1&0&-1&-3&-2\\
0&0&-1&-2&-4&-4
\end{array}\right)
\to
\left(\begin{array}{rrr|rrr}
1&0&0&0&1&1\\
0&1&0&-1&-3&-2\\
0&0&1&2&4&4
\end{array}\right).
\]
因此从基 $\{f_1,f_2,f_3\}$ 到基 $\{g_1,g_2,g_3\}$ 的过渡矩阵为
\[
P=\left(\begin{array}{rrr}
0&1&1\\
-1&-3&-2\\
2&4&4
\end{array}\right).
\]

\noindent\textbf{注}\quad 我们在 (3.8.3) 式中采用的坐标向量是列向量. 如果采用行向量，也可有类似的坐标变换公式，但在形式上略有不同. 我们把结论列出如下，而把证明留给读者. 这时 (3.8.2) 式的矩阵表示应改为
\[
(\lambda_1,\lambda_2,\cdots,\lambda_n)=(\mu_1,\mu_2,\cdots,\mu_n)
\left(\begin{array}{cccc}
a_{11}&a_{12}&\cdots&a_{1n}\\
a_{21}&a_{22}&\cdots&a_{2n}\\
\vdots&\vdots&&\vdots\\
a_{n1}&a_{n2}&\cdots&a_{nn}
\end{array}\right).
\]
我们注意到上式正好是 (3.8.3) 式的转置.

\subsection*{习题 3.8}

1. 求向量 $\alpha$ 在下列基下的坐标，$\{e_1,e_2,e_3,e_4\}$ 是 $\mathbb{K}^4$ 的基：

(1) $\alpha=(1,2,1,3),e_1=(1,0,0,0),e_2=(1,0,1,0),e_3=(0,1,0,-1),e_4=(1,0,1,1)$；

(2) $\alpha=(1,0,0,0),e_1=(1,1,0,1),e_2=(2,1,3,1),e_3=(1,1,0,0),e_4=(0,1,-2,-1)$.

2. 求从基 $\{e_1,e_2,e_3,e_4\}$ 到基 $\{f_1,f_2,f_3,f_4\}$ 的过渡矩阵：

(1)
\[
\left\{\begin{array}{l}
e_1=(1,0,0,0),\\
e_2=(0,1,0,0),\\
e_3=(0,0,1,0),\\
e_4=(0,0,0,1);
\end{array}\right.
\qquad
\left\{\begin{array}{l}
f_1=(1,0,0,1),\\
f_2=(0,0,1,-1),\\
f_3=(2,1,0,3),\\
f_4=(-1,0,1,2);
\end{array}\right.
\]

(2)
\[
\left\{\begin{array}{l}
e_1=(1,1,0,1),\\
e_2=(2,1,2,0),\\
e_3=(1,1,0,0),\\
e_4=(0,1,-1,-1);
\end{array}\right.
\qquad
\left\{\begin{array}{l}
f_1=(1,0,0,1),\\
f_2=(0,0,1,-1),\\
f_3=(2,1,0,3),\\
f_4=(-1,0,1,2).
\end{array}\right.
\]

3. 求向量 $\alpha$ 在第 2 题中两组基下的坐标：

(1) $\alpha=(1,0,0,1)$；\qquad (2) $\alpha=(3,-1,0,2)$.

4. 设 $a$ 为常数，求向量 $\alpha=(a_1,a_2,\cdots,a_n)$ 在基
\[
\{f_1=(a^{n-1},a^{n-2},\cdots,a,1),f_2=(a^{n-2},a^{n-3},\cdots,1,0),\cdots,f_n=(1,0,\cdots,0,0)\}
\]
下的坐标.

5. 设 $V$ 是由次数不超过 $n$ 的实多项式全体组成的线性空间，求从基 $\{1,x,x^2,\cdots,x^n\}$ 到基 $\{1,x-a,(x-a)^2,\cdots,(x-a)^n\}$ 的过渡矩阵，并以此证明多项式的 Taylor（泰勒）公式：
\[
f(x)=f(a)+\frac{f'(a)}{1!}(x-a)+\frac{f''(a)}{2!}(x-a)^2+\cdots+\frac{f^{(n)}(a)}{n!}(x-a)^n,
\]
其中 $f^{(n)}(x)$ 表示 $f(x)$ 的 $n$ 次导数.

\section*{§3.9\quad 子空间}

在空间解析几何中，读者已经知道，通常的三维空间中任一经过原点的平面有这样的性质：任意以原点为始点、终点在该平面内的向量之和仍在该平面内；平面内任一向量的数乘也仍在该平面内，即这个平面关于向量的线性组合封闭. 这样的平面称为三维空间的子空间. 我们把这一概念推广到一般的线性空间.

\noindent\textbf{定义 3.9.1}\quad 设 $V$ 是数域 $\mathbb{K}$ 上的线性空间，$V_0$ 是 $V$ 的非空子集，且对 $V_0$ 中的任意两个向量 $\alpha,\beta$ 及 $\mathbb{K}$ 中任一数 $k$，总有 $\alpha+\beta\in V_0$ 及 $k\alpha\in V_0$，则称 $V_0$ 是 $V$ 的线性子空间，简称子空间.

\noindent\textbf{命题 3.9.1}\quad 定义 3.9.1 中的 $V_0$ 在 $V$ 的加法及数乘下是数域 $\mathbb{K}$ 上的线性空间.

\noindent\textbf{证明}\quad 我们要证明 $V_0$ 适合定义 3.3.1 中的 8 条规则. 因为 $V_0$ 是 $V$ 的子集，因此 $V_0$ 中向量的加法适合定义规则 (1)，(2). 由于 $V_0$ 非空，$0=\alpha+(-1)\alpha$，因此规则 (3) 成立. 又 $-\alpha=(-1)\alpha$，因此规则 (4) 也成立. 规则 (5)～(8) 显然对 $V_0$ 成立，因此 $V_0$ 是 $\mathbb{K}$ 上的线性空间. $\square$

由上述命题，我们称定义 3.9.1 中的 $V_0$ 为线性子空间是合理的. 由定义不难证明，对 $V_0$ 中任意有限个向量 $\alpha_1,\alpha_2,\cdots,\alpha_m$，它们任意的线性组合
\[
\lambda_1\alpha_1+\lambda_2\alpha_2+\cdots+\lambda_m\alpha_m
\]
仍属于 $V_0$.

任一线性空间 $V$ 至少有两个子空间，一是由零向量 $\{0\}$ 组成的子空间，称为零子空间（维数规定为 0）；另一个是 $V$ 自身，这两个子空间通常称为平凡子空间.

如果 $V$ 是 $n$ 维线性空间，则由推论 3.5.1 知道，$V$ 的任一子空间的维数不超过 $n$. 若 $V_0$ 是 $V$ 的非平凡子空间，则
\[
0<\dim V_0<\dim V=n.
\]
事实上，若 $\dim V_0=n$，则 $V_0$ 有一组由 $n$ 个向量组成的基. 但任意 $n$ 个线性无关的向量均可组成 $V$ 的一组基，因此这时 $V_0=V$，与 $V_0$ 是非平凡子空间矛盾.

在通常的三维空间中，通过原点的平面是二维子空间，过原点的直线是一维子空间. 反之，二维子空间必是过原点的平面，一维子空间必是过原点的直线.

\noindent\textbf{定义 3.9.2}\quad 若 $V_1,V_2$ 是 $V$ 的子空间，定义它们的交为既在 $V_1$ 中又在 $V_2$ 中的全体向量组成的集合 $V_1\cap V_2$. 定义它们的和为
\[
V_1+V_2=\{\alpha+\beta\mid \alpha\in V_1,\beta\in V_2\},
\]
即所有形如 $\alpha+\beta$ 的向量的集合，其中要求 $\alpha\in V_1,\beta\in V_2$.

\noindent\textbf{命题 3.9.2}\quad $V_1\cap V_2,V_1+V_2$ 都是 $V$ 的子空间.

\noindent\textbf{证明}\quad 按定义 3.9.1 直接验证即可. $\square$

类似地，可以定义 $m$ 个子空间的交
\[
V_1\cap V_2\cap\cdots\cap V_m
\]
为属于所有 $V_i(i=1,2,\cdots,m)$ 的向量全体组成的子集，它也是 $V$ 的子空间. 定义 $m$ 个子空间的和为
\[
V_1+V_2+\cdots+V_m=\{\alpha_1+\alpha_2+\cdots+\alpha_m\mid \alpha_i\in V_i,\ i=1,2,\cdots,m\},
\]
$V_1+V_2+\cdots+V_m$ 也是 $V$ 的子空间.

\noindent\textbf{例 3.9.1}\quad 在三维行向量空间 $\mathbb{K}^3$ 中，定义 $V_1=\{(a,0,0)\mid a\in\mathbb{K}\}$，$V_2=\{(0,b,0)\mid b\in\mathbb{K}\}$，$V_3=\{(0,0,c)\mid c\in\mathbb{K}\}$，$V_4=\{(a,b,0)\mid a,b\in\mathbb{K}\}$，$V_5=\{(0,b,c)\mid b,c\in\mathbb{K}\}$，则 $V_4=V_1+V_2,V_5=V_2+V_3,V_2=V_4\cap V_5,\mathbb{K}^3=V_3+V_4=V_1+V_5=V_4+V_5=V_1+V_2+V_3$.

\noindent\textbf{定义 3.9.3}\quad 设 $S$ 是线性空间 $V$ 的子集，记 $L(S)$ 为 $S$ 中向量所有可能的线性组合构成的子集，则由定义 3.9.1 不难看出，$L(S)$ 是 $V$ 的一个子空间，称为由 $S$ 生成的子空间，或由 $S$ 张成的子空间.

\noindent\textbf{定理 3.9.1}\quad 设 $S$ 是线性空间 $V$ 的子集，$L(S)$ 为由 $S$ 张成的子空间.

(1) $S\subseteq L(S)$ 且若 $V_0$ 是包含集合 $S$ 的子空间，则 $L(S)\subseteq V_0$，也即 $L(S)$ 是包含 $S$ 的 $V$ 的最小子空间；

(2) $L(S)$ 的维数等于 $S$ 中极大无关组所含向量的个数，且若 $\alpha_1,\alpha_2,\cdots,\alpha_m$ 是 $S$ 的极大无关组，则
\[
L(S)=L(\alpha_1,\alpha_2,\cdots,\alpha_m).
\]

\noindent\textbf{证明}\quad (1) 显然 $S\subseteq L(S)$. 设 $\beta\in L(S)$，则 $\beta$ 是 $S$ 中若干个向量 $\alpha_1,\alpha_2,\cdots,\alpha_r$ 的线性组合：
\[
\beta=\lambda_1\alpha_1+\lambda_2\alpha_2+\cdots+\lambda_r\alpha_r.
\]
因为 $S\subseteq V_0$，由子空间的定义可知 $\beta\in V_0$，所以 $L(S)\subseteq V_0$.

(2) 设 $\alpha_1,\alpha_2,\cdots,\alpha_m$ 是 $S$ 的极大无关组，则 $S$ 中任一向量都是 $\alpha_1,\alpha_2,\cdots,\alpha_m$ 的线性组合，即
\[
S\subseteq L(\alpha_1,\alpha_2,\cdots,\alpha_m).
\]
因此
\[
L(S)\subseteq L(\alpha_1,\alpha_2,\cdots,\alpha_m).
\]
另一方面，显然有
\[
L(\alpha_1,\alpha_2,\cdots,\alpha_m)\subseteq L(S),
\]
因此
\[
L(S)=L(\alpha_1,\alpha_2,\cdots,\alpha_m).
\]
$L(S)$ 的维数等于 $m$. $\square$

\noindent\textbf{例 3.9.2}\quad 在三维实空间中，由一个非零向量生成的子空间为这个向量所在的直线. 由不在同一条直线上的两个向量生成的子空间为这两个向量所在的平面，由 3 个不在同一平面内的向量生成的子空间即为整个空间.

\noindent\textbf{例 3.9.3}\quad 设 $V_1,V_2$ 是线性空间 $V$ 的子空间，则 $L(V_1\cup V_2)=V_1+V_2$.

\noindent\textbf{证明}\quad 由生成的定义，对任意的 $\alpha_1\in V_1,\alpha_2\in V_2,\alpha_1+\alpha_2\in L(V_1\cup V_2)$，故 $V_1+V_2\subseteq L(V_1\cup V_2)$. 另一方面，因为 $V_1\subseteq V_1+V_2,V_2\subseteq V_1+V_2$，由定理 3.9.1 (1)，$L(V_1\cup V_2)\subseteq V_1+V_2$. 于是 $L(V_1\cup V_2)=V_1+V_2$. $\square$

\noindent\textbf{注}\quad 一般地，不难证明：若 $V_1,V_2,\cdots,V_m$ 是 $V$ 的子空间，则
\[
L(V_1\cup V_2\cup\cdots\cup V_m)=V_1+V_2+\cdots+V_m.
\]

下面我们给出和空间与交空间之间的维数公式.

\noindent\textbf{定理 3.9.2}\quad 设 $V_1,V_2$ 是线性空间 $V$ 的子空间，则
\begin{equation}
\dim(V_1+V_2)=\dim V_1+\dim V_2-\dim(V_1\cap V_2).
\tag{3.9.1}
\end{equation}

\noindent\textbf{证明}\quad 设 $\dim V_1=n_1,\dim V_2=n_2,\dim(V_1\cap V_2)=m$. 取 $V_1\cap V_2$ 的一组基 $\{\alpha_1,\cdots,\alpha_m\}$，由于 $V_1\cap V_2$ 是 $V_1$ 的子空间，故可添上 $V_1$ 中的向量 $\alpha_{m+1},\cdots,\alpha_{n_1}$，使 $\{\alpha_1,\cdots,\alpha_m,\alpha_{m+1},\cdots,\alpha_{n_1}\}$ 是 $V_1$ 的一组基. 同样道理，可添上 $\beta_{m+1},\cdots,\beta_{n_2}$，使 $\{\alpha_1,\cdots,\alpha_m,\beta_{m+1},\cdots,\beta_{n_2}\}$ 成为 $V_2$ 的一组基. 显然，$V_1+V_2$ 中的向量均可由向量组
\begin{equation}
\alpha_1,\cdots,\alpha_m,\alpha_{m+1},\cdots,\alpha_{n_1},\beta_{m+1},\cdots,\beta_{n_2}
\tag{3.9.2}
\end{equation}
的线性组合给出. 如能证明上式中的向量线性无关，则它们构成 $V_1+V_2$ 的一组基，由此即可推出所要的结论. 现假设
\[
\lambda_1\alpha_1+\cdots+\lambda_m\alpha_m+\lambda_{m+1}\alpha_{m+1}+\cdots+\lambda_{n_1}\alpha_{n_1}+\mu_{m+1}\beta_{m+1}+\cdots+\mu_{n_2}\beta_{n_2}=0,
\]
则
\[
\lambda_1\alpha_1+\cdots+\lambda_m\alpha_m+\lambda_{m+1}\alpha_{m+1}+\cdots+\lambda_{n_1}\alpha_{n_1}=-(\mu_{m+1}\beta_{m+1}+\cdots+\mu_{n_2}\beta_{n_2}).
\]
上式左端属于 $V_1$，右端属于 $V_2$，故
\[
\mu_{m+1}\beta_{m+1}+\cdots+\mu_{n_2}\beta_{n_2}\in V_1\cap V_2,
\]
即存在 $\xi_1,\cdots,\xi_m\in\mathbb{K}$，使
\[
\mu_{m+1}\beta_{m+1}+\cdots+\mu_{n_2}\beta_{n_2}=\xi_1\alpha_1+\cdots+\xi_m\alpha_m.
\]
但 $\alpha_1,\cdots,\alpha_m,\beta_{m+1},\cdots,\beta_{n_2}$ 是 $V_2$ 的基，因此 $\mu_{m+1}=\cdots=\mu_{n_2}=\xi_1=\cdots=\xi_m=0$. 再由 $\alpha_1,\cdots,\alpha_m,\alpha_{m+1},\cdots,\alpha_{n_1}$ 线性无关得 $\lambda_1=\cdots=\lambda_m=\lambda_{m+1}=\cdots=\lambda_{n_1}=0$. $\square$

若 $V_1\cap V_2=0$，则由维数公式马上得到 $\dim(V_1+V_2)=\dim V_1+\dim V_2$，即和空间的维数等于维数的和. 我们可以自然地问：若 $V_1,V_2,\cdots,V_m$ 是 $V$ 的子空间，当它们之间满足怎样的关系时，$\dim(V_1+V_2+\cdots+V_m)=\dim V_1+\dim V_2+\cdots+\dim V_m$ 才会成立呢？为了回答这个问题，我们需要先引入直和的概念，它是子空间的一类特别重要的和.

\noindent\textbf{定义 3.9.4}\quad 设 $V_1,V_2,\cdots,V_m$ 是线性空间 $V$ 的子空间，若对一切 $i(i=1,2,\cdots,m)$，
\[
V_i\cap(V_1+\cdots+V_{i-1}+V_{i+1}+\cdots+V_m)=0,
\]
则称和 $V_1+V_2+\cdots+V_m$ 为直接和，简称直和，记为
\[
V_1\oplus V_2\oplus\cdots\oplus V_m.
\]

\noindent\textbf{例 3.9.4}\quad 同例 3.9.1 的假设，则 $V_4=V_1\oplus V_2,V_5=V_2\oplus V_3,\mathbb{K}^3=V_3\oplus V_4=V_1\oplus V_5=V_1\oplus V_2\oplus V_3$ 都是直和，但 $\mathbb{K}^3=V_4+V_5$ 不是直和.

\noindent\textbf{定理 3.9.3}\quad 设 $V_1,V_2,\cdots,V_m$ 是线性空间 $V$ 的子空间，$V_0=V_1+V_2+\cdots+V_m$，则下列命题等价：

(1) $V_0=V_1\oplus V_2\oplus\cdots\oplus V_m$ 是直和；

(2) 对任意的 $2\leq i\leq m$，
\[
V_i\cap(V_1+V_2+\cdots+V_{i-1})=0;
\]

(3) $\dim(V_1+V_2+\cdots+V_m)=\dim V_1+\dim V_2+\cdots+\dim V_m$；

(4) $V_1,V_2,\cdots,V_m$ 的一组基可以拼成 $V_0$ 的一组基；

(5) $V_0$ 中的向量表示为 $V_1,V_2,\cdots,V_m$ 中的向量之和时其表示唯一，即若 $\alpha\in V_0$ 且
\[
\alpha=v_1+v_2+\cdots+v_m=u_1+u_2+\cdots+u_m,
\]
其中 $v_i,u_i\in V_i$，则 $u_i=v_i(i=1,2,\cdots,m)$.

\noindent\textbf{证明}\quad (1) $\Rightarrow$ (2)：显然.

(2) $\Rightarrow$ (3)：由维数公式可知，对任意的 $2\leq i\leq m$，
\[
\dim(V_1+V_2+\cdots+V_i)=\dim(V_1+V_2+\cdots+V_{i-1})+\dim V_i.
\]
不断迭代下去即得
\[
\dim(V_1+V_2+\cdots+V_m)=\dim V_1+\dim V_2+\cdots+\dim V_m.
\]

(3) $\Rightarrow$ (4)：依次取 $V_1,V_2,\cdots,V_m$ 的一组基为
\begin{equation}
e_{11},\cdots,e_{1n_1};e_{21},\cdots,e_{2n_2};\cdots;e_{m1},\cdots,e_{mn_m},
\tag{3.9.3}
\end{equation}
则 $\dim V_0=n_1+n_2+\cdots+n_m$，且 $V_0$ 中任一向量均可由 (3.9.3) 式中向量的线性组合来表示. 由定理 3.5.3 知 (3.9.3) 式中的向量构成了 $V_0$ 的一组基，即 $V_1,V_2,\cdots,V_m$ 的一组基可以拼成 $V_0$ 的一组基.

接下去我们先证明 (5) 等价于一个弱一点的条件：

$(5')$ $V_0$ 中的零向量表示为 $V_1,V_2,\cdots,V_m$ 中的向量之和时其表示唯一，即若
\[
0=v_1+v_2+\cdots+v_m,
\]
其中 $v_i\in V_i$，则 $v_i=0(i=1,2,\cdots,m)$.

(5) $\Rightarrow$ $(5')$：显然. 反之，由向量 $\alpha$ 的两个表示可得
\[
0=(u_1-v_1)+(u_2-v_2)+\cdots+(u_m-v_m),
\]
从而由 $(5')$ 可推出 (5) 的结论.

(4) $\Rightarrow$ $(5')$：依次取 $V_1,V_2,\cdots,V_m$ 的一组基如 (3.9.3) 式所设，于是它们拼成了 $V_0$ 的一组基，特别地，(3.9.3) 式中诸向量线性无关. 对任意的 $1\leq i\leq m$，设
\[
v_i=\lambda_{i1}e_{i1}+\lambda_{i2}e_{i2}+\cdots+\lambda_{in_i}e_{in_i},
\]
则
\[
\begin{aligned}
0={}&\lambda_{11}e_{11}+\lambda_{12}e_{12}+\cdots+\lambda_{1n_1}e_{1n_1}\\
&+\lambda_{21}e_{21}+\lambda_{22}e_{22}+\cdots+\lambda_{2n_2}e_{2n_2}+\cdots\\
&+\lambda_{m1}e_{m1}+\lambda_{m2}e_{m2}+\cdots+\lambda_{mn_m}e_{mn_m}.
\end{aligned}
\]
由 (3.9.3) 式中向量的线性无关性，即得 $\lambda_{ij}=0$. 因此 $v_i=0(i=1,2,\cdots,m)$.

$(5')\Rightarrow$ (1)：任取 $v\in V_i\cap(V_1+\cdots+V_{i-1}+V_{i+1}+\cdots+V_m)$，则
\[
v=v_1+\cdots+v_{i-1}+v_{i+1}+\cdots+v_m,
\]
其中 $v_j\in V_j(j=1,\cdots,i-1,i+1,\cdots,m)$，于是
\[
0=v_1+\cdots+v_{i-1}+(-v)+v_{i+1}+\cdots+v_m.
\]
注意到 $v\in V_i$，由 $(5')$ 可得 $v=0$，从而对任意的 $1\leq i\leq m$，
\[
V_i\cap(V_1+\cdots+V_{i-1}+V_{i+1}+\cdots+V_m)=0,
\]
即 $V_1+V_2+\cdots+V_m$ 是直和. $\square$

\subsection*{习题 3.9}

1. 判断 $n$ 维实行向量空间 $\mathbb{R}^n$ 的下列子集是否子空间：

(1) $S=\{(a_1,a_2,\cdots,a_n)\mid \sum_{i=1}^n a_i=0\}$；\qquad
(2) $S=\{(a_1,a_2,\cdots,a_n)\mid \sum_{i=1}^n a_i^2=1\}$；

(3) $S=\{(a_1,a_2,\cdots,a_n)\mid a_1=0\}$；\qquad
(4) $S=\{(a_1,a_2,\cdots,a_n)\mid a_i\geq 0\}$；

(5) $S=\{(a_1,a_2,\cdots,a_n)\mid a_1=a_2=\cdots=a_n\}$.

2. 试求 $\mathbb{R}^3$ 中下列子空间的维数：

(1) $V=\{(a,0,b)\mid a,b\in\mathbb{R}\}$；\qquad
(2) $V=\{(a,2a,b)\mid a,b\in\mathbb{R}\}$；

(3) $V=L((1,0,1),(1,2,3))$；\qquad
(4) $V=L((1,2,-1),(3,6,-3))$.

3. 若 $V_1,V_2$ 是线性空间 $V$ 的子空间且 $V_1\subseteq V_2$. 求证：$V_1=V_2$ 成立的充分必要条件是 $\dim V_1=\dim V_2$.

4. 设 $V=M_n(\mathbb{K})$ 是数域 $\mathbb{K}$ 上的 $n$ 阶矩阵全体组成的线性空间，$A\in V$，求证：与 $A$ 乘法可交换的矩阵全体 $C(A)$ 组成 $V$ 的子空间且其维数不为零. 又若 $T$ 是 $V$ 的非空子集，求证：与 $T$ 中任一矩阵乘法可交换的矩阵全体 $C(T)$ 也构成 $V$ 的子空间且其维数不为零.

5. 求 §2.2 习题 13 中子空间 $C(A)$ 的一组基和维数.

6. 设 $V_1,V_2,V_3$ 是线性空间 $V$ 的子空间.

(1) 举例说明：$V_1\cap(V_2+V_3)=V_1\cap V_2+V_1\cap V_3$ 未必成立；

(2) 若 $V_1$ 包含 $V_2$ 或 $V_1$ 包含 $V_3$，求证：$V_1\cap(V_2+V_3)=V_1\cap V_2+V_1\cap V_3$ 成立. 进一步，若 $V_2+V_3$ 是直和，求证：$V_1\cap(V_2\oplus V_3)=(V_1\cap V_2)\oplus(V_1\cap V_3)$ 成立.

7. 设 $\alpha_1=(1,0,-1,0),\alpha_2=(0,1,2,1),\alpha_3=(2,1,0,1)$ 是四维实行向量空间 $V$ 中的向量，它们生成的子空间为 $V_1$，又向量 $\beta_1=(-1,1,1,1),\beta_2=(1,-1,-3,-1),\beta_3=(-1,1,-1,1)$ 生成的子空间为 $V_2$，求子空间 $V_1+V_2$ 和 $V_1\cap V_2$ 的基.

8. 设 $\alpha_1,\alpha_2,\alpha_3$ 是线性空间 $V$ 中的向量，若它们两两线性无关但全体线性相关，求证：
\[
L(\alpha_1,\alpha_2)=L(\alpha_1,\alpha_3)=L(\alpha_2,\alpha_3).
\]

9. 若 $V$ 的子空间 $V_1,V_2,V_3$ 满足 $V_i\cap V_j=0(1\leq i<j\leq 3)$，问：$V_1+V_2+V_3$ 是否直和？

10. 设 $V$ 是数域 $\mathbb{K}$ 上 $n$ 阶矩阵组成的向量空间，$V_1$ 和 $V_2$ 分别是 $\mathbb{K}$ 上对称阵和反对称阵组成的子集，求证：$V_1$ 和 $V_2$ 都是 $V$ 的子空间且 $V=V_1\oplus V_2$.

11. 设 $U,V$ 是数域 $\mathbb{K}$ 上的两个线性空间，$W=U\times V$ 是 $U$ 和 $V$ 的积集合，即 $W=\{(u,v)\mid u\in U,v\in V\}$. 现在 $W$ 上定义加法和数乘：
\[
(u_1,v_1)+(u_2,v_2)=(u_1+u_2,v_1+v_2),\quad k(u,v)=(ku,kv).
\]
验证：$W$ 是 $\mathbb{K}$ 上的线性空间 (这个线性空间称为 $U$ 和 $V$ 的外直和). 又若设 $U'=\{(u,0)\mid u\in U\},V'=\{(0,v)\mid v\in V\}$，求证：$U',V'$ 是 $W$ 的子空间，$U'$ 和 $U$ 同构，$V'$ 和 $V$ 同构，并且 $W=U'\oplus V'$.

12. 设 $U$ 是 $V$ 的子空间，求证：存在 $V$ 的子空间 $W$，使得 $V=U\oplus W$. 这样的子空间 $W$ 称为子空间 $U$ 在 $V$ 中的补空间.

13. 若 $V=U\oplus W$，且 $U=U_1\oplus U_2$，求证：$V=U_1\oplus U_2\oplus W$.

14. 求证：每一个 $n$ 维线性空间均可表示为 $n$ 个一维子空间的直和.

15. 设 $V_1,V_2,\cdots,V_m$ 是数域 $\mathbb{K}$ 上向量空间 $V$ 的 $m$ 个非平凡子空间，证明：在 $V$ 中必存在一个向量 $\alpha$，它不属于任何一个 $V_i$.

\section*{§3.10\quad 线性方程组的解}

用矩阵秩的概念我们很容易给出一般线性方程组解的判定定理.

\noindent\textbf{定理 3.10.1}\quad 设有 $n$ 个未知数 $m$ 个方程式组成的线性方程组：
\begin{equation}
\left\{\begin{array}{c}
a_{11}x_1+a_{12}x_2+\cdots+a_{1n}x_n=b_1,\\
a_{21}x_1+a_{22}x_2+\cdots+a_{2n}x_n=b_2,\\
\cdots\cdots\cdots\cdots\\
a_{m1}x_1+a_{m2}x_2+\cdots+a_{mn}x_n=b_m,
\end{array}\right.
\tag{3.10.1}
\end{equation}
它的系数矩阵记为 $A$，增广矩阵记为 $\widetilde{A}$，即
\[
\widetilde{A}=\left(\begin{array}{ccccc}
a_{11}&a_{12}&\cdots&a_{1n}&b_1\\
a_{21}&a_{22}&\cdots&a_{2n}&b_2\\
\vdots&\vdots&&\vdots&\vdots\\
a_{m1}&a_{m2}&\cdots&a_{mn}&b_m
\end{array}\right),
\]
则有下列结论：

(1) 若 $\widetilde{A}$ 与 $A$ 的秩都等于 $n$，则该方程组有唯一一组解；

(2) 若 $\widetilde{A}$ 与 $A$ 的秩相等但小于 $n$，即 $\mathrm{r}(\widetilde{A})=\mathrm{r}(A)<n$，则该方程组有无穷多组解；

(3) 若 $\widetilde{A}$ 与 $A$ 的秩不相等，则该方程组无解.

\noindent\textbf{证明}\quad (1) 首先我们证明方程组 (3.10.1) 有解的充分必要条件是
\[
\mathrm{r}(\widetilde{A})=\mathrm{r}(A).
\]
如同第 4 节那样把方程组 (3.10.1) 写成向量形式就是
\begin{equation}
x_1\alpha_1+x_2\alpha_2+\cdots+x_n\alpha_n=\beta,
\tag{3.10.2}
\end{equation}
其中 $\alpha_i$ 为矩阵 $A$ 的第 $i$ 个列向量，$\beta$ 为常数项向量. 方程组 (3.10.1) 有解等同于 $\beta$ 是 $\alpha_1,\alpha_2,\cdots,\alpha_n$ 的线性组合. 因此 $\mathrm{r}(\widetilde{A})=\mathrm{r}(A)$.

反过来，若 $\mathrm{r}(\widetilde{A})=\mathrm{r}(A)$，则 $A$ 的列向量的极大线性无关组就是 $\widetilde{A}$ 的列向量的极大线性无关组. 因此 $\beta$ 可表示为 $A$ 的列向量的线性组合，即方程组 (3.10.2) 有解.

若再有 $\mathrm{r}(\widetilde{A})=\mathrm{r}(A)=n$，此时 $A$ 的 $n$ 个列向量线性无关，所以 $\beta$ 只有唯一一种方法表示为 $\alpha_1,\alpha_2,\cdots,\alpha_n$ 的线性组合，即方程组只有唯一一组解.

(2) 若 $\mathrm{r}(\widetilde{A})=\mathrm{r}(A)=r<n$，则 $\alpha_1,\alpha_2,\cdots,\alpha_n$ 线性相关，即存在不全为零的数 $k_1,k_2,\cdots,k_n$，使得
\[
k_1\alpha_1+k_2\alpha_2+\cdots+k_n\alpha_n=0.
\]
这时对任意的数 $k$，
\[
kk_1\alpha_1+kk_2\alpha_2+\cdots+kk_n\alpha_n=0.
\]
因此若 $x_1=c_1,x_2=c_2,\cdots,x_n=c_n$ 是方程组的解，则 $x_1=kk_1+c_1,x_2=kk_2+c_2,\cdots,x_n=kk_n+c_n$ 都是解，显然这样的解有无穷多组. $\square$

当线性方程组 (3.10.1) 有无穷多组解时，我们能否用有限组解来把握所有的解，这是我们要解决的问题. 我们先研究最简单的情形，即所谓齐次线性方程组的解.

\noindent\textbf{定义 3.10.1}\quad 若线性方程组 (3.10.1) 的所有常数项 $b_i$ 都为零，则称之为齐次线性方程组，否则称之为非齐次线性方程组.

齐次线性方程组的矩阵形式为
\begin{equation}
Ax=0.
\tag{3.10.3}
\end{equation}
注意到增广矩阵 $\widetilde{A}=(A\mid 0)$ 的秩与 $A$ 的秩显然相等，且 $x_1=x_2=\cdots=x_n=0$ 总是方程组 (3.10.3) 的解. 若 $\mathrm{r}(A)<n$，则方程组 (3.10.3) 有无穷多组解；当 $\mathrm{r}(A)=n$ 时只有零解，这时，我们称方程组 (3.10.3) 只有平凡解. 我们的目的是在方程组有非平凡解时找出所有的解.

从向量空间的观点来看方程组 (3.10.3)，它的一个解可以看成是 $n$ 维列向量空间中的一个元素. 若 $\alpha,\beta$ 是方程组 (3.10.3) 的解，即 $A\alpha=0,A\beta=0$，则 $A(\alpha+\beta)=0$. 对任意的数 $k,A(k\alpha)=kA\alpha=0$. 因此，$\alpha+\beta$ 及 $k\alpha$ 均是方程组 (3.10.3) 的解. 这一事实表明方程组 (3.10.3) 的全部解构成 $n$ 维列向量空间的一个子空间. 这个子空间称为齐次线性方程组 (3.10.3) 的解空间. 这个解空间的维数是多少呢？我们如何来求出该子空间的一组基 (这时可将方程组的任意一组解表示为基向量的线性组合，从而达到用有限多组解表示无穷多组解的目的)？

设 $\mathrm{r}(A)=r<n$，则 $A$ 有 $r$ 个行向量线性无关，因此剔除了多余的方程式以后，还剩 $r$ 个方程式，不妨就设为前 $r$ 个方程式. 如果允许未知数的对换，则我们可假设这 $r$ 个方程式系数矩阵的前 $r$ 个列向量线性无关，将 (3.10.3) 式化为
\begin{equation}
\left\{\begin{array}{c}
a_{11}x_1+\cdots+a_{1r}x_r=-a_{1,r+1}x_{r+1}-\cdots-a_{1n}x_n,\\
\cdots\cdots\cdots\\
a_{r1}x_1+\cdots+a_{rr}x_r=-a_{r,r+1}x_{r+1}-\cdots-a_{rn}x_n.
\end{array}\right.
\tag{3.10.4}
\end{equation}
将上述方程组看成是 $r$ 个未知数的线性方程组，注意到它的系数行列式不等于零. 用 Cramer 法则将方程组 (3.10.4) 解出来，其解含有 $n-r$ 个参数，不妨设为
\begin{equation}
\left\{\begin{array}{c}
x_1=c_{1,r+1}x_{r+1}+\cdots+c_{1n}x_n,\\
\cdots\cdots\cdots\\
x_r=c_{r,r+1}x_{r+1}+\cdots+c_{rn}x_n,
\end{array}\right.
\tag{3.10.5}
\end{equation}
其中 $x_{r+1},\cdots,x_n$ 可取任何数，依次取
\[
\left\{\begin{array}{c}
x_{r+1}=1,x_{r+2}=0,\cdots,x_n=0,\\
x_{r+1}=0,x_{r+2}=1,\cdots,x_n=0,\\
\cdots\cdots\cdots\\
x_{r+1}=0,x_{r+2}=0,\cdots,x_n=1,
\end{array}\right.
\]
便得到 $n-r$ 个解：
\[
\eta_1=\left(\begin{array}{c}
c_{1,r+1}\\
\vdots\\
c_{r,r+1}\\
1\\
0\\
\vdots\\
0
\end{array}\right),\quad
\eta_2=\left(\begin{array}{c}
c_{1,r+2}\\
\vdots\\
c_{r,r+2}\\
0\\
1\\
\vdots\\
0
\end{array}\right),\quad \cdots,\quad
\eta_{n-r}=\left(\begin{array}{c}
c_{1n}\\
\vdots\\
c_{rn}\\
0\\
0\\
\vdots\\
1
\end{array}\right).
\]
不难看出 $\eta_1,\eta_2,\cdots,\eta_{n-r}$ 线性无关.

我们希望这 $n-r$ 个解向量就是解空间的基. 为此，我们只需证明方程组的任意一个解可以表示为这 $n-r$ 个解的线性组合即可. 设 $\eta$ 是齐次线性方程组 (3.10.3) 的任一解向量，且设
\[
\eta=\left(\begin{array}{c}
a_1\\
a_2\\
\vdots\\
a_n
\end{array}\right),
\]
则 $a_1,a_2,\cdots,a_n$ 必须适合 (3.10.5) 式，即
\[
\left\{\begin{array}{c}
a_1=c_{1,r+1}a_{r+1}+\cdots+c_{1n}a_n,\\
\cdots\cdots\cdots\\
a_r=c_{r,r+1}a_{r+1}+\cdots+c_{rn}a_n.
\end{array}\right.
\]
于是
\[
\eta=a_{r+1}\eta_1+a_{r+2}\eta_2+\cdots+a_n\eta_{n-r}.
\]
这即表明方程组 (3.10.3) 的任一解均可表示为 $\eta_1,\eta_2,\cdots,\eta_{n-r}$ 的线性组合，因此 $\eta_1,\eta_2,\cdots,\eta_{n-r}$ 是方程组 (3.10.3) 解空间的一组基. 由此即知方程组 (3.10.3) 的解空间维数为 $n-r$. 一个齐次线性方程组解空间的基又称为该方程组的基础解系. 我们把上面的论证总结为如下定理.

\noindent\textbf{定理 3.10.2}\quad 设有齐次线性方程组
\begin{equation}
Ax=0,
\tag{3.10.6}
\end{equation}
其中 $A=(a_{ij})$ 是 $m\times n$ 矩阵. 若 $\mathrm{r}(A)=r<n$，则上述方程组有非零解. 它的解构成 $n$ 维列向量空间的一个 $n-r$ 维子空间. 也就是说，存在由 $n-r$ 个向量构成的基础解系 $\{\eta_1,\eta_2,\cdots,\eta_{n-r}\}$，使方程组 (3.10.6) 的任一组解均可表示为 $\{\eta_1,\eta_2,\cdots,\eta_{n-r}\}$ 的线性组合.

现在我们转而考虑非齐次线性方程组，它的矩阵形式为
\begin{equation}
Ax=\beta.
\tag{3.10.7}
\end{equation}
称齐次线性方程组
\begin{equation}
Ax=0
\tag{3.10.8}
\end{equation}
为方程组 (3.10.7) 的相伴齐次线性方程组 (或导出组). 若 $\gamma$ 是方程组 (3.10.7) 的一个解，且 $\alpha$ 是其相伴齐次线性方程组的解，则
\[
A(\alpha+\gamma)=A\alpha+A\gamma=0+\beta=\beta.
\]
因此相伴齐次线性方程组的任一解与方程组 (3.10.7) 的解之和仍是方程组 (3.10.7) 的解. 现固定方程组 (3.10.7) 的一个解 $\gamma$，称之为方程组 (3.10.7) 的一个特解，我们要求出方程组 (3.10.7) 的一切解. 设 $\xi$ 是方程组 (3.10.7) 的任一解，则
\[
A\xi=\beta=A\gamma,
\]
于是
\[
A(\xi-\gamma)=0,
\]
即 $\xi-\gamma$ 是相伴齐次线性方程组的解. 因此只要知道方程组 (3.10.7) 的一个特解和相伴齐次线性方程组的一切解，那么它的任一解都可以知道了. 由定理 3.10.2，相伴齐次线性方程组 (3.10.8) 有一组基础解系 $\{\eta_1,\eta_2,\cdots,\eta_{n-r}\}$，因此方程组 (3.10.7) 的解可表示为
\[
k_1\eta_1+k_2\eta_2+\cdots+k_{n-r}\eta_{n-r}+\gamma,
\]
于是我们得到了非齐次线性方程组解的结构定理.

\noindent\textbf{定理 3.10.3}\quad 设有非齐次线性方程组 (3.10.7)，它的系数矩阵 $A$ 及其增广矩阵 $\widetilde{A}$ 的秩都等于 $r(r<n)$. 假设方程组 (3.10.7) 相伴的齐次线性方程组有基础解系 $\{\eta_1,\eta_2,\cdots,\eta_{n-r}\}$，又 $\gamma$ 是方程组 (3.10.7) 的任一特解，则其所有解均可表示为如下形状：
\begin{equation}
k_1\eta_1+k_2\eta_2+\cdots+k_{n-r}\eta_{n-r}+\gamma,
\tag{3.10.9}
\end{equation}
其中 $k_1,k_2,\cdots,k_{n-r}$ 可取任何数.

\noindent\textbf{注}\quad 我们在讨论线性方程组的解时并没有特别强调在哪个数域中求解，事实上，如果我们限定在数域 $\mathbb{K}$ 中求解方程组 (3.10.7) (这时 $A$ 中的元素当然必须属于 $\mathbb{K}$)，那么 (3.10.9) 式中的 $k_i$ 就必须取自 $\mathbb{K}$ 中.

我们已经从理论上完全解决了求解线性方程组的问题，下面举例说明如何具体地求解一个线性方程组. 我们用初等变换的方法来做.

\noindent\textbf{例 3.10.1}\quad 求解下列线性方程组：
\[
\left\{\begin{array}{r}
x_1+3x_2-2x_3+4x_4+x_5=2,\\
2x_1+6x_2-3x_3+5x_4+2x_5=5,\\
4x_1+11x_2-7x_3+15x_4+5x_5=3,\\
x_1+3x_2-x_3+x_4+x_5=3.
\end{array}\right.
\]

\noindent\textbf{解}\quad 对方程组的增广矩阵作如下初等变换：
\[
\left(\begin{array}{rrrrr|r}
1&3&-2&4&1&2\\
2&6&-3&5&2&5\\
4&11&-7&15&5&3\\
1&3&-1&1&1&3
\end{array}\right)
\to
\left(\begin{array}{rrrrr|r}
1&3&-2&4&1&2\\
0&0&1&-3&0&1\\
0&-1&1&-1&1&-5\\
0&0&1&-3&0&1
\end{array}\right)
\to
\left(\begin{array}{rrrrr|r}
1&3&-2&4&1&2\\
0&-1&1&-1&1&-5\\
0&0&1&-3&0&1\\
0&0&1&-3&0&1
\end{array}\right)
\]
\[
\to
\left(\begin{array}{rrrrr|r}
1&3&-2&4&1&2\\
0&1&-1&1&-1&5\\
0&0&1&-3&0&1\\
0&0&0&0&0&0
\end{array}\right)
\to
\left(\begin{array}{rrrrr|r}
1&0&1&1&4&-13\\
0&1&-1&1&-1&5\\
0&0&1&-3&0&1\\
0&0&0&0&0&0
\end{array}\right)
\to
\left(\begin{array}{rrrrr|r}
1&0&0&4&4&-14\\
0&1&0&-2&-1&6\\
0&0&1&-3&0&1\\
0&0&0&0&0&0
\end{array}\right).
\]
由此可得原方程组的特解 $\gamma$ 及相伴齐次方程组的基础解系 $\{\eta_1,\eta_2\}$ 如下：
\[
\gamma=\left(\begin{array}{r}
-14\\
6\\
1\\
0\\
0
\end{array}\right),\quad
\eta_1=\left(\begin{array}{r}
-4\\
2\\
3\\
1\\
0
\end{array}\right),\quad
\eta_2=\left(\begin{array}{r}
-4\\
1\\
0\\
0\\
1
\end{array}\right).
\]
原方程组的全部解为 $k_1\eta_1+k_2\eta_2+\gamma$.

上例中我们只用初等行变换就把增广矩阵的左半部分化为对角形. 但有时光靠初等行变换不一定能做到这一点，这时我们可用列的对换. 注意列的对换仅改变未知数的次序而不影响方程组的解，因此是允许的. 我们不必用其他两类列变换就可求出解来.

\noindent\textbf{例 3.10.2}\quad 求解下列线性方程组：
\[
\left\{\begin{array}{r}
x_1+2x_2-x_3+3x_4+x_5=2,\\
2x_1+4x_2-2x_3+6x_4+3x_5=6,\\
-x_1-2x_2+x_3-x_4+3x_5=4.
\end{array}\right.
\]

\noindent\textbf{解}\quad 对增广矩阵进行如下初等变换：
\[
\left(\begin{array}{rrrrr|r}
1&2&-1&3&1&2\\
2&4&-2&6&3&6\\
-1&-2&1&-1&3&4
\end{array}\right)
\to
\left(\begin{array}{rrrrr|r}
1&2&-1&3&1&2\\
0&0&0&0&1&2\\
0&0&0&2&4&6
\end{array}\right)
\to
\left(\begin{array}{rrrrr|r}
1&2&-1&3&1&2\\
0&0&0&0&1&2\\
0&0&0&1&2&3
\end{array}\right)
\]
\[
\to
\left(\begin{array}{rrrrr|r}
1&3&1&2&-1&2\\
0&0&1&0&0&2\\
0&1&2&0&0&3
\end{array}\right)
\to
\left(\begin{array}{rrrrr|r}
1&3&1&2&-1&2\\
0&1&2&0&0&3\\
0&0&1&0&0&2
\end{array}\right)
\to
\left(\begin{array}{rrrrr|r}
1&0&-5&2&-1&-7\\
0&1&2&0&0&3\\
0&0&1&0&0&2
\end{array}\right)
\]
\[
\to
\left(\begin{array}{rrrrr|r}
1&0&-5&2&-1&-7\\
0&1&0&0&0&-1\\
0&0&1&0&0&2
\end{array}\right)
\to
\left(\begin{array}{rrrrr|r}
1&0&0&2&-1&3\\
0&1&0&0&0&-1\\
0&0&1&0&0&2
\end{array}\right).
\]
在上面作初等变换的过程中，我们进行了两次列对换，即第二列和第四列对换以及第三列和第五列对换，这相当于未知数作了如下的变换：
\[
x_1=y_1,\ x_2=y_4,\ x_3=y_5,\ x_4=y_2,\ x_5=y_3.
\]
根据最后一个矩阵，我们可取一组特解 $\gamma$ 以及相伴齐次方程组的基础解系 $\{\eta_1,\eta_2\}$ 为 (关于 $y_i$ 的)：
\[
\gamma=\left(\begin{array}{r}
3\\
-1\\
2\\
0\\
0
\end{array}\right),\quad
\eta_1=\left(\begin{array}{r}
-2\\
0\\
0\\
1\\
0
\end{array}\right),\quad
\eta_2=\left(\begin{array}{r}
1\\
0\\
0\\
0\\
1
\end{array}\right).
\]
将 $y_i$ 换成 $x_i$，得到特解 $\delta$ 以及基础解系 $\{\xi_1,\xi_2\}$ 为：
\[
\delta=\left(\begin{array}{r}
3\\
0\\
0\\
-1\\
2
\end{array}\right),\quad
\xi_1=\left(\begin{array}{r}
-2\\
1\\
0\\
0\\
0
\end{array}\right),\quad
\xi_2=\left(\begin{array}{r}
1\\
0\\
1\\
0\\
0
\end{array}\right).
\]
原线性方程组的解为
\[
k_1\xi_1+k_2\xi_2+\delta,
\]
或
\[
\left\{\begin{array}{l}
x_1=-2k_1+k_2+3,\\
x_2=k_1,\\
x_3=k_2,\\
x_4=-1,\\
x_5=2.
\end{array}\right.
\]

现在我们把求一般线性方程组解的办法总结如下：

第一步：对增广矩阵 $\widetilde{A}$ 进行初等行变换及第一类初等列变换 (列对换)，注意用虚线将常数列隔开且不允许常数列与其他列对换. 经过若干次初等变换后将 $\widetilde{A}$ 化成下列形状：
\[
\left(\begin{array}{ccccccc|c}
1&0&\cdots&0&c_{1,r+1}&\cdots&c_{1n}&d_1\\
0&1&\cdots&0&c_{2,r+1}&\cdots&c_{2n}&d_2\\
\vdots&\vdots&&\vdots&\vdots&&\vdots&\vdots\\
0&0&\cdots&1&c_{r,r+1}&\cdots&c_{rn}&d_r\\
&&&&&&&*
\end{array}\right).
\]
上述矩阵空白部分全为零，若 $*$ 部分有元素不等于零，则表示 $\mathrm{r}(\widetilde{A})>\mathrm{r}(A)$，因此原方程组无解；若 $*$ 部分所有元素全为零，则原方程组有解.

第二步：写出特解 (注意：若有列变换，这还不是原方程组的特解) 以及相伴齐次线性方程组的基础解系 (若有列变换，这也不是原方程组相伴齐次线性方程组的基础解系)：
\[
\gamma=\left(\begin{array}{c}
d_1\\
\vdots\\
d_r\\
0\\
\vdots\\
0
\end{array}\right),\quad
\eta_1=\left(\begin{array}{c}
-c_{1,r+1}\\
\vdots\\
-c_{r,r+1}\\
1\\
\vdots\\
0
\end{array}\right),\quad \cdots,\quad
\eta_{n-r}=\left(\begin{array}{c}
-c_{1n}\\
\vdots\\
-c_{rn}\\
0\\
\vdots\\
1
\end{array}\right).
\]

第三步：根据列对换情况，调整 $\gamma,\eta_1,\cdots,\eta_{n-r}$ 的各分量，得到原方程组的特解 $\delta$ 以及原方程组相伴齐次线性方程组的基础解系 $\xi_1,\cdots,\xi_{n-r}$. 最后写出原方程组的解：
\[
k_1\xi_1+\cdots+k_{n-r}\xi_{n-r}+\delta,
\]
其中 $k_1,\cdots,k_{n-r}$ 为参变数.

\subsection*{习题 3.10}

1. 求解下列线性方程组：
\[
(1)\left\{\begin{array}{r}
x_1-x_2+2x_3=3,\\
x_1-x_2+x_3-2x_4=2,\\
-x_1+x_2-x_3+2x_4=0,\\
-3x_1+x_2-8x_3-10x_4=-1;
\end{array}\right.
\]
\[
(2)\left\{\begin{array}{r}
x_1+2x_2-3x_3-2x_4=-6,\\
-2x_1-5x_2+8x_3+3x_4=14,\\
x_1+2x_2-5x_3=0,\\
x_1+2x_2-2x_3-3x_4=-9;
\end{array}\right.
\]
\[
(3)\left\{\begin{array}{r}
x_1+2x_2+x_3-3x_4+2x_5=4,\\
2x_1+x_2+x_3+x_4-4x_5=4,\\
x_1+x_2+2x_3+2x_4-2x_5=4,\\
2x_1+3x_2-5x_3-17x_4+8x_5=0;
\end{array}\right.
\]
\[
(4)\left\{\begin{array}{r}
x_1+2x_2+x_3-x_4=2,\\
2x_1+4x_2+3x_3-3x_4-x_5=5,\\
-x_1-x_2-3x_3+3x_4+2x_5=-2;
\end{array}\right.
\]
\[
(5)\left\{\begin{array}{r}
x_1-x_2+2x_3+2x_4+3x_5=1,\\
2x_1-2x_2+4x_3+5x_4+7x_5=3,\\
-x_1+x_2-2x_3-x_5=1.
\end{array}\right.
\]

2. 求解下列线性方程组，其中 $\lambda$ 为参数：
\[
\left\{\begin{array}{r}
2x_1+3x_2+x_3+x_4=1,\\
x_1+2x_2-x_3+4x_4=2,\\
x_1+3x_2-4x_3+11x_4=\lambda.
\end{array}\right.
\]

3. 求解下列线性方程组，其中 $\lambda$ 为参数：
\[
\left\{\begin{array}{r}
\lambda x_1+x_2+x_3+x_4=1,\\
x_1+\lambda x_2+x_3+x_4=\lambda,\\
x_1+x_2+\lambda x_3+x_4=\lambda^2,\\
x_1+x_2+x_3+\lambda x_4=\lambda^3.
\end{array}\right.
\]

4. 求解下列线性方程组，其中 $a,b$ 为参数：
\[
\left\{\begin{array}{r}
x_1+2x_2+3x_3+3x_4+7x_5=b,\\
3x_1+2x_2+x_3+x_4+(a-5)x_5=b+2,\\
x_2+2x_3+2x_4+3ax_5=b,\\
5x_1+4x_2+3x_3+3x_4+(1-a)x_5=-b.
\end{array}\right.
\]

5. 判定下列向量 $\beta$ 能否用向量组 $\alpha_1,\alpha_2,\alpha_3$ 线性表示：

(1) $\beta=(5,4,-2,4)$；$\alpha_1=(1,0,-1,3)$，$\alpha_2=(2,1,0,1)$，$\alpha_3=(0,-2,1,1)$；

(2) $\beta=(1,1,1,1)$；$\alpha_1=(-1,2,-1,1)$，$\alpha_2=(4,0,1,-1)$，$\alpha_3=(3,2,0,0)$.

6. 设 $\alpha_1,\alpha_2$ 是某个齐次线性方程组的基础解系，问 $\alpha_1+\alpha_2$ 及 $2\alpha_1-\alpha_2$ 是否也是这个齐次线性方程组的基础解系？为什么？

7. 设 $Ax=0$ 是由 $n$ 个未知数 $n$ 个方程式组成的线性方程组. 证明：$Ax=0$ 只有零解当且仅当对任意的正整数 $k$，线性方程组 $A^kx=0$ 也只有零解.

8. 设 $A=(a_{ij})$ 是一个 $m\times n$ 矩阵，记 $\alpha_i$ 为 $A$ 的第 $i$ 个行向量，$\beta=(b_1,b_2,\cdots,b_n)$. 求证：如果齐次线性方程组 $Ax=0$ 的解全是方程 $b_1x_1+b_2x_2+\cdots+b_nx_n=0$ 的解，则 $\beta$ 是 $\alpha_1,\alpha_2,\cdots,\alpha_m$ 的线性组合.

9. 设 $Ax=\beta$ 是 $m$ 个方程式 $n$ 个未知数的线性方程组，求证：它有解的充分必要条件是方程组 $A'y=0$ 的任一解 $\alpha$ 均适合等式 $\alpha'\beta=0$.

10. 设有两个线性方程组：
\[
\left\{\begin{array}{c}
a_{11}x_1+a_{12}x_2+\cdots+a_{1n}x_n=b_1,\\
a_{21}x_1+a_{22}x_2+\cdots+a_{2n}x_n=b_2,\\
\cdots\cdots\cdots\\
a_{m1}x_1+a_{m2}x_2+\cdots+a_{mn}x_n=b_m;
\end{array}\right.
\tag{1}
\]
\[
\left\{\begin{array}{c}
a_{11}x_1+a_{21}x_2+\cdots+a_{m1}x_m=0,\\
a_{12}x_1+a_{22}x_2+\cdots+a_{m2}x_m=0,\\
\cdots\cdots\cdots\\
a_{1n}x_1+a_{2n}x_2+\cdots+a_{mn}x_m=0,\\
b_1x_1+b_2x_2+\cdots+b_mx_m=1.
\end{array}\right.
\tag{2}
\]
求证：方程组 (1) 有解的充分必要条件是方程组 (2) 无解.

11. 设 $A$ 是秩为 $r$ 的 $m\times n$ 矩阵，$\alpha_1,\cdots,\alpha_{n-r}$ 与 $\beta_1,\cdots,\beta_{n-r}$ 是齐次线性方程组 $Ax=0$ 的两个基础解系. 求证：必存在一个 $n-r$ 阶非异阵 $P$，使
\[
(\beta_1,\cdots,\beta_{n-r})=(\alpha_1,\cdots,\alpha_{n-r})P.
\]

12. 如果 $n$ 阶实方阵 $A=(a_{ij})$ 适合条件：
\[
|a_{ii}|>\sum_{j=1,j\neq i}^n |a_{ij}|,\quad i=1,2,\cdots,n,
\]
则称 $A$ 为严格对角占优阵. 求证：严格对角占优阵必是非异阵. 若上述条件改为
\[
a_{ii}>\sum_{j=1,j\neq i}^n |a_{ij}|,\quad i=1,2,\cdots,n,
\]
求证：$|A|>0$.

13. 设数域 $\mathbb{K}$ 上的 $n$ 阶方阵 $A=(a_{ij})$ 满足：$|A|=0$ 且某个元素 $a_{ij}$ 的代数余子式 $A_{ij}\neq 0$. 求证：齐次线性方程组 $Ax=0$ 的所有解都可写为下列形式：
\[
k\left(\begin{array}{c}
A_{i1}\\
A_{i2}\\
\vdots\\
A_{in}
\end{array}\right),\quad k\in\mathbb{K}.
\]

14. 设 $n$ 阶方阵 $A$ 的行列式等于零，证明：伴随矩阵 $A^*$ 的秩不超过 1.

\section*{历史与展望}

数学家们创建了行列式理论和矩阵理论，但这还不足以完全解决线性方程组的求解问题，为此还需要创建线性空间理论. 现在的教科书都是先给出线性空间的公理化定义，然后逐步给出向量的线性关系、线性空间的基和维数等概念. 但回顾线性空间的发展历史，其公理化定义经历了漫长的过程，而线性无关、基和维数等概念在线性空间的定义尚未严格建立时就得到了应用.

事实上到 1880 年为止，线性代数的许多重要结果都已经建立了，但它们还不能作为线性空间理论的组成部分. 其主要原因是，作为这一理论的核心，线性空间的定义尚未给出. 而线性空间的严格定义在 1888 年才由 Peano（皮亚诺）给出.

矢量（向量）最早的定义起源于物理学，它意味着既有大小又有方向的量 (如速度或力). 矢量的平行四边形法则定义了两个矢量的加法，这一思想在 17 世纪末得到了确立. 由此，矢量的加法和数乘具有了明确的物理意义.

向量的数学定义起源于复数的几何表示，在 18 世纪末到 19 世纪初，由多位数学家独立地给出. 1797 年 Wessel (韦塞尔) 和 1831 年 Gauss 将复数几何地表示为复平面上的点或有向线段；1835 年 Hamilton 将复数代数地定义为有序实数对，并带有通常的加法、数乘以及乘法. 他指出这些有序实数对满足运算封闭性、交换律、结合律和分配律，有零元、加法逆元和乘法逆元.

将向量的思想扩张到三维空间是一个重要的发展. Hamilton 在他的四元数系统中构造了一个向量代数，每个向量表示为 $ai+bj+ck$ 的形式，其中 $a,b,c$ 是实数，$\mathrm{i},\mathrm{j},\mathrm{k}$ 是四元数单位，它们就是这个三维线性空间的一组基.

在 1840 年代，Hamilton、Cayley 和 Grassmann（格拉斯曼）将三维空间进一步推广到高维空间，这是线性空间理论发展的里程碑. 1843 年，Hamilton 构造了四元数，这是一个四维向量空间 (也是一个可除代数)，此后他花了 20 年时间从事四元数的研究与应用. Cayley 关于维数的思想出现在 1843 年的论文《$n$ 维解析几何章节》中. Grassmann 在 1844 年的论文《线性扩张学说》中阐述了一些开创性的思想. 他的目标是建立一个坐标自由的 $n$ 维向量空间，其中包含了许多线性代数的基本思想，例如 $n$ 维向量空间、子空间、线性扩张集合、线性无关、基、维数和线性变换的概念等. 他还证明了关于向量空间的很多结果，例如交空间与和空间的维数公式等. 由于包含很多用哲学语言表达的新思想，Grassmann 的论文深奥难懂，一直被数学界所忽略，直到 1862 年版本的出现.

受 Grassmann 思想的启发，1888 年 Peano 在著作《几何演算》的最后一章“线性系统的变换”中给出了实数域上向量空间的公理化定义，他称向量空间为线性系统. Peano 给出的向量空间的公理化定义非常接近我们现在看到的向量空间的定义. 他给出了向量空间的例子，例如：实数域、复数域、平面和三维欧氏空间，两个向量空间之间线性映射全体，一元多项式全体. 他还定义了线性代数中的其他概念，例如维数和线性变换，证明了一些重要定理.

20 世纪上半叶，数学公理化的方法才得到长足的发展，群、域、正整数 (Peano 公理) 和射影几何的公理化定义相继出现. 而在 19 世纪末的 20 年，Peano 关于向量空间的公理化定义很大程度上被忽视了. 1918 年，Weyl (外尔) 在著作《空间、时间、物质》的第一章“仿射几何基础”中给出了有限维实向量空间的公理化定义. 显然，他并没有注意到 Peano 的工作.

1920 年 Banach (巴拿赫) 在他的博士论文中给出了实数域上完备赋范线性空间 (后称为 Banach 空间) 的公理化定义，其中前 13 个公理就是线性空间的公理. 1921 年 Noether (诺特) 在论文《环中的理想理论》中引入了模的概念，而将线性空间看成是模的特例. 因此，我们看到线性空间来源于三个不同的数学分支：几何、分析和代数. 1930 年，Van der Waerden (范德瓦尔登) 在其经典教材《近世代数》中单独设有一章“线性代数”，这是“线性代数”这个词语第一次出现在近现代数学的著作中.

在行列式理论、矩阵理论和线性空间理论相继建立之后，线性方程组的求解问题得到了圆满的解决. 回顾历史，上述三种理论在近现代数学中体现的价值远远大于解决线性方程组的求解问题本身. 特别地，线性空间理论更是成为了近现代数学的基石之一.

\section*{复习题三}

1. 设向量组 $\alpha_1,\alpha_2,\cdots,\alpha_r$ 线性无关，又
\[
\left\{\begin{array}{c}
\beta_1=a_{11}\alpha_1+a_{12}\alpha_2+\cdots+a_{1r}\alpha_r,\\
\beta_2=a_{21}\alpha_1+a_{22}\alpha_2+\cdots+a_{2r}\alpha_r,\\
\cdots\cdots\cdots\\
\beta_r=a_{r1}\alpha_1+a_{r2}\alpha_2+\cdots+a_{rr}\alpha_r.
\end{array}\right.
\]
求证：$\beta_1,\beta_2,\cdots,\beta_r$ 线性相关的充分必要条件是系数矩阵 $A=(a_{ij})_{r\times r}$ 的行列式等于零.

2. 设 $\alpha_1,\alpha_2,\cdots,\alpha_m$ 是一组线性无关的向量，若向量组 $\beta_1,\beta_2,\cdots,\beta_k$ 可用 $\alpha_1,\alpha_2,\cdots,\alpha_m$ 线性表示如下：
\[
\left\{\begin{array}{c}
\beta_1=a_{11}\alpha_1+a_{12}\alpha_2+\cdots+a_{1m}\alpha_m,\\
\beta_2=a_{21}\alpha_1+a_{22}\alpha_2+\cdots+a_{2m}\alpha_m,\\
\cdots\cdots\cdots\\
\beta_k=a_{k1}\alpha_1+a_{k2}\alpha_2+\cdots+a_{km}\alpha_m.
\end{array}\right.
\]
记表示矩阵 $A=(a_{ij})_{k\times m}$. 求证：向量组 $\beta_1,\beta_2,\cdots,\beta_k$ 的秩等于 $\mathrm{r}(A)$.

3. 设 $V$ 是实数域上连续函数全体构成的实线性空间，求证下列函数线性无关：

(1) $\sin x,\sin 2x,\cdots,\sin nx$；\qquad
(2) $1,\cos x,\cos 2x,\cdots,\cos nx$；

(3) $1,\sin x,\cos x,\sin 2x,\cos 2x,\cdots,\sin nx,\cos nx$.

4. 设 $\alpha_1,\alpha_2,\cdots,\alpha_m$ 是向量空间 $V$ 中一组向量，向量组 $\beta_1,\beta_2,\cdots,\beta_k$ 可用 $\alpha_1,\alpha_2,\cdots,\alpha_m$ 线性表示，求证：向量组 $\beta_1,\beta_2,\cdots,\beta_k$ 的秩小于等于向量组 $\alpha_1,\alpha_2,\cdots,\alpha_m$ 的秩.

5. 设 $\alpha_1,\alpha_2,\cdots,\alpha_m$ 是向量空间 $V$ 中一组向量且其秩等于 $r$，$\alpha_{i_1},\alpha_{i_2},\cdots,\alpha_{i_r}$ 是其中 $r$ 个向量. 假设下列条件之一成立：

(1) $\alpha_{i_1},\alpha_{i_2},\cdots,\alpha_{i_r}$ 线性无关；

(2) 任一 $\alpha_i$ 均可由 $\alpha_{i_1},\alpha_{i_2},\cdots,\alpha_{i_r}$ 线性表示.

求证：$\alpha_{i_1},\alpha_{i_2},\cdots,\alpha_{i_r}$ 是向量组的极大无关组.

6. 证明：线性空间 $V$ 中两个向量组 $A,B$ 等价的充分必要条件是它们的秩相同，且向量组 $A$ 可用向量组 $B$ 线性表示. 举例说明秩相同的两个向量组未必等价.

7. 设 $\mathbb{K}_1,\mathbb{K}_2,\mathbb{K}_3$ 是数域且 $\mathbb{K}_1\subseteq\mathbb{K}_2\subseteq\mathbb{K}_3$，若将 $\mathbb{K}_2$ 看成是 $\mathbb{K}_1$ 上的线性空间，其维数为 $m$，又将 $\mathbb{K}_3$ 看成是 $\mathbb{K}_2$ 上的线性空间，其维数为 $n$，求证：如将 $\mathbb{K}_3$ 看成是 $\mathbb{K}_1$ 上的线性空间，则其维数为 $mn$.

8. 设 $\alpha_1,\alpha_2,\cdots,\alpha_m$ 是数域 $\mathbb{F}$ 上的 $n$ 维线性空间 $V$ 中选定的 $m$ 个向量且已知它们的秩等于 $r$. 求证：全体适合 $x_1\alpha_1+x_2\alpha_2+\cdots+x_m\alpha_m=0$ 的列向量 $(x_1,x_2,\cdots,x_m)'(x_i\in\mathbb{F})$ 构成数域 $\mathbb{F}$ 上 $m$ 维列向量空间 $\mathbb{F}^m$ 的 $m-r$ 维子空间.

9. 设 $V_1,V_2$ 分别是数域 $\mathbb{K}$ 上齐次线性方程组 $x_1=x_2=\cdots=x_n$ 与 $x_1+x_2+\cdots+x_n=0$ 的解空间，求证：$\mathbb{K}^n=V_1\oplus V_2$.

10. 设 $V_1,V_2,\cdots,V_m$ 是线性空间 $V$ 的 $m$ 个非平凡子空间，证明：$V$ 中必有一组基，使得每个基向量都不在诸 $V_i$ 的并中.

11. 设 $V$ 是数域 $\mathbb{K}$ 上的线性空间，$U$ 是 $V$ 的子空间. 对任意的 $v\in V$，集合 $v+U:=\{v+u\mid u\in U\}$ 称为 $v$ 的 $U$-陪集. 在所有 $U$-陪集构成的集合 $S=\{v+U\mid v\in V\}$ 中，定义加法和数乘如下，其中 $v_1,v_2\in V,k\in\mathbb{K}$：
\[
(v_1+U)+(v_2+U):=(v_1+v_2)+U,\quad k\cdot(v_1+U):=k\cdot v_1+U.
\]
证明下列结论成立：

(1) $U$-陪集之间的关系是：作为集合或者相等，或者不相交.

(2) $v_1+U=v_2+U$ (作为集合相等) 当且仅当 $v_1-v_2\in U$. 特别地，$v+U$ 是 $V$ 的子空间当且仅当 $v\in U$.

(3) $S$ 中的加法以及 $\mathbb{K}$ 关于 $S$ 的数乘不依赖于代表元的选取，即若 $v_1+U=v_1'+U$ 以及 $v_2+U=v_2'+U$，则 $(v_1+U)+(v_2+U)=(v_1'+U)+(v_2'+U)$，以及 $k\cdot(v_1+U)=k\cdot(v_1'+U)$.

(4) $S$ 在上述加法和数乘下成为数域 $\mathbb{K}$ 上的线性空间，称为 $V$ 关于子空间 $U$ 的商空间，记为 $V/U$.

12. 设 $V$ 是数域 $\mathbb{K}$ 上的 $n$ 维线性空间，$U$ 是 $V$ 的子空间，$W$ 是 $U$ 的补空间，证明：$\dim V/U=\dim V-\dim U$，并且存在线性同构 $\varphi:W\to V/U$.

13. 设 $A=(a_{ij}),B=(b_{ij})$ 是 $m\times n$ 矩阵，且 $b_{ij}=(-1)^{i+j}a_{ij}$，求证：$\mathrm{r}(B)=\mathrm{r}(A)$.

14. 设 $A_1,A_2,\cdots,A_m$ 为 $n$ 阶方阵，求证：
\[
\mathrm{r}(A_1)+\mathrm{r}(A_2)+\cdots+\mathrm{r}(A_m)\leq (m-1)n+\mathrm{r}(A_1A_2\cdots A_m).
\]

15. 证明：$\mathrm{r}(ABC)\geq \mathrm{r}(AB)+\mathrm{r}(BC)-\mathrm{r}(B)$.

16. 设有分块矩阵 $M=\left(\begin{array}{cc}A&B\\ C&D\end{array}\right)$，证明：

(1) 若 $A$ 可逆，则 $\mathrm{r}(M)=\mathrm{r}(A)+\mathrm{r}(D-CA^{-1}B)$；

(2) 若 $D$ 可逆，则 $\mathrm{r}(M)=\mathrm{r}(D)+\mathrm{r}(A-BD^{-1}C)$；

(3) 若 $A,D$ 都可逆，则 $\mathrm{r}(A)+\mathrm{r}(D-CA^{-1}B)=\mathrm{r}(D)+\mathrm{r}(A-BD^{-1}C)$.

17. 设
\[
M=\left(\begin{array}{cccc}
a_1^2&a_1a_2+1&\cdots&a_1a_n+1\\
a_2a_1+1&a_2^2&\cdots&a_2a_n+1\\
\vdots&\vdots&&\vdots\\
a_na_1+1&a_na_2+1&\cdots&a_n^2
\end{array}\right),
\]
证明：$\mathrm{r}(M)\geq n-1$，等号成立当且仅当 $|M|=0$.

18. 设 $A,B$ 都是数域 $\mathbb{K}$ 上的 $n$ 阶矩阵且 $AB=BA$，证明：
\[
\mathrm{r}(A+B)\leq \mathrm{r}(A)+\mathrm{r}(B)-\mathrm{r}(AB).
\]

19. 设 $A$ 是 $m\times n$ 实矩阵，求证：$\mathrm{r}(A'A)=\mathrm{r}(AA')=\mathrm{r}(A)$.

20. 设 $A,B$ 是 $\mathbb{K}$ 上的 $n$ 阶矩阵，若线性方程组 $Ax=0$ 和 $Bx=0$ 同解，且每个方程组的基础解系含有 $m$ 个线性无关的向量，求证：$\mathrm{r}(A-B)\leq n-m$.

21. 设 $A$ 是 $m\times n$ 矩阵，$B$ 是 $n\times k$ 矩阵，证明：方程组 $ABx=0$ 和方程组 $Bx=0$ 同解的充分必要条件是 $\mathrm{r}(AB)=\mathrm{r}(B)$.

22. 设 $A$ 是 $m\times n$ 矩阵，$B$ 是 $n\times k$ 矩阵. 若 $AB$ 和 $B$ 有相同的秩，求证：对任意的 $k\times l$ 矩阵 $C$，矩阵 $ABC$ 和矩阵 $BC$ 也有相同的秩.

23. 设 $A$ 是 $n$ 阶实反对称阵，$D=\operatorname{diag}\{d_1,d_2,\cdots,d_n\}$ 是同阶对角阵且主对角元素全大于零，求证：$|A+D|>0$. 特别，$|I_n\pm A|>0$，从而 $I_n\pm A$ 都是非异阵.

24. 设 $A$ 是 $n$ 阶实对称阵，求证：$I_n+\mathrm{i}A$ 和 $I_n-\mathrm{i}A$ 都是非异阵.

25. 设 $A$ 是矩阵，$|D|$ 是 $A$ 的 $r$ 阶子式，$A$ 中所有包含 $|D|$ 为 $r$ 阶子式的 $r+1$ 阶子式称为 $|D|$ 的 $r+1$ 阶加边子式. 求证：矩阵 $A$ 的秩等于 $r$ 的充分必要条件是 $A$ 存在一个 $r$ 阶子式 $|D|$ 不等于零而 $|D|$ 的所有 $r+1$ 阶加边子式全等于零.

26. 设 $m\times n$ 矩阵 $A$ 的 $m$ 个行向量为 $\alpha_1,\alpha_2,\cdots,\alpha_m$，且 $\alpha_{i_1},\alpha_{i_2},\cdots,\alpha_{i_r}$ 是其极大无关组. 又设 $A$ 的 $n$ 个列向量为 $\beta_1,\beta_2,\cdots,\beta_n$，且 $\beta_{j_1},\beta_{j_2},\cdots,\beta_{j_r}$ 是其极大无关组. 证明：$\alpha_{i_1},\alpha_{i_2},\cdots,\alpha_{i_r}$ 和 $\beta_{j_1},\beta_{j_2},\cdots,\beta_{j_r}$ 交叉点上的元素组成的子矩阵 $D$ 的行列式 $|D|\neq 0$.

27. 设 $A$ 是一个方阵，$A$ 的第 $i_1,\cdots,i_r$ 行和第 $i_1,\cdots,i_r$ 列交叉点上的元素组成的子式称为 $A$ 的一个主子式. 若 $A$ 是对称阵或反对称阵且秩为 $r$，求证：$A$ 必有一个 $r$ 阶主子式不等于零.

28. 证明：反对称阵的秩必为偶数.

29. 求证：秩等于 $r$ 的矩阵可以表示为 $r$ 个秩等于 1 的矩阵之和，但不能表示为少于 $r$ 个秩为 1 的矩阵之和.

30. 设 $A$ 为 $m\times n$ 矩阵，证明：存在 $n\times m$ 矩阵 $B$，使得 $ABA=A$.

31. 设 $A,B$ 分别是 $3\times 2,2\times 3$ 矩阵且满足
\[
AB=\left(\begin{array}{rrr}
8&2&-2\\
2&5&4\\
-2&4&5
\end{array}\right),
\]
试求 $BA$.

32. 设 $A$ 是 $n$ 阶方阵且 $\mathrm{r}(A)=r$，求证：$A^2=A$ 的充分必要条件是存在秩等于 $r$ 的 $n\times r$ 矩阵 $S$ 和秩等于 $r$ 的 $r\times n$ 矩阵 $T$，使 $A=ST,TS=I_r$.

33. 设 $A$ 是秩为 $r$ 的 $m\times n$ 矩阵，求证：必存在秩为 $n-r$ 的 $n\times(n-r)$ 矩阵 $B$，使 $AB=O$.

34. 设
\[
A=\left(\begin{array}{cccc}
a_{11}&a_{12}&\cdots&a_{1n}\\
a_{21}&a_{22}&\cdots&a_{2n}\\
\vdots&\vdots&&\vdots\\
a_{m1}&a_{m2}&\cdots&a_{mn}
\end{array}\right)\quad (m<n),
\]
已知 $Ax=0$ 的基础解系为 $\beta_i=(b_{i1},b_{i2},\cdots,b_{in})'(i=1,2,\cdots,n-m)$，试求线性方程组
\[
\sum_{j=1}^n b_{ij}y_j=0\quad (i=1,2,\cdots,n-m)
\]
的基础解系.

35. 设 $V_0$ 是 $\mathbb{K}_n$ 的一个非平凡子空间，求证：必存在矩阵 $A$，使 $V_0$ 是 $n$ 元齐次线性方程组 $Ax=0$ 的解空间.

36. 设 $A,B$ 为 $m\times n$ 和 $m\times p$ 矩阵，$X$ 为 $n\times p$ 未知矩阵，证明：矩阵方程 $AX=B$ 有解的充分必要条件是 $\mathrm{r}(A;B)=\mathrm{r}(A)$.

37. 设
\[
A=\left(\begin{array}{rrrr}
1&1&2&1\\
1&2&3&3\\
2&3&5&4\\
3&5&8&7
\end{array}\right),\quad
B=\left(\begin{array}{rr}
1&1\\
5&-1\\
6&0\\
11&-1
\end{array}\right),
\]
$X$ 为 $4\times 2$ 未知矩阵，试求矩阵方程 $AX=B$ 的解.

38. 设 $A,B$ 为 $m\times n$ 和 $n\times p$ 矩阵，证明：存在 $p\times n$ 矩阵 $C$，使 $ABC=A$ 成立的充分必要条件是 $\mathrm{r}(A)=\mathrm{r}(AB)$.

39. 设有两个非齐次线性方程组 (I) 和 (II)，它们的通解分别为
\[
\gamma+t_1\eta_1+t_2\eta_2;\quad \delta+k_1\xi_1+k_2\xi_2,
\]
其中
\[
\gamma=\left(\begin{array}{r}
5\\
-3\\
0\\
0
\end{array}\right),\quad
\eta_1=\left(\begin{array}{r}
-6\\
5\\
1\\
0
\end{array}\right),\quad
\eta_2=\left(\begin{array}{r}
-5\\
4\\
0\\
1
\end{array}\right);
\quad
\delta=\left(\begin{array}{r}
-11\\
3\\
0\\
0
\end{array}\right),\quad
\xi_1=\left(\begin{array}{r}
8\\
-1\\
1\\
0
\end{array}\right),\quad
\xi_2=\left(\begin{array}{r}
10\\
-2\\
0\\
1
\end{array}\right).
\]
试求这两个方程组的公共解.

40. 设有非齐次线性方程组 (I)：
\[
\left\{\begin{array}{r}
7x_1-6x_2+3x_3=b,\\
8x_1-9x_2+ax_4=7.
\end{array}\right.
\]
又已知方程组 (II) 的通解为
\[
(1,1,0,0)'+t_1(1,0,-1,0)'+t_2(2,3,0,1)'.
\]
若这两个方程组有无穷多组公共解，求出 $a,b$ 的值并求出公共解.

41. 求平面上 $n$ 个点 $(x_1,y_1),(x_2,y_2),\cdots,(x_n,y_n)$ 位于同一条直线上的充分必要条件.

42. 求三维实空间中 4 个点 $(x_i,y_i,z_i)(i=1,2,3,4)$ 共面的充分必要条件.

43. 证明：通过平面内不在一条直线上的 3 点 $(x_1,y_1),(x_2,y_2),(x_3,y_3)$ 的圆方程为
\[
\left|\begin{array}{cccc}
x^2+y^2&x&y&1\\
x_1^2+y_1^2&x_1&y_1&1\\
x_2^2+y_2^2&x_2&y_2&1\\
x_3^2+y_3^2&x_3&y_3&1
\end{array}\right|=0.
\]
以此证明：通过具有有理坐标的 3 点的圆，其圆心也是有理坐标.

44. 求平面上不在一条直线上的 4 个点 $(x_1,y_1),(x_2,y_2),(x_3,y_3),(x_4,y_4)$ 位于同一个圆上的充分必要条件.

45. 已知平面上两条不同的二次曲线 $a_ix^2+b_ixy+c_iy^2+d_ix+e_iy+f_i=0(i=1,2)$ 交于 4 个不同的点 $(x_i,y_i)(i=1,2,3,4)$. 求证：过这 4 个点的二次曲线均可写为如下形状：
\[
\lambda_1(a_1x^2+b_1xy+c_1y^2+d_1x+e_1y+f_1)+\lambda_2(a_2x^2+b_2xy+c_2y^2+d_2x+e_2y+f_2)=0.
\]

\end{document}
